%% ---------------------------------------------------------------------------
%% Springer Nature LaTeX template (sn-jnl.cls), IJCV submission.
%% Migrated from Markdown by _migrate/migrate.py -- see MIGRATION_REPORT.md.
%% Reference style: author--year (sn-mathphys-ay), the IJCV convention.
%% ---------------------------------------------------------------------------
\documentclass[pdflatex,sn-mathphys-ay]{sn-jnl}

\usepackage{iftex}
\ifPDFTeX
  \usepackage[T1]{fontenc}
  \usepackage[utf8]{inputenc}
  %% The two non-ASCII markers kept literally in the source text.
  \DeclareUnicodeCharacter{2705}{\ensuremath{\checkmark}}
  \DeclareUnicodeCharacter{26A0}{\textbf{!}}
\else
  \usepackage{newunicodechar}
  \newunicodechar{^^^^2705}{\ensuremath{\checkmark}}
  \newunicodechar{^^^^26a0}{\textbf{!}}
\fi
\usepackage{graphicx}
\usepackage{array}
\usepackage{longtable}
\usepackage{booktabs}
\usepackage{amsmath,amssymb,amsfonts}
\usepackage{amsthm}
\usepackage{mathrsfs}
\usepackage{textcomp}
\usepackage{multirow}
\usepackage{xcolor}

\raggedbottom

\begin{document}

%% Author metadata is not present in the Markdown source; the double-blind
%% placeholders below are to be replaced at submission.
\author*[1]{\fnm{Author names} \sur{withheld for double-blind review}}
\affil*[1]{\orgname{Affiliations withheld for double-blind review}}

\title[Supplementary material]{Supplementary material --- A measurement-and-audit protocol for budgeted fine-tuning of lightweight detectors}

\maketitle

\tableofcontents
\clearpage

This file carries the appendices of the manuscript: \textbf{supplementary material}, not typeset into the article, published online as supplied, and \textbf{not part of the manuscript's page budget}. The letters are unchanged from the manuscript's cross-references, so "Appendix J.4" in the main text resolves to the section of that name here. The appendices run \textbf{A--U}. (The reference-status record of an earlier 74-entry numbering, and the internal compression record it named, are \textbf{not} part of this file and are not supplied.)

\begin{quote}
⚠ \textbf{On the section-number pointers in this file.} A pointer of the form \emph{"\S{}8.6"} points into \textbf{an earlier version of the manuscript} and is \textbf{not} a live cross-reference: in the present article \S{}5--\S{}9 carry no subsections, so it will not resolve there. The same holds for every "\S{}x.y" form here: it either names a subsection \emph{of this supplement} (where the letter is an appendix letter, e.g. "P.1") or is a historical pointer of that kind, and is \textbf{never a live pointer into the article's \S{}5--\S{}9}. Pointers of the form \textbf{\S{}4.x do resolve}: the article's \S{}4 has subsections 4.1--4.6. The \textbf{appendix letters (A--U)} are the live pointers into this file: they are what the main text cites.
\end{quote}

\subsection*{Figures carried in this supplement, and how they are cited}
\addcontentsline{toc}{subsection}{Figures carried in this supplement, and how they are cited}

Five figures are carried here rather than in the article file, so that the article is paginated to the journal's length rule. \textbf{As the article stands, one of them (Fig. S3) is cited from the main text} and the other four are not: the readings they illustrate are reported in the article in tables and in words, and the figures are the released record of those readings. Their numbering is therefore internal to the supplement, Fig. S1 to Fig. S5. (The protocol schematic, \textbf{Fig. 1}, is carried in the article file and is the article's only figure.)

\textbf{Fig. S1. Budget curve, SHWD\ensuremath{\rightarrow}SFCHD} --- point estimates +1.00 / +0.65 / +0.51 pp at 30 / 50 / 100 epochs (3 seeds each): the gain decreasing as the budget grows, then plateauing (+0.52 pp at 200 epochs).

\textbf{Fig. S2. Shift ordering (heuristic)} --- +0.07 (D = 9.44) / +0.51 (D = 12.76) / +1.78 (D = 12.99); a two-segment heuristic (D \textless{} 10 vs. D \ensuremath{\approx} 13) on three points whose label spaces are not equivalent.

\textbf{Fig. S3. Dual-metric separation} --- at 100 epochs strategy gains turn negative in mAP50 while staying non-negative in mAP50-95 in \textbf{4 of the 15} long-budget cells (per-cell signs and marginals in Table T12).

\textbf{Fig. S4. Measurement-hierarchy schematic} --- trajectory-evaluated first-order reads (worst-direction increments, gradient norms, alignment cosines) against second-order stationary reads; the proposition bans a \emph{use}, not the quantities.

\textbf{Fig. S5. Prediction--evidence--status matrix} --- per-prediction status across \S{}6--\S{}7 (four evidence lines; SNR floor 18/18 \textless{} 2).

\section*{Appendix A. Lemmas used by the bound, and the linear displacement refinement}
\addcontentsline{toc}{section}{Appendix A. Lemmas used by the bound, and the linear displacement refinement}

\textbf{Statements (assumption numbering H1--H10 as in Appendix J.1).} The derivations are \textbf{released} as \texttt{provenance/theory\_supplement/} (three working notes; Lemma 1 and Lemma 2 are textbook Kantorovich--Rubinstein results and are cited as such, and \textbf{no result of \S{}4 depends on any non-standard step} --- Lemma 7 and Remark 6 carry the Appendix B anisotropy remark and the Appendix J diagnostic only). The lemmas are stated below without their derivations, which the released notes supply.

\begin{itemize}
\item \textbf{Lemma 1 (Kantorovich--Rubinstein duality; H1).} For any $\theta$, $|R_T(\theta)-R_S(\theta)| \le L\,W_1(\mu_S,\mu_T)$.
\item \textbf{Lemma 2 (gradient-domain KR; H2).} $\|\nabla R_T(\theta)-\nabla R_S(\theta)\| \le L'\,W_1(\mu_S,\mu_T)$. \emph{(The interchange $\nabla R = \mathbb E[\nabla\ell]$ requires measurability of $\nabla\ell(\theta,\cdot)$, integrability and the usual dominated-convergence conditions along the differentiating directions; H2 is a statement about $\nabla\ell$ and presupposes its differentiability in $\theta$ rather than implying it.)}
\item \textbf{Lemma 3 (optimal-displacement bound; H1 and H5 in its ball form, together with the localization of $\theta^*_S$ --- H8 does not belong to this footprint).} $\|\theta^*_S-\theta^*_T\| \le \sqrt{4LW_1/\mu_T}$. \emph{(Two corrections of record. Constant: $\sqrt{4LW_1/\mu_T}$ follows from the \textbf{linear} branch, $\text{(4)}\le2LW_1$, with the two-point strong-convexity inequality, $\tfrac{\mu_T}{2}\|\Delta\|^2\le\text{(4)}\le2LW_1$; routing the same inequality through the single-KR bound $\text{(4)}\le LW_1$ gives the sharper $\sqrt{2LW_1/\mu_T}$, and the looser form is printed because it matches H10's normalisation $W_1\le\mu_T r_T^2/4L$, so H10 and this lemma use one radius scale rather than two. Hypothesis list: the derivation applies the two-point inequality to the pair $(\theta^*_S,\theta^*_T)$, so it \textbf{presupposes} the localization stated in Appendix J and Appendix C step 1 and cannot be used to establish it --- no circularity is available. It uses no Taylor identity and no $C^2$ regularity, so H8 is not in the footprint.)}
\item \textbf{Lemma 4 (second-order expansion; H8, not H3 --- label corrected).} $R_T(\theta^*_S)-R_T(\theta^*_T) = \nabla R_T(\theta^*_S)^\top\Delta - \tfrac12 \Delta^\top H_T(\xi)\Delta$ with $\Delta := \theta^*_S-\theta^*_T$ and $\xi$ on the segment between $\theta^*_S$ and $\theta^*_T$. Expanding about $\theta^*_S$ and evaluating at $\theta^*_T = \theta^*_S - \Delta$ makes the second-order term carry a minus sign; the opposite sign, $+\tfrac12$, is wrong. The correction is \textbf{load-bearing}: the theorem now bounds $(4)$ itself rather than $|(4)|$, and the sign of the quadratic term decides whether a curvature term is needed (Appendix C step 4). \emph{(Label corrected in the same pass: the form printed is an intermediate-point identity, which needs $C^2$ regularity (H8); $\beta$-smoothness (H3) yields an integral-form remainder instead of a value at some $\xi$. The lemma is unused by Corollary 1 and Proposition 1 in any case, which use the defining strong-convexity inequality rather than a Taylor expansion --- Appendix C step 3.)}
\item \textbf{Lemma 5 (flatness--spectrum relation; not used by Corollary 1 or Proposition 1 --- it is the measurement note of \S{}4.2, and belongs to the flatness protocol and the diagnostic of Appendix J).} At a stationary point ($\|\nabla R_D(\theta)\|\approx0$): $\max_{\|\epsilon\|\le\rho} R_D(\theta+\epsilon)-R_D(\theta) = \tfrac12\rho^2\lambda_{\max}(H_D(\theta)) + o(\rho^2)$; at non-stationary points the read carries a first-order term $\rho\|\nabla R_D(\theta)\|$.
\item \textbf{Remark 6 (former Lemma 6; PAC-Bayes source complexity).} The source-complexity term $C_2S_S$ is derivable via a PAC-Bayes argument but is deliberately not part of the theorem (\S{}4.2, first omission).
\item \textbf{Lemma 7 (curvature transfer; H4; retained for the anisotropic remark of Appendix B and the diagnostic of Appendix J --- not used by the theorem).} $\lambda_{\max}(H_T(\xi)) \le \lambda_{\max}(H_T(\theta)) + \beta_H(\|\xi-\theta\|)$, hence the target--source spectral excess at matched points is bounded by $L''W_1$ (Lemma 2 raised one derivative order).
\item \textbf{Lemma 8 (optimisation error to parameter distance; H9).} $\|\theta-\theta^*_S\| \le \sqrt{2\varepsilon^{\mathrm{pop}}_{\mathrm{opt}}/\mu_S}$.
\end{itemize}

\textbf{Small-shift condition (H10).} H10, $W_1\le\mu_T r_T^2/(4L)$, is a \emph{compatibility} condition between the shift scale and the growth scale of H5; it does \textbf{not} by itself establish that $\theta^*_S$ lies in the strong-convexity ball (a distant source well only $O(LW_1)$ worse is not excluded), so localization is a separate assumption (Appendix C step 1) with a basin-separation certificate as a check. The displacement bound does not depend on the $\alpha'$ bookkeeping: the sharp branch $L'^2W_1^2/(2\mu_T)$ follows from the strong-convexity inequality and the linear branch $2LW_1$ from Lemma 1 alone, and the $W_1^{3/2}$ estimate the $\alpha'=1/2$ bookkeeping produced is withdrawn as dominated.

\section*{Appendix A-bis. The measurement-hierarchy proposition: derivation, regimes and falsification}
\addcontentsline{toc}{section}{Appendix A-bis. The measurement-hierarchy proposition: derivation, regimes and falsification}

\begin{quote}
\textbf{Why this appendix exists.} The article states the proposition, the two regimes and the falsification condition, and points here for the derivation: the second-order expansion and its assumption footprint, the two thresholds and what each licenses, the placement argument, and what the statement does \textbf{not} claim.
\end{quote}

\textbf{The proposition, stated so that it can be refuted.} Write $R_D$ for the risk of the loss on domain $D$ and $H_D$ for its Hessian in the parameters. Every quantity in the family is read by perturbing a released checkpoint: at a parameter point $\theta$, along a unit direction $u$ and at the \textbf{relative radius} the protocol fixes ($\epsilon=\rho\|\theta\|u$, Appendix M.1), the worst-direction increment is

\[\Delta_D(\theta,\rho,u)\;:=\;R_D(\theta+\rho\|\theta\|\,u)-R_D(\theta)\;=\;\underbrace{\rho\|\theta\|\,\langle\nabla R_D(\theta),u\rangle}_{\text{first order,}\;\propto\;\rho}\;+\;\underbrace{\tfrac12\rho^2\|\theta\|^2\,u^\top H_D(\theta+\tau\rho\|\theta\|u)u}_{\text{second order,}\;\propto\;\rho^2},\qquad \tau\in(0,1),\]

the ordinary second-order expansion to an intermediate point --- an identity needing $C^2$ regularity in a neighbourhood of $\theta$ and a bounded Hessian there (\textbf{H8} and \textbf{H6} in the notation of Appendix A) and no Lipschitz constant of the loss in the data, i.e. \textbf{not H1}, which is what Corollary 1 uses and which plays no part here. Bounding the Hessian at the intermediate point by its value at the base point is a step we do \textbf{not} take: the ratio of the two terms is stated at the base point, and it is one number,

\[r_D(\theta,\rho)\;:=\;\frac{\|\nabla R_D(\theta)\|}{\|H_D(\theta)\|\,\rho\|\theta\|},\]

the local slope in units of the local curvature scale at the operational radius. \textbf{In the aligned case with a constant Hessian and $u^\top H u>0$} (the conditions under which the main text states the same ratio, \S{}6), if $r_D\le\tfrac1{16}$ the quadratic term is at least eight times the linear one, so the read is curvature-dominated with exponent $2$ in $\rho$; if $r_D\ge1$ the linear term is at least twice the quadratic one, so the read is gradient-dominated with exponent $1$, the two being equal at $r_D=\tfrac12$. The regimes are separated by a band, not a single threshold, and \textbf{the operational consequence is what matters}: a perturbation-based diagnostic is readable as a curvature value only where $r_D\le\tfrac1{16}$; outside, the local slope is comparable with or larger than the curvature term.

\textbf{Statement.} Under \textbf{H6} and \textbf{H8} of Appendix A --- a bounded parameter domain and $C^2$ regularity, in a parameter ball where the Hessian is bounded in operator norm and the protocol radius is small against that bound --- the read $\Delta_D$ is curvature-dominated only where $r_D(\theta,\rho)\le\tfrac1{16}$, and where $r_D$ is larger the read is dominated by, or comparable with, the local slope; so a finite-SGD checkpoint with a non-negligible residual gradient will generally produce a slope-dominated read, while the direction of comparison at any particular endpoint is a matter for measurement. \textbf{We measure $r_D$ in both forms, and the two disagree, so the statement is relative to the form asked}: the pointwise ratio at the released endpoints has median \textbf{0.028} (87\% below $1/16$, none above $1$), placing those cells in the curvature-dominated regime, while the path form has median \textbf{10.8} and places them in the slope-dominated one (Appendix U). The statement is therefore about the \textbf{$\varepsilon$-scaling channel} it is derived for; it does not adjudicate between the two forms, and where they disagree the regime label belongs to the form named rather than to the cell. What it contributes is the \textbf{derivation of the two regimes and the judgement of which regime the \ensuremath{\varepsilon}-scaling reading can speak to}, not a measurement of the boundary between them. We also do not claim that every fine-tuned endpoint here lies in one regime: non-stationarity is not by itself dominance, and the placement is reported as a measurement rather than argued from the checkpoint's provenance.

Three of our readings bear on that statement, reported as evidence about the placement rather than as illustrations of it. The \textbf{\ensuremath{\varepsilon}-scaling exponents of 0.56--0.81 (median 0.59)} are read at the two customary radii: well below the $\approx2$ curvature domination would require, so curvature domination is excluded there --- with the qualification that these are effective exponents over a finite radius interval rather than the asymptotic exponent of either branch, and that at the third radius (0.01) the increments depart from the power law outright, which is why the protocol fixes the primary radius at 0.001. \textbf{That exclusion is the whole of what the first reading establishes, and it is a direct test of the regime}: an effective exponent near $2$ would place those cells in the curvature branch and would falsify the placement on its own. The \textbf{SNR floor (18/18 cells $<2$, band 0.76--1.47)} is a second, independent obstacle of a different kind: a failure of the gate for \emph{using} the diagnostic as an outcome predictor, not evidence about first- against second-order, and the two must not be pooled. The \textbf{evaluator and reporting sensitivity} of \S{}6(b) is a third and again independent obstacle: what a slope-dominated read does when the reference it is ratioed against is itself a trajectory quantity --- the ranking flips with the evaluator, and in the single cleanest instance the sign of the reported gain flips with the \textbf{reporting column alone}, with runs, seeds and fine-tuning untouched. What we do not claim is impossibility: a larger gradient batch or a variance-reduced estimator changes the residual gradient $r_D$ is built from, and with it the regime (Appendix H.1); the statement is about this protocol and this estimator.

\section*{Appendix B. Anisotropic curvature form}
\addcontentsline{toc}{section}{Appendix B. Anisotropic curvature form}

\textbf{This appendix is retained as a record but licenses no reported number.} It sets out the anisotropic curvature form used by the earlier derivation of the bound, and \textbf{the article does not cite it}. The derivation it served is superseded by the treatment in Appendix J, which states the conditions the bound actually uses. \textbf{No claim in the paper depends on this appendix}, and nothing in it is used to license a number.

\section*{Appendix C. Term-(4) sensitivity: the sharp bound, the withdrawn interaction, and the algebra}
\addcontentsline{toc}{section}{Appendix C. Term-(4) sensitivity: the sharp bound, the withdrawn interaction, and the algebra}

Term (4) is the target-domain optimal displacement cost $R_T(\theta^*_S)-R_T(\theta^*_T)$.

\textbf{The withdrawn interaction, recorded.} An earlier version of the bound carried an interaction term between term (4) and the curvature. \textbf{That term is withdrawn}, and no statement in this paper relies on it. We record the withdrawal here because a reader who saw the earlier form is entitled to know which part of it no longer holds.

\textbf{Sharpness is not claimed.} The bound was described as sharp contingent on a family of constructions that the earlier algebra took for granted. \textbf{We no longer claim sharpness}, because the construction that would attain the bound is not exhibited: a point-mass family was referred to but never written down, and we will not assert an attainment we cannot show.

\textbf{A correction to the algebra.} Where the derivation writes the operator norm of the Hessian, a \textbf{spectral radius} is required: Appendix G records a negative smallest eigenvalue at \textbf{every} measured endpoint, and at such a point the two quantities differ. The earlier form is therefore not correct as written; the correction is stated here rather than silently applied.

\textbf{No number in the paper depends on this appendix} --- it is cited zero times in the article --- so nothing above licenses or withdraws a reported value. It is kept as a record of what was claimed and what is no longer claimed.

\section*{Appendix D. Pipeline and seed protocol rules}
\addcontentsline{toc}{section}{Appendix D. Pipeline and seed protocol rules}

Dual-layer seed mechanism: the framework's \texttt{seed} argument is recorded in args.yaml, but the data-pipeline RNG does not consume it on these datasets (bit-identical replay); a \texttt{--shuffle-seed} CLI flag (deterministic permutation of the train split file order, validation untouched) restores seed validity --- same value reproducible, different values provably different (controlled test). args.yaml does \textbf{not} record shuffle-seed; per-run registrations, with \texttt{train\_obj} version numbers, are archived with this release. Pipeline-version rule: old/new pipeline mAP50-95 offsets measured 0.08--1.12 pp; every gain statement is same-pipeline paired; cross-pipeline numbers appear only in discovery narratives.

\textbf{Augmentation-RNG control (second seed layer).} A third seed layer exists and was quantified rather than assumed away: the released stack fixes the DataLoader generator (\texttt{manual\_seed(6148914691236517205 + RANK)}), so with \texttt{--shuffle-seed} fixed the \emph{augmentation} stream is bit-identical across runs and the registered \ensuremath{\sigma} sees data order only. A \texttt{--aug-seed} flag re-seeds the \textbf{worker-process} RNGs (random/numpy/torch) while leaving the generator, sampler, split-order permutation and initialisation untouched. Applied to SHWD\ensuremath{\rightarrow}SFCHD at 100 epochs (20\% budget, \texttt{--shuffle-seed 42} fixed), three augmentation seeds give baseline-arm test mAP50-95 of 43.810 / 43.930 / 43.800 (SD 0.072 pp) and strategy-arm (lr0.005) values of 44.500 / 44.350 / 44.020 (SD 0.246 pp), i.e. gains +0.690 / +0.420 / +0.220 pp against the matched same-budget baseline --- all positive, mean +0.443 pp. The two arms ran on 4090s under ultralytics 8.4.120 \citep{bjorck2025} / torch 2.4.0+cu121 with identical training code and a verified-equal dataset file list; the anchor configuration reproduced the test endpoint to four decimals (mAP50 0.7371, mAP50-95 0.4450, \ensuremath{\Delta} = 0.000 pp), the comparability check for this control. Released: \texttt{a5\_train\_obj\_aug.py} (trainer with \texttt{--aug-seed}), per-run \texttt{results.csv}/\texttt{args.yaml}.

\section*{Appendix E. Reversal list}
\addcontentsline{toc}{section}{Appendix E. Reversal list}

\textbf{Table T7.} The reversal list --- five conclusions reversed by prospectively frozen rules and one further result invalidated for same-value duplication --- with the layer that caught each.

{\footnotesize
\setlength{\tabcolsep}{4pt}
\begin{longtable}{p{0.022\textwidth}p{0.180\textwidth}p{0.353\textwidth}p{0.198\textwidth}p{0.141\textwidth}}
\toprule
\# & conclusion reversed or invalidated & the reading it replaced & disposition & caught by \\
\midrule
\endhead
1 & pipeline-pairing artefact & a "+0.80 pp steady-state gain" manufactured by a base-side pipeline effect of +1.12 pp & cell retired as an artefact & internal re-analysis \\
2 & headline magnitude & quoted in turn as +0.62 pp \ensuremath{\rightarrow} +0.40 pp \ensuremath{\rightarrow} +0.45 pp \ensuremath{\rightarrow} \textbf{+0.51 pp}, as the base was re-run and paired & corrected & internal re-analysis \\
3 & alignment rank correlation & +1.00 & \textbf{\ensuremath{-}1.00} --- the sign flips with the evaluator (\S{}7) & internal re-analysis \\
4 & loss-prior "positive tendency" & a single positive reading & dissolved into three zero/negative multi-seed estimates & internal re-analysis \\
5 & historical module degradation & \ensuremath{-}10 pp & exposed as a silent mis-construction of the networks & checkpoint audit during internal review \\
6 & same-value duplication (invalidated, not reversed) & a dead-layer insertion producing weights identical to the baseline & result invalidated & blinded external audit pass \\
\botrule
\end{longtable}}

\subsection*{E.1 The reversal list, triggers and defensive rules}
\addcontentsline{toc}{subsection}{E.1 The reversal list, triggers and defensive rules}

A \textbf{second seed-replay instance is established}: three \texttt{mendeley} baseline runs nominally at seeds 42/43/44 carry bit-identical epoch sequences (and one further run is a re-labelled copy), i.e. the \texttt{--shuffle-seed} injection had not taken effect for them. The screening rule the first replay introduced --- compare per-run epoch sequences by hash before forming any pair --- is mandatory, and the affected runs were re-run under corrected seeds rather than overwritten.

\textbf{What the audit layer is.} A \textbf{blinded external audit pass} over the draft: the full draft, the grading table and the reference list are submitted to an independently hosted reviewing system with no project background and no access to the repository, and its unedited responses are released. It is \emph{external} in that the reviewing system is not author-controlled, and \emph{not independent} in the sense of an unaffiliated human auditor --- the authors chose the material and wrote the prompt, so its findings are advisory. Two assumption-level failures (the retracted family-control claim and the Appendix J.2 erratum) and the dead-layer duplication were established through this pass, and every finding was adjudicated by the authors, with rejections as well as acceptances. The process therefore mandates, as its fourth rule, a "key assumptions \ensuremath{\times} plausible ranges" table at pre-registration time (Appendix G). The reversals are the record of what the tests caught --- evidence about our \emph{process}, not about the validity of the survivors: the evidence for the surviving claims is \S{}4.6 and \S{}8, with its limits stated (the density caveats of \S{}4.6 and \S{}9).

\textbf{One correction is mathematical rather than empirical, and is a result of the same kind.} The interaction term $C_3W_1\lambda_{\max}(H_T(\theta))$ has been \textbf{removed} (Appendix C step 5): term (4) is non-negative by target optimality, so once H5 is read as genuine local strong convexity the curvature contribution is slack. The objection that forced the correction is that the printed chain bounds the curvature magnitude by $\lambda_{\max}$, not a valid modulus for $|\Delta^\top H\Delta|$ whenever a negative eigenvalue exceeds $\lambda_{\max}$ in magnitude --- precisely where a curvature term would be doing work. The diagnostic $\kappa^-_\Delta$ is retained, and the assumption the correction turns on (H5 as the two-point inequality, not one-point growth) is stated in Appendix G.

\textbf{Status of the surviving claims.} Three statements, each checkable against the released files. \emph{(i) All cited geometry is canonical-protocol re-measurement, not legacy data.} No legacy flatness reading enters any claim: the same checkpoint reads 5.19\ensuremath{\times}10\textsuperscript{5} under the legacy protocol and 342.3 under the canonical one, and legacy values appear only inside the retraction of Appendix J.2. \emph{(ii) All outcome claims are post-repair, same-pipeline and paired.} Every gain statement in \S{}8 comes from runs made after the \texttt{--shuffle-seed} repair, paired against a same-family, same-budget baseline; cross-pipeline readings survive only in discovery narratives and are marked as such. The pipeline offset is bounded empirically: re-running the 20\% label fraction under the current pipeline reproduced two registered single runs to within 0.01 pp on the baseline and exactly on the strategy arm (0.4396 and 0.4426 against the registered 0.4397 and 0.4426; the ten-seed cell means are 0.4386 and 0.4440) --- a single-run reproduction of two runs, not of the cell means. \emph{(iii) No surviving claim depends on a withdrawn cell.} Of the events above, one cell was retired as a pipeline-pairing artefact, one magnitude was corrected, and one cell is retained only as a retraction. This is provenance, not completeness: we can show the survivors were re-measured under the repaired protocol, and cannot show that protocol is free of defects we have not yet imagined --- which is why every surviving claim is stated together with the measurement that could overturn it.

\subsection*{E.2 The selection premium: per-cell restatements and the cells the coincidence affects}
\addcontentsline{toc}{subsection}{E.2 The selection premium: per-cell restatements and the cells the coincidence affects}

\textbf{The premium was measured per cell; the measurement and its direction are reported in the companion paper on evaluation validity rather than here} --- that paper owns the per-cell and per-arm values, the endpoint restatements, and the enumeration of which cells the \texttt{val}/\texttt{test} coincidence affects. What this paper retains is the contrast that conditions its headline numbers: the published-protocol magnitudes \textbf{+1.76 pp} (smoke\ensuremath{\rightarrow}SFCHD) and \textbf{+0.54 pp} (SHWD\ensuremath{\rightarrow}SFCHD) are best-checkpoint levels on the split that selected them, not held-out test estimates. \textbf{The clean three-way protocol and its disclosures.} A three-way protocol (label subset / validation carve-out from the remaining training pool / untouched test split) removes that structure, and we report it rather than merely proposing it: on the clean protocol both headline cells hold --- smoke\ensuremath{\rightarrow}SFCHD \textbf{+1.434 pp at ten seeds} (baseline 42.062 \ensuremath{\rightarrow} strategy 43.496; SD 0.383 pp; paired t = \textbf{11.84}; sign-flip permutation p = \textbf{0.0020}, the n = 10 floor, all ten seeds positive) and SHWD\ensuremath{\rightarrow}SFCHD \textbf{+0.507 pp} (43.559 \ensuremath{\rightarrow} 44.066; SD 0.488 pp; t = \textbf{3.29}; p = \textbf{0.0098}, nine of ten positive) --- and these, not the published-protocol values of \S{}4.4, are the paper's primary readings of the two headline cells. \textbf{The eleven- and thirteen-seed extension, on which the family-level status rests:} smoke\ensuremath{\rightarrow}SFCHD \textbf{+1.409 pp} (baseline 42.108 \ensuremath{\rightarrow} strategy 43.517; SD 0.372 pp; paired t = \textbf{12.55}; permutation p = \textbf{0.000977}, the n = 11 floor, eleven of eleven positive) and SHWD\ensuremath{\rightarrow}SFCHD \textbf{+0.628 pp} (43.484 \ensuremath{\rightarrow} 44.112; SD 0.489 pp; t = \textbf{4.63}; p = \textbf{0.001221}, twelve of thirteen positive). \textbf{The extension moved the point estimate by \ensuremath{-}0.025 pp and +0.121 pp on the two cells} --- it tightened the interval rather than raising the effect. \textbf{At ten seeds neither clean-tier test clears the paper's own roster-wise correction} (q = 0.052 and 0.127); \textbf{extended to eleven and thirteen seeds both clear it (q = 0.0159)}, and that extension was undertaken with the ten-seed values known, so this correction cannot control the FDR of the family that generated the selection. \textbf{Fresh-seed check.} Restricting the extension to the seven seeds that were \emph{not} part of the screening reading leaves the three gains intact: +1.75 / +0.55 / +3.60 pp for smoke\ensuremath{\rightarrow}SFCHD / SHWD\ensuremath{\rightarrow}SFCHD / visdrone\ensuremath{\rightarrow}dota15 (paired t = 26.50 / 9.50 / 48.71; p = 1.9\ensuremath{\times}10\textsuperscript{-7} / 7.8\ensuremath{\times}10\textsuperscript{-5} / 5.0\ensuremath{\times}10\textsuperscript{-9}; strategy-side effects 9.3\ensuremath{\sigma} / 3.0\ensuremath{\sigma} / 26.6\ensuremath{\sigma}; exact 7-vs-7 two-sample permutation p = 5.8\ensuremath{\times}10\textsuperscript{-4} --- the complete-separation value, with paired sign-flip floor $2/2^7=0.016$ --- in all three), so the replication does not rest on the three seeds it reuses and the corroborating cell is carried by fresh data alone (released script \texttt{fresh7\_seed\_analysis.py}). Within the clean protocol the estimate is also \textbf{not monotone in the seed count}: +1.562 / +0.693 pp at six seeds, +1.420 / +0.492 at nine, +1.434 / +0.507 at ten, the direction selection inflation predicts, which is why the ten-seed value is the one reported. The selection premium itself is decomposed in \textbf{Appendix N.2}.

\begin{quote}
\textbf{A note on the lettering.} Appendix E carries E.1 and E.2; the surviving subsections keep their original numbers so that the article's and this file's cross-references stay stable.
\end{quote}

\begin{quote}
\textbf{The assigner-level out-of-memory behaviour of the training framework, declared once.} On a dense batch the task-aligned assigner catches a CUDA out-of-memory error, falls back for that batch, and in a few runs the effective batch size is briefly reduced before training continues --- a \textbf{documented property of Ultralytics rather than a deviation from our protocol}. Because such an event touches an isolated batch, and the reading a run reports is a \textbf{maximum over epochs}, its effect is \textbf{inside the noise and error range}: a full-log scan bounds it at a median \textbf{\ensuremath{-}0.0006 pp} over 20 cells against a data-order $\hat{\sigma}$ of \textbf{0.166--0.190 pp}; the origin (process concurrency) was removed at source. This is the only place the subject is declared.
\end{quote}

\section*{Appendix F. Number freeze table (v16) and sensitivity}
\addcontentsline{toc}{section}{Appendix F. Number freeze table (v16) and sensitivity}

\textbf{Null exceedance of the frozen rule (derivation).} Write the pass condition as $\overline{\Delta} > 3\hat\sigma_{\rm strategy}$ with $n = 3$ paired seeds. Under the null $\overline{\Delta} \sim N(0, \sigma_d^2/n)$, and $\hat\sigma_{\rm strategy}/\sigma_{\rm strategy} = s$ is independent of $\overline{\Delta}$ (mean and variance of a normal sample are independent), with $2s^2 \sim \chi^2_2$, i.e. density $f(s) = 2s e^{-s^2}$. Hence, with $q = \sigma_d/\sigma_{\rm strategy}$, \[P_{\rm null}(q) = \int_0^\infty 2s e^{-s^2}\Big[1 - \Phi\big(3\sqrt{3}\,s/q\big)\Big] ds,\] whose two limits matter more than the interior values, and neither is small: $P_{\rm null}(q) \to \tfrac12$ as $q \to \infty$ --- when the paired noise swamps the strategy-side $\hat{\sigma}$ the threshold becomes negligible and the rule degenerates towards a coin flip --- and $P_{\rm null}(q) \approx q^2/54$ as $q \to 0$ (numerically $1.66\times10^{-3}$ at $q = 0.3$ against $q^2/54 = 1.67\times10^{-3}$). \emph{The rate is set by the \ensuremath{\sigma} ratio, not by the rule}, and for large-$q$ cells the rule is close to vacuous, which is why the secondary statistics rather than the criterion carry the inference there. Numerically (closed form and a paired-arm Monte Carlo with $1.5\times10^6$ draws agree to three decimals; released script \texttt{null\_exceedance.py}):

{\footnotesize
\setlength{\tabcolsep}{4pt}
\begin{longtable}{p{0.348\textwidth}p{0.051\textwidth}p{0.051\textwidth}p{0.051\textwidth}p{0.051\textwidth}p{0.051\textwidth}p{0.051\textwidth}p{0.051\textwidth}p{0.051\textwidth}p{0.051\textwidth}}
\toprule
q & 0.6 & 0.8 & 0.9 & 1.0 & 1.2 & 1.4 & 1.6 & 2.0 & 2.5 \\
\midrule
\endhead
per-cell null exceedance & 0.65\% & 1.14\% & 1.44\% & 1.75\% & 2.47\% & 3.28\% & 4.16\% & 6.08\% & 8.66\% \\
expected null passes in 26 configurations & 0.17 & 0.30 & 0.37 & 0.46 & 0.64 & 0.85 & 1.08 & 1.58 & 2.25 \\
P(at least one null pass) & 15.7\% & 25.9\% & 31.3\% & 36.9\% & 47.8\% & 57.9\% & 66.9\% & 80.4\% & 90.5\% \\
\botrule
\end{longtable}}

The family rows are computed on the 26-configuration enumerated roster used everywhere else in this paper (Appendix I), not on the recorded count of 28; the quoted range (0.37--1.65 expected passes, 31--82\%) is \textbf{derived from the null model rather than measured}, and spans the q values the graded cells actually realise Two corrections of record: the family rows use the 26-cell denominator for both family rows (0.40--2.42; 33.3--92.1\%), which contradicted the roster rule stated in the main text (Appendix I); and its q = 0.6 column of the family-wise row read 9.7\% where the closed form and the Monte Carlo both give 16.8\% on 28 cells (15.7\% on the 26-configuration roster). The per-cell column was unaffected and reproduces to three decimals against a paired-arm Monte Carlo with $1.5\times10^6$ draws. \textbf{The rule therefore does not by itself control family-wise error}: at the observed q values the expected number of null passes among the enumerated family cells is 0.37--1.65 and the probability of at least one is 31--82\%. The count distribution matters more than its mean, because two cells is what we observe: under the global null, and on the enumerated roster, the probability of \emph{at least two} passes is 5.3\% at q = 0.9, 7.6\% at q = 1.0, 20.9\% at q = 1.4 and 49.7\% at q = 2.05 (one or more: 31 / 37 / 58 / 82\%). A two-pass screening outcome is therefore not by itself evidence against the null, and we do not present it as such: \emph{given} that these two arms were the ones taken forward, the probability that both would clear their paired test at ten seeds under the null \textbf{if the two cells' tests were independent} is $6.4\times10^{-10} \times 6.2\times10^{-6} = 4\times10^{-15}$ --- and they are not independent, since the two cells share the target corpus, the target label subset and the pipeline, so this product is an upper bound on the evidence rather than a calibrated level. The family is therefore \emph{not} protected by the rule at any plausible q, and the tier structure (ten-seed replication extension, permutation tests, FDR robustness read of \S{}9.3) carries the inferential weight.

Canonical anchors as listed in the grading table v16 (\S{} Anchors). Sensitivity: the corroborating cell's strategy-side $\hat{\sigma}$ is 0.166 pp at ten seeds (0.13 pp at the registered three seeds; $\hat{\sigma}$ CI [0.69\ensuremath{\times}, 1.83\ensuremath{\times}] at n = 10), and the strong cell's is 0.190 pp (CI [0.131, 0.347]). At the ten-seed $\hat{\sigma}$ upper bound the corroborating cell falls to 1.8\ensuremath{\sigma} on the strategy-side rule while its paired test (t = 9.35, p = 6.2\ensuremath{\times}10\textsuperscript{-6}) and exact permutation p (1.1\ensuremath{\times}10\textsuperscript{-5}) are unaffected --- both conventions are reported for every passing cell rather than averaged. \textbf{The \ensuremath{\sigma}-convention, in full here because the main text carries it compressed (\S{}9.3).} The paired-difference \ensuremath{\sigma} measures the contrast's total noise and is the scale an NHST decision on the paired difference would use; the strategy-side \ensuremath{\sigma} measures whether the \emph{strategy's} gain replicates across seeds, which is the reproducibility question this paper asks --- so the criterion's frozen choice follows the question rather than the test. At three seeds the strategy-side $\hat{\sigma}$ CI is [0.5\ensuremath{\times}, 6.3\ensuremath{\times}], at which both cells fall below the 3\ensuremath{\sigma} line (disclosed rather than averaged away), and the ten-seed extension tightens it to [0.69\ensuremath{\times}, 1.83\ensuremath{\times}]. Assumptions\ensuremath{\times}ranges table: see Appendix G. Additional anchors from the same campaign: dose--response 30 ep +0.86/+0.87/0.00/\ensuremath{-}0.76 pp and 100 ep +0.07/+0.29/+0.53/\ensuremath{-}0.38 pp at 2.5\ensuremath{\times}/5\ensuremath{\times}/10\ensuremath{\times}/20\ensuremath{\times}; label-budget baselines 0.4128/0.4397/0.4502/0.4632 with strategy gains +0.51/+0.29/+0.70/+1.25 pp at 10/20/30/50\%; controls cosine +0.12/+0.16 pp, frozen backbone \ensuremath{-}4.94/\ensuremath{-}3.95 pp, AdamW +0.33/+0.31 pp; 200-episode baseline 0.4385 (early-stopped) vs lr0.005 +0.52 pp; within-domain dota15 +1.12/+0.95/+0.93 pp; pre-training far-field reads (\ensuremath{\rho} = 0.01): SHWD\ensuremath{\rightarrow}SFCHD 7,334, MAFA\ensuremath{\rightarrow}mendeley 2,534, mendeley\ensuremath{\rightarrow}mendeley 323, dota 297, visdrone 283, firesmoke 207, mask 59,248.

\subsection*{F.1 Multiplicity: units, p-value set and sensitivity to the family definition}
\addcontentsline{toc}{subsection}{F.1 Multiplicity: units, p-value set and sensitivity to the family definition}

\textbf{The 3-vs-3 defence chain (recomputed, not transcribed).} Where a tier is carried by three seeds per arm, the complete-separation event and its null probability are reported together, so the separation cannot be read as stronger than it is: the corroborating cell's corroborating tier has complete separation \textbf{0.21 pp}, and under permutation the 3-vs-3 complete-separation null probability is \textbf{$2/\binom{6}{3}=2/20=0.10$ two-sided} (\textbf{0.05 one-sided}) --- the same two-sided convention as the paired sign-flip floors. Complete separation at three seeds is therefore a weak event on its own, and the tiers that carry inference are the paired tests.

\textbf{Seed-\ensuremath{\sigma} composition.} The augmentation-RNG component has since been \textbf{measured on the strong cell as well} (smoke\ensuremath{\rightarrow}SFCHD, 100 epochs, 20\% budget, \texttt{--shuffle-seed 42} fixed, \texttt{--aug-seed} 101/202/303): the baseline arm reads 0.4230 / 0.4250 / 0.4270 and the strategy arm 0.4410 / 0.4410 / 0.4370 on the target test split, i.e. gains of \textbf{+1.80 / +1.60 / +1.00 pp (mean +1.47 pp, 3/3 positive, paired t = 6.10)}, with endpoint SDs of 0.20 pp (baseline) and 0.23 pp (strategy). That is the \emph{same order} as the corroborating cell's augmentation component, so on the strong cell the gain is not an artefact of the augmentation stream either. Three components have been separated by direct measurement on the same cell (SHWD\ensuremath{\rightarrow}SFCHD, 100 epochs, 20\% subset). (i) \emph{Training-split order}, which is what the registered \ensuremath{\sigma} varies (0.18--0.21 pp paired-difference, \S{}8.2). (ii) \emph{Augmentation RNG}, measured with three worker-process seeds while the DataLoader generator, the split-order seed and the initialisation are held fixed: strategy-arm SD \textbf{0.246 pp}, baseline-arm SD \textbf{0.072 pp}, i.e. 1.07--1.89\ensuremath{\times} the registered \ensuremath{\sigma}, which makes the 3\ensuremath{\sigma} rule \emph{more permissive} than its nominal calibration; the strong cell's margin absorbs this while the corroborating cell clears the registered threshold and not the augmentation-inflated one, so that cell is carried by its conventional statistics (\S{}8.2, \S{}8.3). The rule is \textbf{not} recalibrated for this --- the shortfall (1.07--1.89\ensuremath{\times}) sits below the 3\ensuremath{\times} trigger set before the measurement (Appendix G) --- so the two statements are complementary: the line stands, \emph{and} this one cell is carried by its conventional statistics instead of by it. (iii) \emph{Initialisation}, measured by training the source endpoint from scratch three times and re-running both arms from each: gain SD \textbf{0.146 pp} (3/3 positive, mean +0.437 pp) and absolute-endpoint SD \textbf{0.209 pp}. \textbf{Seed-count convention for this appendix: every component above is estimated at n = 3}, and where a $\hat{\sigma}$ is also quoted at ten seeds (\S{}8.2) the two estimates are not interchangeable. The \textbf{largest is the augmentation component} (0.246 pp), with the data-order component next (0.18--0.21 pp) and initialisation smallest (0.146 pp); none of the three is large enough to manufacture an effect of \ensuremath{\approx}1.4--1.8 pp. The controls cover one cell, one pair and one budget. \textbf{The consequence for the two cells, stated rather than averaged:} the strong cell is unaffected --- its +1.76 pp gain clears 3$\hat{\sigma}$ = 0.57 pp on the registered ten-seed $\hat{\sigma}$ = 0.190 pp and 3 \ensuremath{\times} 0.246 = 0.74 pp on the augmentation-inflated \ensuremath{\sigma} --- while the corroborating cell is \emph{not}: its +0.54 pp gain clears 3$\hat{\sigma}$ = 0.50 pp on the registered data-order $\hat{\sigma}$ = 0.166 pp and falls below 3 \ensuremath{\times} 0.246 = 0.74 pp once the measured augmentation component is included. What carries that cell is therefore its \emph{conventional} statistics --- paired t = 9.35 (p = 6.2\ensuremath{\times}10\textsuperscript{-6}) at ten seeds and +0.55 pp on the seven fresh seeds (p = 7.8\ensuremath{\times}10\textsuperscript{-5}), neither of which uses \ensuremath{\sigma} --- and its criterion pass is reported as conditional on the data-order \ensuremath{\sigma} convention.

\textbf{Family-level analysis: the three family definitions}, reported as robustness rather than as a family-wise guarantee --- paired tests with BH-FDR (q = 0.05) over the 26-configuration roster, given three ways because the family choice is not innocent and because the roster was enumerated after the outcomes. (i) The two extended \emph{family} arms, the remaining 24 roster configurations entering at p = 1 (m = 26, the Appendix I roster): q = 1.7\ensuremath{\times}10\textsuperscript{-8} (smoke\ensuremath{\rightarrow}SFCHD), 8.1\ensuremath{\times}10\textsuperscript{-5} (SHWD\ensuremath{\rightarrow}SFCHD). (ii) The three extended arms alone (m = 3): q = 6.2\ensuremath{\times}10\textsuperscript{-6} and 9.6\ensuremath{\times}10\textsuperscript{-10} for the two family arms. The family-external visdrone\ensuremath{\rightarrow}dota15 arm is not one of the two extended \emph{family} arms --- "family member" here means a member of that extended-arm set, not a row of the 26-configuration roster (its configuration is row 4 of Appendix I) --- so its raw p = 1.7\ensuremath{\times}10\textsuperscript{-13} is the primary number we report; including it in that m = 3 completion would give q = 5.0\ensuremath{\times}10\textsuperscript{-13}. (iii) The two family arms without completion (m = 2, listed p-values): q = 1.3\ensuremath{\times}10\textsuperscript{-9} (smoke\ensuremath{\rightarrow}SFCHD) and 6.2\ensuremath{\times}10\textsuperscript{-6} (SHWD\ensuremath{\rightarrow}SFCHD). Every family-level value survives by orders of magnitude. It is nonetheless a \emph{robustness read, not a family guarantee}: the extension was triggered by the screening outcome, so correcting within the selected arms cannot control the FDR of the family that generated the selection, and the conventional analysis does not rescue the registered tier either (BH q = 0.29 and 1.00 at three seeds) --- hence the frozen criterion is called the \emph{screening} instrument and the tier structure carries the inference. At three seeds the same procedure is not decisive (paired p = 0.011 does not survive the correction; Welch p = 0.0008 does). The two $\hat{\sigma}$ conventions also diverge at ten seeds: the strong cell remains at 5.1\ensuremath{\sigma} while the corroborating cell falls to 1.8\ensuremath{\sigma} on the strategy-side rule (its 10-vs-10 two-sample permutation p = 1.1\ensuremath{\times}10\textsuperscript{-5}, paired sign-flip floor $2/2^{10}=2.0\times10^{-3}$).

\textbf{What the multiplicity family is.} The BH correction is applied to \emph{configurations}, not to arms: the unit is one (source endpoint \ensuremath{\rightarrow} target corpus, budget) configuration, tested by the paired contrast between its registered strategy arm and its baseline, and the other 24 configurations enter with p = 1 because they carry either a single arm or no paired seeds --- an after-the-fact sensitivity convention rather than a prospectively defined family. The p-value set entering the correction is therefore two values, 0.011 and 0.081 at three seeds (and 6.4\ensuremath{\times}10\textsuperscript{-10} / 6.2\ensuremath{\times}10\textsuperscript{-6} at ten seeds; 1.9\ensuremath{\times}10\textsuperscript{-7} / 7.8\ensuremath{\times}10\textsuperscript{-5} on the fresh seven). Sensitivity to the family definition: over the \textbf{26-configuration roster} (Appendix I) the ten-seed values are q = 1.7\ensuremath{\times}10\textsuperscript{-8} / 8.1\ensuremath{\times}10\textsuperscript{-5}; over the \textbf{recorded registry count of 28} they are 1.8\ensuremath{\times}10\textsuperscript{-8} / 8.7\ensuremath{\times}10\textsuperscript{-5}; and over the \textbf{arm-level count of 138 contrasts} in the roster they are 8.8\ensuremath{\times}10\textsuperscript{-8} / 4.3\ensuremath{\times}10\textsuperscript{-4} --- every version clears by orders of magnitude, and none rescues the registered three-seed tier (q = 0.29 / 0.31 / 1.00 across the three definitions). All three variants are recomputed under one stated convention --- two real paired tests, every other configuration entering at p = 1 --- by the released script \texttt{recompute\_family\_sensitivity.py}. \textbf{The clean-protocol tier under the same convention, stated because it does \emph{not} pass:} the three-way replication of \S{}8.6 gives the two headline cells p = 0.0020 and 0.0098; under the identical roster-wise convention (m = 26, every other configuration at p = 1) those become \textbf{q = 0.052} and \textbf{q = 0.127} at ten seeds --- \textbf{neither clears q = 0.05}, the first only because 26\ensuremath{\times}(2/2\textsuperscript{10}) = 0.0508, i.e. because the permutation floor is coarse at n = 10 rather than because the evidence is weak --- and at the eleven- and thirteen-seed extension \textbf{q = 0.0159} for both, which clears q = 0.05.

\textbf{Table T8.} The \ensuremath{\sigma} components behind the tier vocabulary, per cell (data-order seeds, n = 3; percentage points). \ensuremath{\sigma} is a (strategy \ensuremath{\times} domain-pair) property, computed on the strategy side unless a row is labelled otherwise.

{\footnotesize
\setlength{\tabcolsep}{4pt}
\begin{longtable}{p{0.385\textwidth}p{0.226\textwidth}p{0.319\textwidth}}
\toprule
cell (strategy \ensuremath{\times} domain pair) & \ensuremath{\sigma}\_b, baseline arm & \ensuremath{\sigma}\_strategy, strategy arm \\
\midrule
\endhead
SHWD\ensuremath{\rightarrow}SFCHD, lr0.005 & 0.18 & 0.13 \\
smoke\ensuremath{\rightarrow}SFCHD, lr0.005 & 0.25 & 0.23 \\
MAFA\ensuremath{\rightarrow}mendeley, sns & --- & 0.26 \\
\botrule
\end{longtable}}

The n = 3 sd estimation CI is \textbf{[0.5\ensuremath{\times}, 6.3\ensuremath{\times}]}. The two components the registered design excludes are reported separately in Appendix G: augmentation-RNG 0.246 pp (strategy arm) and 0.072 pp (baseline arm); initialisation 0.146 pp (gain SD) and 0.209 pp (absolute endpoint SD).

\subsection*{F.2 The \ensuremath{\sigma} components in full}
\addcontentsline{toc}{subsection}{F.2 The \ensuremath{\sigma} components in full}

\emph{\ensuremath{\sigma} controls (the components the registered design excludes).} Three components are separated by direct measurement on the same cell (SHWD\ensuremath{\rightarrow}SFCHD, 100 epochs, 20\% subset): data order (0.18--0.21 pp paired-difference), which is what the registered \ensuremath{\sigma} varies; augmentation RNG (strategy-arm \textbf{lr0.005} SD \textbf{0.246 pp}, baseline-arm SD \textbf{0.072 pp}, gains \textbf{+0.69 / +0.42 / +0.22 pp}, mean \textbf{+0.443 pp}, 3/3 positive); and initialisation (\textbf{SD 0.146 pp} across three from-scratch source endpoints, gains \textbf{+0.570 / +0.280 / +0.460 pp}, mean \textbf{+0.437 pp}, 3/3 positive), the absolute endpoint moving with SD \textbf{0.209 pp}. Initialisation is the \textbf{smallest} of the three, so it does not widen the \ensuremath{\sigma} bound, though on the corroborating cell the augmentation component does mean that it clears the registered threshold and \textbf{not} the augmentation-inflated one and is therefore carried by its conventional statistics; the full composition, with every sample size, is in \textbf{Appendix F.1}.

\subsection*{F.3 The family-level reads}
\addcontentsline{toc}{subsection}{F.3 The family-level reads}

\textbf{Family-level analysis.} The Benjamini--Hochberg correction over the 26-configuration roster is reported as \emph{robustness, not a family-wise guarantee}; its p-values, family sizes and the sensitivity table are in \textbf{Appendix F.1}.

\section*{Appendix G. Key assumptions \ensuremath{\times} plausible ranges (the fourth rule)}
\addcontentsline{toc}{section}{Appendix G. Key assumptions \ensuremath{\times} plausible ranges (the fourth rule)}

\textbf{The \ensuremath{\varepsilon}-scaling exponents behind \S{}7.2(1), per cell (target-side worst direction, customary radii).} Fitting $S(\rho) \propto \rho^{k}$ on the two customary radii (0.001 \ensuremath{\rightarrow} 0.003) of the released target-domain flatness readings gives: 0.66 (\texttt{sio\_a\_base\_100ep}), 0.56 (\texttt{sio\_a\_jps\_100ep}), 0.56 (\texttt{sio\_a\_jps\_30ep}), 0.59 (\texttt{sio\_a\_pws\_100ep}), 0.63 (\texttt{sio\_css\_30ep}), 0.58 (\texttt{sio\_pws\_30ep}), 0.81 (\texttt{sio\_sns\_30ep}) --- median 0.59, all well below the \ensuremath{\approx}2 a curvature-dominated read would require; the two pathological \texttt{spi} arms (a deliberate negative control) give 4.57 and 4.31, and at the third radius (0.01) the increments diverge by up to seven orders of magnitude, which is why the primary radius is fixed at 0.001. Source-side readings have only the 0.001 radius in the release, so no exponent is claimed for them.

\textbf{What was measured, and what was not.} The hierarchy is a statement about the quantities \emph{this study measured}: trajectory-evaluated first-order (gradient) reads and same-protocol relative diagnostics. A spectral read at the fine-tuned endpoints now exists (\S{}4.2 and this appendix: Lanczos with 20 iterations \citep{hernandez2021} over 220 endpoint pairs, Ritz residuals reported), but it is a \emph{stochastic estimate at a computed checkpoint} --- batch-draw dependent, reported as a distribution and used only as a diagnostic --- not a population spectrum. What it establishes is negative and about the assumption rather than the outcome: the \emph{ball} form of H5 is refuted at every measured endpoint ($\lambda_{\min}<0$ in 220/220), so a curvature estimate from these artefacts would not supply the modulus the removed interaction term needed. Still unmeasured, and not measurable from these artefacts: $L$, $L'$, $W_1$, $r_T$ and the localization placement assumed in \S{}4.2. The claim is that \emph{these} diagnostics are not outcome predictors, not that spectral quantities would fail as well; the stationary-side reads we do report (source-side spectral sharpness, \S{}8.6) are computed by perturbation, not from a spectrum, and are deliberately not used as predictors.

{\footnotesize
\setlength{\tabcolsep}{4pt}
\begin{longtable}{p{0.095\textwidth}p{0.352\textwidth}p{0.483\textwidth}}
\toprule
Assumption & Registered/plausible range & Sensitivity statement \\
\midrule
\endhead
L (loss Lipschitz const.) & order-dependent; not measured & enters C\textsubscript{1} only (additive), no grading impact \\
\ensuremath{\mu}\_T (target strong-convexity) & \textbf{measured at the checkpoints, and the ball form is refuted there}: $\lambda_{\min}(H_T)<0$ in \textbf{220 of 220} endpoint pairs; the ordered-pair form is what remains testable, and its secant modulus $\mu_{sec}$ is measured and changes sign across configurations (\textbf{75 negative / 145 positive}) & \textbf{the reading matters} --- two-point strong convexity, not one-point growth; it scales the quadratic branch of Proposition 1 and $T_*$, and it is what makes a curvature term redundant under H5 while the weaker one-point reading does not (Appendix C step 5). Because the ball form fails at every measured endpoint, \S{}4.2 states the quadratic branch under the \textbf{ordered-pair} form rather than under a ball \\
r\_T (H5 radius) & not measured, and at the measured endpoints it is moot: a ball carrying H5's content requires $\lambda_{\min}$ bounded below by a positive $\mu_T$, and $\lambda_{\min}<0$ holds in 220 of 220 endpoint pairs, so no such ball exists there & H10 compatibility requires r\_T \ensuremath{\geq} \ensuremath{\surd}(4LW\textsubscript{1}/\ensuremath{\mu}\_T); at observed D this is a mid-range requirement \\
H10 small shift & D = 9.4--13.0 observed (W\textsubscript{1} proxy) & \textbf{a compatibility check only}: it relates the shift scale to the growth scale of H5, does \emph{not} establish localization (Appendix C step 1), and is not a hypothesis of Corollary 1 or Proposition 1 \\
\ensuremath{\alpha}\ensuremath{'} (displacement exponent) & \emph{withdrawn} --- this row belonged to the generic-displacement estimate $2L'\sqrt{L/\mu_T}W_1^{3/2}$, which Appendix J.4 removes as never the lower envelope of the printed bound; a historical entry & the printed displacement bound uses the strong-convexity quadratic branch and Lemma 1's linear branch, and contains no \ensuremath{\alpha}\ensuremath{'} \\
\ensuremath{\rho} (perturbation radius) & 0.001 primary (protocol fact) & first-order dominated at target; three-tier decomposition provided \\
\ensuremath{\sigma}\_b, \ensuremath{\sigma}\_strategy & measured per Table T8 (n = 3) & headline-cell sensitivity table in \S{}9.3 \\
augmentation-RNG component (the \ensuremath{\sigma} component the registered design excludes) & \textbf{measured}: strategy arm 0.246 pp, baseline arm 0.072 pp (n = 3 augmentation seeds, SHWD\ensuremath{\rightarrow}SFCHD 100 ep) & \ensuremath{\sigma} is a lower bound by 1.07--1.89\ensuremath{\times} of the data-order \ensuremath{\sigma}; below the 3\ensuremath{\times} recalibration trigger for the 3\ensuremath{\sigma} line and the +0.4 pp module threshold (\S{}8.2--\S{}8.3, Appendix D) --- so no recalibration is triggered, and the one cell whose margin was \ensuremath{\approx}1.1\ensuremath{\times} the line is instead carried by its conventional statistics (Appendix F.1) \\
initialisation component (the other \ensuremath{\sigma} component the registered design excludes) & \textbf{measured}: gain SD 0.146 pp and absolute-endpoint SD 0.209 pp across 3 independent source trainings (3/3 positive gains, mean +0.437 pp; SHWD\ensuremath{\rightarrow}SFCHD 100 ep) & the smallest of the three components: below the augmentation component (0.246 pp) and below the across-seed spread of the same clean-protocol cell (0.49 pp, \S{}8.6), so it does not widen the \ensuremath{\sigma} bound; the endpoint movement (43.250--43.630\%) is why every comparison here is a same-source-endpoint paired contrast \\
\botrule
\end{longtable}}

\textbf{Which result uses which assumption.} Because the results no longer contain a curvature term and no longer use a Taylor expansion, their assumption footprint is much smaller than it would otherwise be, and the two branches of Proposition 1 have \emph{different} footprints.

{\footnotesize
\setlength{\tabcolsep}{4pt}
\begin{longtable}{p{0.234\textwidth}p{0.064\textwidth}p{0.074\textwidth}p{0.154\textwidth}p{0.115\textwidth}p{0.055\textwidth}p{0.164\textwidth}}
\toprule
Assumption & Corollary 1 (gap) & Prop. 1, linear branch & Prop. 1, quadratic branch & Eq. (8) schematic (\S{}6) & Lemma 8 (\S{}5) & $\kappa^-_\Delta$ diagnostic (App. J.2) \\
\midrule
\endhead
H1 ($L$) & \textbf{yes} --- Lemma 1 & \textbf{yes} --- Lemma 1 twice & no &  &  &  \\
H2 ($L'$) & no & no & \textbf{yes} --- Lemma 2 &  &  &  \\
H3 ($\beta$-smoothness) & no & no & no &  &  &  \\
H4 ($M$) & no & no & no &  &  & yes --- transferring curvature from $\xi$ to $\theta$ \\
H5 ($\mu_T,r_T$; \textbf{two-point}, and only for the ordered pair) & no & no & \textbf{yes} --- the defining first-order inequality &  &  &  \\
H6 (bounded parameter domain) & no & no & no &  &  &  \\
H7 (isotropic curvature) & no & no & no & \textbf{yes} --- the isotropic special case &  &  \\
H8 ($C^2$ regularity) & no & no & no & yes &  & yes \\
H9 ($\mu_S$) & no & no & no &  & \textbf{yes} --- Lemma 8 &  \\
H10 (small shift) & no & no & no --- \textbf{a compatibility check, not a hypothesis} &  &  &  \\
\textbf{localization} $\theta^*_S\in B(\theta^*_T,r_T)$ & no & no & \textbf{yes} (assumption) &  &  &  \\
\textbf{source optimality} (used to drop the middle telescoping term) & no & \textbf{yes} & no &  &  &  \\
\textbf{source stationarity} $\nabla R_S(\theta^*_S)=0$, or the directional VI & no & no & \textbf{yes} (assumption) &  &  &  \\
\textbf{non-negative loss} $\ell\ge0$ & \textbf{yes} (assumption) & no & no &  &  &  \\
\botrule
\end{longtable}}

Four consequences. (i) \textbf{Corollary 1 needs H1 and $\ell\ge0$ only} --- no strong convexity, no localization, no stationarity, no displacement machinery, and (if the gap is defined via an infimum) not even existence of $\theta^*_T$. (ii) \textbf{The two branches of Proposition 1 are not the same result}: the linear branch $(4)\le2LW_1$ uses H1 and source-optimality alone, while the quadratic branch uses H2, H5, localization and stationarity and does \textbf{not} use H1. (iii) \textbf{H10 is not a hypothesis of any result here}: it relates the shift scale to the growth scale of H5 and does not establish localization, and the earlier claim that it did is withdrawn. (iv) \textbf{H6 is likewise not a logical assumption} of the abstract statements: a bounded domain is one way to make the Lipschitz constants uniform over $\Theta$ (which the "for any $\theta\in\Theta$" quantifier wants) and one way to make the sharpness family satisfy H1, but neither is a proof step. H3 and H8 are simply not needed by these results, since the displacement bound uses the strong-convexity inequality rather than a Taylor expansion; they remain in the list because other remarks invoke them.

\textbf{What the endpoint read suite measures, and what it leaves open.} \textbf{Protocol.} 220 endpoint pairs over \textbf{119 endpoints} --- the 99 endpoints representing the frozen roster (one baseline and one strategy arm per configuration, plus the full ten-seed arm of the two outcome-significant cells) and the twenty endpoints of the \texttt{dota15\ensuremath{\rightarrow}aitod} replication cell --- each read on 1, 2 or 4 batch draws of 4 batches \ensuremath{\times} 4 images at 640 px, with the loss evaluated with the model in training mode and every normalisation layer in evaluation mode, Lanczos with 20 iterations at the target endpoint, and the source objective evaluated wherever the source and target classification heads are compatible. Every quantity is an empirical value \textbf{at a computed checkpoint}, not at a population minimiser. Released scripts: \texttt{theory\_read\_endpoint.py} (the read) and \texttt{analyze\_reads.py} (the aggregation); the \ensuremath{\varepsilon}-scaling and encoder-Jacobian tables of \S{}7.2(1) are released alongside them as \texttt{theory\_e2.csv} and \texttt{theory\_e3.csv}. One correction belongs with the protocol, because it changed numbers: in the first form of the read suite the loss was evaluated with the model in \texttt{eval()} mode, where the detection head runs NMS and the value is not the training risk. The correction changed the E2 table and left the forward-only E3 table byte-identical, and the pre-correction tables are retained in the release rather than deleted.

{\footnotesize
\setlength{\tabcolsep}{4pt}
\begin{longtable}{p{0.254\textwidth}p{0.223\textwidth}p{0.453\textwidth}}
\toprule
Quantity & Status & Measured over the 220 pairs \\
\midrule
\endhead
$\lambda_{\min}(H_T(\theta_T))$ & measured, \textbf{diagnostic only} & \textbf{negative in 220/220} \\
$\lambda_{\max}(H_T(\theta_T))$ & measured, \textbf{diagnostic only} & $5.2\times10^3$ to $1.1\times10^6$; not a modulus for $\|\Delta^{\top}H_T\Delta\|$ (\textbf{Appendix J.4}) \\
Ritz residuals of both extremes & measured (convergence check) & relative median $4.6\times10^{-4}$ and $6.1\times10^{-4}$; 90th percentile $2.7\times10^{-3}$ and $3.5\times10^{-3}$ \\
$\|\Delta\|$ & measured & median 4.5 (range 0.39--508) \\
$\widehat{(4)}=R_T(\theta_S)-R_T(\theta_T)$ & measured & \textbf{negative in 26/220} --- so it is reported as a checkpoint loss difference, not as a non-negative cost \\
$\mu_{sec}$, the ordered-pair secant modulus & measured & \textbf{75 negative / 145 positive}, median 2.19 \\
source optimality $R_S(\theta_S)-R_S(\theta_T)$ & measured on the 164 head-compatible pairs & holds in \textbf{59/164} \\
source directional VI $\langle\nabla R_S(\theta_S),\Delta\rangle$ & same subset & holds in \textbf{32/164}; both hold together in \textbf{14/164} \\
$L$, $L'$, $W_1$, $r_T$, localization & \textbf{not measured, and not measurable from these artefacts} & Eqs. (4)--(5) remain scaling statements \\
\botrule
\end{longtable}}

Three consequences. (i) \textbf{The ball form of H5 is refuted at every endpoint measured}: $\lambda_{\min}<0$ in 220 of 220 pairs, so no ball around a fine-tuned endpoint carries H5's content; what remains testable is the \textbf{ordered-pair} inequality, which is what $\mu_{sec}$ measures. (ii) \textbf{Both source-side premises fail more often than they hold}, the empirical reason the $\varepsilon$-forms of \S{}4.2 and Appendix C steps 2--5 exist: they replace the two exact conditions with measured residual terms, and are therefore not cosmetic. (iii) \textbf{A quarter of the pairs cannot form the $\varepsilon$-certificate at all} (56/220), because the source checkpoint's classification head has a different class count and is re-initialised, so $\nabla R_S$ is not a gradient in the target's parameter space --- the limitation \S{}4.2 states. Two caveats on the whole table: the Hessian quantities are batch-draw dependent (the relative spread of $\lambda_{\max}$ across draws of one endpoint has median 0.52, of $\widehat{(4)}$ 0.26 and of $\mu_{sec}$ 0.36), so they are reported as distributions over draws rather than point values; and the sign of the source-side quantities is not draw-stable in every case (it flips across draws of the same endpoint in 14 of 51 multi-draw endpoints for source optimality and 11 of 51 for the directional VI), so those are reported as rates with that instability disclosed rather than as per-cell verdicts.

\subsection*{G.1 Assumption plausibility under detection losses}
\addcontentsline{toc}{subsection}{G.1 Assumption plausibility under detection losses}

\textbf{The full argument the main text summarises.} The loss family's \textbf{data-Lipschitz constants} are finite and explicit: the box terms are $[0,1]$-valued and piecewise smooth in box coordinates that are themselves Lipschitz in the image through the regression head (needing the usual non-degeneracy of box denominators), and the classification term satisfies $L_{\rm cls}\le\mathrm{Lip}(u(\theta;\cdot))$, since $|\partial\,\mathrm{BCE}(\sigma(u))/\partial u|=|\sigma(u)-y|\le1$ and composition multiplies Lipschitz constants --- so no logit clipping is required for H1 (the loss \emph{value} is unbounded; its \emph{sensitivity to the data} is not, which is what H1 needs). H6 enters here rather than in the derivation: it is what makes these constants \textbf{uniform over the parameter domain}, which the theorem's ``for any $\theta\in\Theta$'' quantifier requires; smoothness (H3) and Hessian-Lipschitz regularity (H4) hold for standard convolutional backbones away from non-differentiable points; \textbf{H5 is not supported at the measured endpoints in its ball form} --- $\lambda_{\min}(H_T)<0$ in 220 of 220 endpoint pairs --- and the bound's quadratic branch uses the \textbf{ordered-pair form} instead, whose secant modulus $\mu_{\rm sec}$ is measured and changes sign across configurations; and H10's compatibility at the observed shift range is an assumption-level statement, tabulated with its plausible ranges above.

\subsection*{G.2 The \ensuremath{\varepsilon}-scaling reading in full}
\addcontentsline{toc}{subsection}{G.2 The \ensuremath{\varepsilon}-scaling reading in full}

(1) \ensuremath{\varepsilon}-scaling: at customary radii the target-domain worst-direction increment is first-order dominated --- the read is $\rho\|\nabla R_T\|$, a slope/distance proxy, not curvature --- and that is \textbf{measured}: the two-point slope of the worst-direction increment against \ensuremath{\rho} at the two customary radii (0.001 \ensuremath{\rightarrow} 0.003) gives exponents of \textbf{0.56--0.81 (median 0.59)}, where the third radius (0.01) diverges by up to seven orders of magnitude --- which is why the protocol fixes the primary radius at 0.001; the two deliberately pathological \texttt{spi} arms are the exception at 4.3--4.6. The per-radius perturbation readings and the encoder-Jacobian spectral widths behind this reading are released as theory\_e2.csv and theory\_e3.csv, recomputed after the evaluation-mode correction recorded above. (2) The cos-proxy flip chain (\S{}6.2). (3) The SNR floor (18/18 \textless{} 2). (4) Gap-vs-flatness linear correlations vanish under the strict protocol. Contrast class: source-side stationary spectral reads hold at 253-284, \ensuremath{\pm}6\% --- 12 worst-direction readings at $\epsilon = 0.001$ on the source side, spanning 252.99--283.96, a relative spread of \ensuremath{\pm}5.8\% --- and source stationary positions are insensitive to the source-training loss (stage ablation). Unit note: source-side reads are worst-direction increments at a \emph{relative} radius \ensuremath{\rho}\ensuremath{\|}\ensuremath{\theta}\ensuremath{\|} (\ensuremath{\|}\ensuremath{\theta}\ensuremath{\|} = 146.2 at the pretrained origin and 148--158 across the fine-tuned endpoints, Appendix M.1) expressed in \emph{loss} units; radius and reading are in different units and must not be compared numerically, and \ensuremath{\rho} is given explicitly wherever a read is quoted.

\section*{Appendix H. cos-proxy protocol details}
\addcontentsline{toc}{section}{Appendix H. cos-proxy protocol details}

Six protocols (v1 test-distribution; v3 train-distribution quasi-stationary; v3b frozen-reference variant of v3; v3c multi-checkpoint; v3d best-neighbour; v3e plateau + SNR gate) --- \S{}6.2 lists the five that enter the SNR table, and v3b appears only in the flip chain, 18 cells each (3 pairs \ensuremath{\times} 2 budgets \ensuremath{\times} 3 strategies), SNR = \ensuremath{\|}d\_\ensuremath{\pi}\ensuremath{\|}/\ensuremath{\|}g\_ref\ensuremath{\|} per cell, gate \ensuremath{\geq} 2 --- the gate asks the signal to exceed the residual reference by a factor large enough that sign, not merely magnitude, is resolvable. The full per-cell table is printed in H.1; the \textbf{flip chain itself (v3b vs v3c/d/e) is not released}, because the three per-point values each cell is built from are not among the released files.

\subsection*{H.1 The cos-proxy protocol chain in full}
\addcontentsline{toc}{subsection}{H.1 The cos-proxy protocol chain in full}

The five protocols that enter the SNR table form the progressively corrected chain below, and the sixth --- v3b, the frozen-reference variant of v3 --- appears only in the flip chain. The five: v1 (test-distribution gradients --- invalidated: at a mismatched point $\mathbb{E}_{\text{test}}[\nabla\ell_0]\neq0$ and the difference gradient degenerates to $d_\pi \approx (s{-}1)g_{\mathrm{ref}}$, $s<1$); v3 (train-distribution quasi-stationary evaluation, where the stationary condition cancels the reference term); v3c (multi-checkpoint averaging); v3d (best-neighbour averaging); v3e (plateau averaging with an SNR gate). Outcome: across all 18 measurement cells (three domain pairs \ensuremath{\times} two budgets \ensuremath{\times} three strategies, every cell under the v3e plateau-averaged protocol with its SNR gate; the v3c/v3d variants each also produce 18 cells) the difference gradient is the same order as the residual reference gradient (\textbf{SNR \ensuremath{\in} 0.76--1.47, all below the prospectively frozen gate of 2}), and \textbf{the ranking is not stable across averaging protocols, and the reversal we once recorded does not reproduce on the protocols whose logs survive.} Pairing the two available protocols --- best-checkpoint (\texttt{A\_v3}) against multi-checkpoint (\texttt{A\_pooled}) --- over the three strategies per evaluation point gives Spearman \textbf{+0.5, +1.0, +0.5, +0.5, \ensuremath{-}0.5, \ensuremath{-}0.5} across the six points (mean \textbf{+0.25}): \textbf{no point reverses}, so the earlier claim of +1.00 against \ensuremath{-}1.00 is \textbf{withdrawn}, and the three values per point are released (\texttt{provenance/appendix\_H/H1\_flip\_chain\_three\_values.csv}). The \texttt{v3d} best-neighbour values the reversal rested on are \textbf{not in this archive}. What survives is the weaker statement that the ranking moves with the averaging choice. The SNR floor holds on all 18 cells and depends on no ordering, but it is \textbf{supporting rather than decisive}: the gate sits only 1.36x above the upper edge of a band whose total width is 1.94x.

\textbf{The eighteen cells, one row each} (protocol \textbf{v3e}: plateau mean over the last five saved epochs, \texttt{SNR = \textbar{}d\_pi\textbar{}/\textbar{}g0\textbar{}}; \texttt{A\_pooled} is the pooled alignment estimate with its bootstrap interval, \texttt{frac\_pos} the fraction of per-batch estimates that are positive).

{\footnotesize
\setlength{\tabcolsep}{4pt}
\begin{longtable}{p{0.222\textwidth}p{0.119\textwidth}>{\raggedleft\arraybackslash}p{0.089\textwidth}>{\raggedleft\arraybackslash}p{0.119\textwidth}>{\raggedleft\arraybackslash}p{0.104\textwidth}>{\raggedleft\arraybackslash}p{0.089\textwidth}>{\raggedleft\arraybackslash}p{0.119\textwidth}}
\toprule
pair & strategy & SNR & A\_pooled & CI\_lo & CI\_hi & frac\_pos \\
\midrule
\endhead
shwd2sf@b30 & css & 1.469 & 0.6767 & \ensuremath{-}0.3540 & 0.5224 & 0.71 \\
shwd2sf@b100 & css & 1.422 & 0.4686 & \ensuremath{-}0.2288 & 0.4251 & 0.73 \\
smoke\ensuremath{\rightarrow}fs@b100 & sns & 1.419 & \ensuremath{-}0.0691 & \ensuremath{-}0.0766 & 0.0346 & 0.25 \\
shwd2sf@b30 & sns & 1.408 & 0.6638 & \ensuremath{-}0.3129 & 0.5427 & 0.65 \\
smoke\ensuremath{\rightarrow}fs@b30 & sns & 1.396 & \ensuremath{-}0.0788 & \ensuremath{-}0.1172 & 0.0734 & 0.42 \\
smoke\ensuremath{\rightarrow}fs@b30 & css & 1.339 & 0.0743 & \ensuremath{-}0.0764 & 0.1209 & 0.69 \\
shwd2sf@b100 & sns & 1.327 & 0.4952 & \ensuremath{-}0.1495 & 0.4551 & 0.65 \\
smoke\ensuremath{\rightarrow}fs@b30 & pws & 1.284 & 0.2407 & \ensuremath{-}0.0141 & 0.1420 & 0.94 \\
smoke\ensuremath{\rightarrow}fs@b100 & css & 1.270 & \ensuremath{-}0.0091 & \ensuremath{-}0.0557 & 0.0475 & 0.40 \\
smoke\ensuremath{\rightarrow}fs@b100 & pws & 1.263 & 0.0574 & \ensuremath{-}0.0253 & 0.0638 & 0.54 \\
shwd2sf@b100 & pws & 1.209 & 0.5065 & \ensuremath{-}0.1461 & 0.4228 & 0.65 \\
mafa\ensuremath{\rightarrow}mende@b30 & css & 1.163 & 0.0971 & \ensuremath{-}0.0968 & 0.1077 & 0.69 \\
shwd2sf@b30 & pws & 1.149 & 0.6978 & \ensuremath{-}0.3293 & 0.5573 & 0.69 \\
mafa\ensuremath{\rightarrow}mende@b30 & pws & 1.036 & 0.0687 & \ensuremath{-}0.0853 & 0.1064 & 0.50 \\
mafa\ensuremath{\rightarrow}mende@b30 & sns & 0.997 & 0.0615 & \ensuremath{-}0.0888 & 0.1021 & 0.50 \\
mafa\ensuremath{\rightarrow}mende@b100 & css & 0.774 & \ensuremath{-}0.0820 & \ensuremath{-}0.1363 & 0.1344 & 0.31 \\
mafa\ensuremath{\rightarrow}mende@b100 & sns & 0.759 & \ensuremath{-}0.1126 & \ensuremath{-}0.1625 & 0.1186 & 0.39 \\
mafa\ensuremath{\rightarrow}mende@b100 & pws & 0.758 & \ensuremath{-}0.0695 & \ensuremath{-}0.1245 & 0.1699 & 0.42 \\
\botrule
\end{longtable}}

\textbf{Why the floor is supporting and not decisive, stated on the table above.} All eighteen SNRs lie below the gate (the largest, 1.469, by a factor 1.36), so the gate marks the whole family as unusable as an outcome predictor here. But the intervals in the last three columns show why that cannot carry the sign question: on \texttt{shwd2sf@b30}/\texttt{css} the pooled estimate is \textbf{+0.677} while the interval is \textbf{[\ensuremath{-}0.354, +0.522]} --- \textbf{because the two are different estimators}: \texttt{A\_pooled} aggregates the batch estimates, whereas the interval is the bootstrap range of the \textbf{per-batch} estimate \texttt{A\_perbatch} (0.145 on that row, inside the interval). Seven of the eighteen rows have the pooled value above the per-batch interval's upper edge; the pooled aggregate and the per-batch interval are not the same object, and on most rows the interval is as wide as the estimate. The floor is therefore a \textbf{bound on the diagnostic's usability}, not a resolution of the sign.

Mechanistically, at a \emph{precise} stationary point $d_\pi$ is well-defined; finite-SGD training only reaches approximate stationarity, so the SNR stays O(1) in our runs --- an empirical obstruction under finite-SGD, finite-batch training rather than an impossibility theorem (larger gradient batches, variance-reduced or fully-converged estimators are outside this protocol).

\section*{Appendix I. The registered roster, enumerated cell-first (26 configurations, two pipelines)}
\addcontentsline{toc}{section}{Appendix I. The registered roster, enumerated cell-first (26 configurations, two pipelines)}

\textbf{The registered family: what we can and cannot reconstruct.} The family was recorded as 28 strategy--pair cells constructed as "lr0.005 and the tested loss families \ensuremath{\times} the seven domain pairs \ensuremath{\times} budget tiers" (the released grading table, family entry), and 28 \ensuremath{\times} 3 = 84 graded runs are reported. Re-deriving the roster from the released material yields \textbf{26 pair-able configurations}: 22 from the B pipeline and 4 from the SFCHD pipeline, whose runs postdate the registration snapshot and which carry \textbf{both} criterion-passing cells. The first enumeration of this appendix covered only the first 26 and therefore excluded the two cells the paper is about; that inconsistency was found in the external audit pass and is corrected here rather than argued away. The recorded count of 28 is reported as the registration record and is used as a denominator nowhere --- and \textbf{this roster is not the family used for the multiple-comparison correction}: the correction of \S{}5 runs over the \textbf{three registered pairs} (m = 3), a different object from the 26/28-configuration roster, so no reported statistic divides by either number; the two constructions we can check --- seven pairs \ensuremath{\times} four arm families (= 28) and the enumerated registry (= 26) --- do not agree, and we cannot recover retrospectively which 28 cells the frozen count referred to, because the registration documents describe the family \emph{by construction} rather than as an enumerated list. This is a defect in the registration record, and we state its consequences rather than absorb them: (i) "registered family" here means the set of cells covered by the frozen criterion, not an independently auditable list; (ii) the family-level quantities --- the expected number of null passes (0.37--1.65), the family-wise probability of at least one (31--82\%), and the FDR family size --- would inherit that uncertainty if computed on an undefined roster, so they are computed on the 26-configuration roster of Table T10 and \textbf{m = 26 is used for every multiplicity calculation}; (iii) the discrepancy is stated wherever the roster is used as a denominator.

Table T10 is the roster that every family-level quantity in this paper is computed on. Its enumeration rule, stated once and applied to both released sources: group the released runs by (\textbf{source pretraining endpoint \ensuremath{\rightarrow} target corpus, budget}); a group containing a \texttt{base*} run (shapeiou @ lr0.001) \textbf{and} at least one strategy arm (\texttt{lr005}, \texttt{sns}, \texttt{css}, \texttt{pws}, \texttt{jps}) is one pair-able configuration; these five are the released \textbf{run-name} vocabulary, the three loss-prior families tested in this paper are \texttt{css}, \texttt{sns} and \texttt{pws}, and \texttt{jps} appears in one configuration's recorded arms only. Arms are identified from the released \textbf{run-name convention} rather than from \texttt{args.yaml}, because the loss prior is injected at runtime (\texttt{losses.patch\_loss}) and therefore never appears as a \texttt{loss:} field in the recorded arguments. Released script \texttt{enumerate\_family\_cellfirst.py} regenerates this table from the released run archive.

Two enumeration attempts were wrong, and we record both errors rather than the correction alone. The first enumerated only the released \emph{registration table} (135 runs), which contains \textbf{no SFCHD row}, and so omitted the two criterion-passing cells; the second attempted to repair that by appending four SFCHD groups but kept the mis-classification of loss-prior arms as baselines, which made several rows arithmetically unable to support the cells the paper reports (\S{}6.3's \texttt{sns} at 3v3 against a single baseline run is the clearest example). Both defects were found in the external audit pass, not by our own checks, and the roster below is the corrected form.

\textbf{Table T10.} The registered roster, enumerated cell-first (26 configurations).

{\footnotesize
\setlength{\tabcolsep}{4pt}
\begin{longtable}{p{0.038\textwidth}p{0.088\textwidth}p{0.176\textwidth}p{0.038\textwidth}p{0.082\textwidth}p{0.401\textwidth}p{0.038\textwidth}}
\toprule
\# & pipeline & source \ensuremath{\rightarrow} target & budget & baseline runs & strategy arms (runs[seeds]) & runs \\
\midrule
\endhead
1 & within-domain & aitod\_pretrain \ensuremath{\rightarrow} aitod & 30 ep & 1 & lr005\ensuremath{\times}1[ns]; pws\ensuremath{\times}1[ns]; sns\ensuremath{\times}1[ns] & 4 \\
2 & within-domain & dota\_pretrain \ensuremath{\rightarrow} dota & 30 ep & 1 & lr005\ensuremath{\times}1[ns]; sns\ensuremath{\times}1[ns] & 3 \\
3 & within-domain & dota15\_pretrain \ensuremath{\rightarrow} dota15 & 30 ep & 1 & lr005\ensuremath{\times}1[ns]; pws\ensuremath{\times}1[ns]; sns\ensuremath{\times}1[ns] & 4 \\
4 & cross-domain & visdrone\_pretrain \ensuremath{\rightarrow} dota15 & 100 ep & 10 & lr005\ensuremath{\times}10[42,43,44,45,46,47,48,49,50,51] & 20 \\
5 & within-domain & dota15\_pretrain \ensuremath{\rightarrow} dota15 & 100 ep & 1 & css\ensuremath{\times}1[ns]; lr005\ensuremath{\times}1[ns]; pws\ensuremath{\times}1[ns]; sns\ensuremath{\times}1[ns] & 5 \\
6 & within-domain & fsin\_pretrain \ensuremath{\rightarrow} firesmoke & 30 ep & 1 & css\ensuremath{\times}1[ns]; lr005\ensuremath{\times}1[ns]; pws\ensuremath{\times}1[ns]; sns\ensuremath{\times}1[ns] & 5 \\
7 & cross-domain & smoke\_pretrain \ensuremath{\rightarrow} firesmoke & 30 ep & 3 & css\ensuremath{\times}1[ns]; lr005\ensuremath{\times}3[43,44]; pws\ensuremath{\times}1[ns]; sns\ensuremath{\times}3[43,44] & 11 \\
8 & within-domain & fsin\_pretrain \ensuremath{\rightarrow} firesmoke & 50 ep & 1 & lr005\ensuremath{\times}1[ns] & 2 \\
9 & cross-domain & smoke\_pretrain \ensuremath{\rightarrow} firesmoke & 50 ep & 1 & lr005\ensuremath{\times}1[ns]; pws\ensuremath{\times}1[ns]; sns\ensuremath{\times}3[43,44] & 6 \\
10 & within-domain & fsin\_pretrain \ensuremath{\rightarrow} firesmoke & 100 ep & 1 & css\ensuremath{\times}1[ns]; pws\ensuremath{\times}1[ns]; sns\ensuremath{\times}1[ns] & 4 \\
11 & cross-domain & smoke\_pretrain \ensuremath{\rightarrow} firesmoke & 100 ep & 1 & css\ensuremath{\times}1[ns]; lr005\ensuremath{\times}2[ns]; pws\ensuremath{\times}1[ns]; sns\ensuremath{\times}3[43,44] & 8 \\
12 & cross-domain & mafa\_pretrain \ensuremath{\rightarrow} mask & 30 ep & 1 & lr005\ensuremath{\times}1[ns] & 2 \\
13 & within-domain & mask\_pretrain \ensuremath{\rightarrow} mask & 30 ep & 3 & css\ensuremath{\times}1[ns]; lr005\ensuremath{\times}1[ns]; pws\ensuremath{\times}1[ns]; sns\ensuremath{\times}3[43,44] & 9 \\
14 & within-domain & mask\_pretrain \ensuremath{\rightarrow} mask & 50 ep & 1 & lr005\ensuremath{\times}1[ns] & 2 \\
15 & within-domain & mask\_pretrain \ensuremath{\rightarrow} mask & 100 ep & 1 & css\ensuremath{\times}1[ns]; pws\ensuremath{\times}1[ns]; sns\ensuremath{\times}1[ns] & 4 \\
16 & cross-domain & mafa\_pretrain \ensuremath{\rightarrow} mende & 30 ep & 1 & css\ensuremath{\times}1[ns]; lr005\ensuremath{\times}1[ns]; pws\ensuremath{\times}1[ns]; sns\ensuremath{\times}1[ns] & 5 \\
17 & within-domain & mendein\_pretrain \ensuremath{\rightarrow} mende & 30 ep & 3 & css\ensuremath{\times}1[ns]; lr005\ensuremath{\times}1[ns]; pws\ensuremath{\times}1[ns]; sns\ensuremath{\times}1[ns] & 7 \\
18 & cross-domain & mafa\_pretrain \ensuremath{\rightarrow} mende & 50 ep & 1 & lr005\ensuremath{\times}1[ns]; pws\ensuremath{\times}1[ns]; sns\ensuremath{\times}1[ns] & 4 \\
19 & within-domain & mendein\_pretrain \ensuremath{\rightarrow} mende & 50 ep & 1 & lr005\ensuremath{\times}1[ns] & 2 \\
20 & cross-domain & mafa\_pretrain \ensuremath{\rightarrow} mende & 100 ep & 4 & css\ensuremath{\times}4[43,44,45]; jps\ensuremath{\times}2[44,45]; pws\ensuremath{\times}1[ns]; sns\ensuremath{\times}4[42,44,45] & 15 \\
21 & within-domain & mendein\_pretrain \ensuremath{\rightarrow} mende & 100 ep & 2 & css\ensuremath{\times}1[ns]; pws\ensuremath{\times}5[42,43,44,45]; sns\ensuremath{\times}1[ns] & 9 \\
22 & SFCHD-pipeline & shwd2sf\_pretrain \ensuremath{\rightarrow} sfchd & 30 ep & 4 & lr005\ensuremath{\times}4[42,43,44] & 8 \\
23 & SFCHD-pipeline & shwd2sf\_pretrain \ensuremath{\rightarrow} sfchd & 50 ep & 3 & lr005\ensuremath{\times}4[42,43,44] & 7 \\
24 & SFCHD-pipeline & shwd2sf\_pretrain \ensuremath{\rightarrow} sfchd & 100 ep & 12 & lr005\ensuremath{\times}13[42,43,44,45,46,47,48,49,50,51]; pws\ensuremath{\times}1[ns]; sns\ensuremath{\times}3[44,45] & 29 \\
25 & SFCHD-pipeline & smoke\_pretrain \ensuremath{\rightarrow} sfchd & 100 ep & 11 & lr005\ensuremath{\times}11[42,43,44,45,46,47,48,49,50,51] & 22 \\
26 & within-domain & visdrone\_pretrain \ensuremath{\rightarrow} visdrone & 30 ep & 1 & lr005\ensuremath{\times}1[ns]; pws\ensuremath{\times}1[ns]; sns\ensuremath{\times}1[ns] & 4 \\
\botrule
\end{longtable}}

Which cells this roster covers. The two criterion-passing cells are rows 24 (\texttt{SHWD\ensuremath{\rightarrow}SFCHD}, 100 epochs, corroborating) and 25 (\texttt{smoke\ensuremath{\rightarrow}SFCHD}, 100 epochs, strong), and rows 22--23 carry the shorter-budget arms that produce the budget curve of \S{}8.3. The loss-prior and module cells of \S{}6.3/\S{}8.4 are rows 20 and 21 (\texttt{MAFA\ensuremath{\rightarrow}mendeley}, \texttt{mendeley\ensuremath{\rightarrow}mendeley}), whose \texttt{sns}/\texttt{css}/\texttt{pws} arms each carry three or more seeds, and row 24 also carries the \texttt{sns} arms used in \S{}8.4's \texttt{SHWD\ensuremath{\rightarrow}SFCHD} negative result. The \texttt{visdrone\ensuremath{\rightarrow}dota15} directional probe is row 4 (ten seeds, both arms), and it is \emph{not} counted as a family member in the multiplicity arithmetic.

Four properties. (i) \textbf{The family size used for every multiplicity calculation is m = 26} --- the number of pair-able configurations enumerated above --- the same 26 used in the BH corrections of \S{}9.3 and in the family rows of Appendix F. The recorded registration count of \textbf{28} is reported as the registration record and used as a denominator nowhere; it differs from 26 because it was written as a construction ("seven pairs \ensuremath{\times} four arm families") rather than enumerated. (ii) Of the 26, \textbf{twenty-two are B-pipeline and four are SFCHD-pipeline}; the SFCHD pipeline is where both passing cells live, a concentration (15\% of the roster carrying 100\% of the positive result) stated here and in \S{}8.2. (iii) The configurations are \emph{configurations}, not arms: the "strategy arms" column lists three to five arm families per configuration, so the arm-level contrast count is larger than 26 --- which is why family-level quantities are stated per configuration. (iv) Rows 4, 20 and 24 carry multi-seed arms (ten, three-plus and ten seeds) because they host the pre-specified seed extension; the seed sets are printed per row.

\subsection*{I.1 Run counts for the released archive}
\addcontentsline{toc}{subsection}{I.1 Run counts for the released archive}

\textbf{The archive's run counts, taken from the released archive rather than the registration record.} The archive contains \textbf{257 training runs with recorded arguments}: 11 source-pretraining runs (stage 1); \textbf{201 runs belonging to the 26 enumerated configurations} (both arms); 14 side-runs inside those configurations that the roster rule excludes (the 150/200-epoch side budgets and the source-transfer controls); and 31 runs from non-family blocks (18 module/architecture-modification arms, 5 un-prefixed MAFA\ensuremath{\rightarrow}mask arms, 4 pipeline-repair runs, 3 three-epoch seed controls and one pretrain run). The registration record's own figures --- 135 registered runs, 84 graded runs (28 cells \ensuremath{\times} 3 seeds), 177 runs in total --- describe that record at its snapshot; they \textbf{under-count the archive} because the SFCHD pipeline and the post-snapshot extensions are not in them, and they are reported as recorded and used as a denominator nowhere. Family-level arithmetic, the BH corrections and the roster all use the archive-based counts; the pre-specified ten-seed extension (42 runs across three arms) is reported alongside the screening readings and is included in the 201.

\textbf{Table T12. The dual-metric accounting behind C4, cell by cell (recomputed from the released runs; same convention as Table T10: best-checkpoint endpoints, the same-budget \texttt{base100} baseline).} The fifteen cells are the 100-epoch strategy arms of the five configurations that carry the three loss-prior families.

{\footnotesize
\setlength{\tabcolsep}{4pt}
\begin{longtable}{p{0.154\textwidth}p{0.071\textwidth}p{0.071\textwidth}p{0.213\textwidth}p{0.249\textwidth}p{0.119\textwidth}}
\toprule
configuration & arm & runs & $\Delta$mAP50 (pp) & $\Delta$mAP50-95 (pp) & sign split \\
\midrule
\endhead
fsin & css & 1 & -0.14 & -0.07 & no \\
fsin & sns & 1 & -0.10 & +0.00 & yes \\
fsin & pws & 1 & -0.24 & -0.05 & no \\
smokefs & css & 1 & -0.17 & +0.08 & yes \\
smokefs & sns & 3 & +0.47 & +0.35 & no \\
smokefs & pws & 1 & +0.78 & +0.44 & no \\
mask & css & 1 & +0.01 & +0.01 & no \\
mask & sns & 1 & +0.00 & +0.05 & no \\
mask & pws & 1 & -0.01 & -0.00 & no \\
mende & css & 4 & -0.28 & +0.27 & yes \\
mende & sns & 4 & -0.33 & -0.20 & no \\
mende & pws & 1 & -0.28 & +0.15 & yes \\
mendein & css & 1 & -0.41 & -0.60 & no \\
mendein & sns & 1 & -0.42 & -0.56 & no \\
mendein & pws & 5 & +1.03 & +0.17 & no \\
\botrule
\end{longtable}}

Marginals over the fifteen cells, under \textbf{one stated sign convention}: the exact zero of $\Delta$mAP50-95 is counted by its rounded sign, so \textbf{\texttt{+0.00} (\texttt{fsin}/\texttt{sns}) counts as non-negative and \texttt{-0.00} (\texttt{mask}/\texttt{pws}) counts as negative} --- the convention the table's \texttt{yes} column follows. Under it \textbf{10 of 15} have $\Delta$mAP50 $<0$ and \textbf{9 of 15} have $\Delta$mAP50-95 $\geq 0$; the number of cells showing both signs is \textbf{4}, equal to the number of \texttt{yes} rows, against an independence expectation of \textbf{6.0}, so the observed count does not exceed chance. Counting \texttt{-0.00} as non-negative instead would give \textbf{10 of 15} and \textbf{five} such cells, with an independent expectation of \textbf{6.7}. The observed four is a \textbf{description of this accounting, not evidence of a mechanism or of a family-level law}; the cell-level mechanism evidence is the per-class decomposition the main text points to, measured on \textbf{two} cells --- the minority-class recall cost on SHWD\ensuremath{\rightarrow}SFCHD and the coverage channel on the dota15 probe. Seed counts are in the \texttt{runs} column: 11 of the 15 cells are single-run, four carry paired seeds, and the table is therefore a sign accounting rather than a test.

\section*{Appendix J. The joint bound of \S{}4, its derivation and its empirical reading}
\addcontentsline{toc}{section}{Appendix J. The joint bound of \S{}4, its derivation and its empirical reading}

This appendix carries the material the main text states only as motivation: the gap corollaries and their elementary character, the sharp optimal-displacement proposition, the correction that \textbf{removed} the interaction term we previously claimed, the assumption conditions the earlier derivation used without naming, and the diagnostic that replaces the curvature term --- including the pre-specified prospective test that partially failed.

\subsection*{J.0 The term-by-term correspondence of the bound of \S{}4}
\addcontentsline{toc}{subsection}{J.0 The term-by-term correspondence of the bound of \S{}4}

\textbf{What this section is, and is not.} It accounts for the terms of the bound of \S{}4 one by one. It is \textbf{not} an explanation of the budget result: the bound's inputs ($L$, $L'$, $W_1$, $\mu_T$, $r_T$) do not vary with the epoch budget, so the bound is constant in it and cannot produce the trend \S{}4.1 reports. The two are reported side by side, not derived from one another.

\textbf{The shift term.} Unbounded here, and the only term with an \emph{independent} measurement attached: the shift distance on which every cell is ordered (\S{}8.2) proxies the transport discrepancy underlying it. Neither bounds the other; both name the same axis, which is why \S{}5.2's ordering is stated and then demoted to a three-point heuristic (\S{}9.4).

\textbf{The source-side residual --- the only term the experiment tests directly.} This is the optimisation error at the reported checkpoint, and the term a strategy acts on: \S{}5.1 resolves it over the budget as $a_\pi\,h(T)$, so the model's prediction about the convergence axis is a prediction about this term, not about the gap. That yields the one quantitative, time-resolved claim this theory contributes: \textbf{under strategy $\pi$ the residual follows $a_\pi h(T)$}, positive and decreasing below the threshold $T_*$ of \S{}5.1 and negative above it. The graded cells are consistent with it in the direction it predicts, the architecture-level repeat locates the crossover for that family at a budget where the baseline has already converged, and the third cell carries it to a target the registered family does not contain (the backbone swap of \S{}4.4). This is a correspondence to a \emph{pattern of readings}, not a fit: the constants of $h$ are not estimated, and no reading here measures that term alone.

\textbf{The source-optimal approximation term.} Nothing here measures it and no claim rests on it; the experiment supplies only \emph{negative} results about interventions that would move it --- the three box-regression loss priors and the structure modules, reported as no detectable steady-state gain (\S{}8.4). Reading those as evidence about this term would require a prior's benefit to land there and nowhere else; \S{}4.3(i) is the general reason we do not assume it.

\textbf{The optimal-displacement cost.} This is the term \S{}4.2's sharp bound prices, in two branches with different assumption footprints. It is \textbf{not measured} here: that would need $L$, $L'$, $W_1$ and the modulus, of which only the modulus is measured, and the conditions for the quadratic branch fail at the computed checkpoints. It is offered as sharpness about an inequality, not as a quantity to be checked.

\textbf{One assumption fails at the measured endpoints, and the bound survives it.} The ordered-pair form of H5 is what the sensitivity bound of \S{}4.2 consumes, and it does not hold where the measurements are taken: $\lambda_{\min}(H_T)<0$ in \textbf{220 of 220} endpoint pairs, so no ball carrying that content exists at any checkpoint we can read (Appendix G). The relevant distinction is between a \textbf{definitional} hypothesis, which states what the setting is (the finite Lipschitz constant of the loss) and is met by construction, and an \textbf{admissibility} condition, which states where a statement is licensed (the ball form and the two conditions behind the quadratic branch) and can simply fail at the points we compute. What follows is about licence, not truth: the quadratic branch is unavailable at these endpoints, the bound is not falsified by that, and the linear branch is unaffected --- it needs the Lipschitz constant and a common parameter space and nothing else. Where the two domains' class counts differ the head is re-initialised and neither branch exists, which is the one case these artefacts cannot repair.

\textbf{The direction of the coupling.} One term of four is tested, as a \emph{response to a strategy} rather than a magnitude; the other three are named, bounded or deliberately left alone, and none is confirmed by any reading here. An empirical regularity can support the claim that the source-side residual responds to the optimisation path; it cannot confirm the other terms, however many cells agree. It also locates where a future experiment would aim --- the displacement cost, the one term with a sharp bound and no measurement.

\subsection*{J.1 The two statements and their scope}
\addcontentsline{toc}{subsection}{J.1 The two statements and their scope}

\begin{quote}
\textbf{Assumptions at a glance (H1--H10)}. H1 loss Lipschitz in data ($L$); H2 gradient Lipschitz in data ($L'$); H3 \ensuremath{\beta}-smoothness (retained for the anisotropic remark; \textbf{not used by the results of \S{}4}); H4 Hessian operator Lipschitz ($M$; retained for the anisotropic remark and the diagnostic, \textbf{not used}); H5 target-side local strong convexity ($\mu_T$, radius $r_T$; read as the \textbf{two-point} inequality, which is what Proposition 1 uses --- Corollary 1 does not use H5 at all --- and the weaker one-point growth reading does not suffice); H6 bounded parameter domain; H7 isotropic curvature approximation; H8 C\textsuperscript{2} regularity (\textbf{not used by the results of \S{}4}); H9 source-side local strong convexity ($\mu_S$); H10 small shift $W_1 \le \mu_T r_T^2/4L$ (\textbf{a compatibility check only} --- it does not by itself place $\theta^*_S$ in the strong-convexity ball). Three conditions the earlier derivation used without naming are stated here and carried as explicit hypotheses: \textbf{localization} $\theta^*_S\in B(\theta^*_T,r_T)$, \textbf{source stationarity or the directional variational inequality} ($\nabla R_S(\theta^*_S)=0$ for an interior minimiser, or $\langle\nabla R_S(\theta^*_S),\Delta\rangle\le0$, which holds when $\Theta$ is convex and is otherwise assumed), and \textbf{non-negative loss} $\ell\ge0$. Neither statement contains \textbf{any curvature term}; the $\lambda_{\max}$ convention belongs to the diagnostic of this appendix, and the measurement protocol's relative radius $\epsilon = \rho\|\theta\|\hat v$ is the convention used in \S{}7. Localization is assumed directly; the earlier claim that $\theta^*_S$ lies in the ball \emph{because} of H10 is withdrawn.
\end{quote}

\textbf{The two statements, restated for reference.} Under H1 and $\ell\ge0$, for \textbf{any} parameter value $\theta\in\Theta$:

\[\mathrm{Gap}_T(\theta) \;\le\; {R}_S(\theta) + L W_1 \qquad (\ell\ge0) \tag{4}\] (Corollary 1), and, under H1, H2, H5 in its two-point form \textbf{with localization}, and source stationarity, \[(4) \;\le\; \min\Bigl\{2LW_1,\; \frac{L'^2}{2\mu_T}W_1^2\Bigr\}. \tag{5}\] (Proposition 1). The usage matrix of Appendix G gives each statement's exact footprint, and the branch-specific footprints differ (the linear branch needs H1 and source-optimality; the quadratic branch needs H2, H5, localization and stationarity, and not H1).

The first inequality is the \textbf{elementary} target-gap corollary --- Lemma 1 plus non-negativity --- and it contains no displacement term. The second is the \textbf{sharp optimal-displacement sensitivity} of Proposition 1, and it does \textbf{not} improve the first: since $R_T(\theta^*_T)\ge0$, the displacement term is redundant for the gap bound, and we say so rather than let a sharp component estimate be read as a sharp overall bound. What the displacement analysis contributes is the magnitude of $(4)$ and the sharpness of its constant --- $1/(2\mu_T)$, attained exactly by the point-mass family of Appendix C step 4. The two branches of Eq. (5) cross at $W_1=4L\mu_T/L'^2$ (quadratic below, linear above), and the $W_1^{3/2}$ estimate is never the lower envelope of the two, so it is not used as a bound anywhere; it is printed only in \textbf{Appendix C}, where the retraction is explained, and is marked withdrawn there. \textbf{No spectral quantity, no transfer constant and no $\mathcal{O}(\cdot)$ remainder appears}, and H3, H4, H8, Lemma 4 and the curvature-transfer lemma are all unused. The source-domain optimisation residual $\varepsilon_{\mathrm{opt}}(\theta) := R_S(\theta) - R_S(\theta^*_S)$ is term (2) of the decomposition (Eq. 3), \emph{contained in} $R_S(\theta)$, whose budget dependence is the convergence axis modelled in \S{}5. Both statements hold for \emph{any} \ensuremath{\theta} --- including $\theta_{\mathrm{src}}$ and the fine-tuned endpoints $\theta_{\pi,T}$ the experiments grade --- and both are in population risks and deterministic; the empirical layer is deferred to the omissions note. \textbf{What kind of statement this is.} Neither is a numerically tight bound. The gap corollary is elementary; the displacement proposition is a sharp statement about a quantity the gap bound does not need. Telescoping Eq. (3) once with the middle term dropped by source-optimality gives the two-KR form $\mathrm{Gap}_T(\theta)\le\varepsilon_{\mathrm{opt}}(\theta)+2LW_1$, also displacement-free; the two forms must not be combined, since doing so charges the source risk twice.

\textbf{Deliberate omissions, and three conditions that used to be omitted assumptions rather than terms.} \emph{(Omitted terms.)} First, a source-complexity term $C_2S_S(\theta)$ is derivable by a PAC-Bayes argument but has no testable content here, so it is a remark (that theory supplement, Remark 6), not a theorem term. Second, the theorem contains \textbf{no curvature term at all}; the one the decomposition invites is removed for the two independent reasons set out below, and where a curvature term is genuinely required (negative directional curvature) the correct modulus is $\kappa^-_\Delta$, kept out of the theorem because it would need the displacement-stability assumption that H5 currently supplies and because it is measured on computed endpoints rather than population ones. The perturbation-measured flatness is likewise a measurement remark --- it matches curvature only at stationary points and is first-order dominated elsewhere. Third, the statements are deterministic and in population risks; the empirical version carries the standard uniform-convergence surcharge, whose constants are never calibrated here (\S{}9.3) and on which no graded claim rests. \emph{(Conditions that were assumptions in disguise.)} \textbf{Localization} ($\theta^*_S$ lies in the strong-convexity ball --- this is the one H10 does \emph{not} supply), \textbf{source stationarity or the directional variational inequality}, and \textbf{non-negativity of the loss}; the assumption-usage matrix of Appendix G is built against this list.

\subsection*{J.2 Empirical reading: bounded, and not a validated predictor}
\addcontentsline{toc}{subsection}{J.2 Empirical reading: bounded, and not a validated predictor}

The strict-protocol regression of the target-gap proxy on flatness gives $r = +0.004$ ($n = 22$ endpoint readings --- source-stationary and fine-tuned endpoints across budgets and strategies, as tabulated in the released flatness files; $p = 0.99$) across budgets --- no linear relationship --- so the flatness proxy is not an empirically calibrated predictor of the target gap within the normal range. This regression was the test we designed for the interaction term; with that term now \textbf{removed} (Appendix C step 5), it stands as a diagnostic result: the endpoint geometry carries no predictive signal on these cells.

\textbf{Erratum for the withdrawn reading.} Earlier versions \emph{regime-ised} this term on a single pathological configuration described as "target-side flatness \ensuremath{\approx} 33\ensuremath{\times} the normal range, collapsing across budgets (\ensuremath{-}2.4 pp at 30 ep, \ensuremath{-}2.0 pp at 100 ep)". Neither element reproduces. (i) \emph{The 33\ensuremath{\times} figure has no traceable measurement}: under the canonical protocol (worst-direction increment at the relative radius $\rho\|\theta\|$, shapeiou reference loss, test split) the ratio between that run and its matched family baseline is \textbf{1.6\ensuremath{\times} / 2.0\ensuremath{\times} / 3.4\ensuremath{\times}} at \ensuremath{\rho} = 0.001 / 0.003 / 0.01, and the legacy readings behind the published number exceed the canonical ones by three to six orders of magnitude at \ensuremath{\rho} \ensuremath{\geq} 0.003, the gap growing with radius --- the signature of a perturbation not normalised to the relative radius. (ii) \emph{The collapse does not reproduce}: over matched same-seed baselines the reweighted-prior variant gains +0.15 pp at 30 epochs (\ensuremath{-}0.06 / +0.35) and +0.01 pp at 100 epochs (\ensuremath{-}0.14 / +0.24 / \ensuremath{-}0.06). (iii) The cell was also mis-attributed: \texttt{mendein-pws} is mendeley\ensuremath{\rightarrow}mendeley, not MAFA\ensuremath{\rightarrow}mendeley. The erratum is recorded in the audit trail (Appendix E) and in the released grading table.

\textbf{A pre-specified prospective test that partially failed.} Before any outcome data for the new pairs existed we specified that the source-endpoint far-field read (worst-direction increment at \ensuremath{\rho} = 0.01, measured \emph{before} fine-tuning) should sort the sign of the lr0.005 gain, and tested it on two further pairs. On those two \textbf{prospective} cells the record is 1/2: it held on visdrone (lowest read of the new pairs, 283 \ensuremath{\rightarrow} measured \ensuremath{-}0.28 pp) and \textbf{failed on mask --- the largest read of the seven (59,248) --- where the measured gain is \ensuremath{-}2.69 pp} (\ensuremath{-}2.64 pp at 30 epochs as well). The three pairs that motivated the rule (smoke\ensuremath{\rightarrow}SFCHD 7,334 \ensuremath{\rightarrow} +1.76 pp at ten seeds, the SHWD\ensuremath{\rightarrow}SFCHD pair being +0.53 pp on the same convention; MAFA\ensuremath{\rightarrow}mendeley 2,534 \ensuremath{\rightarrow} \ensuremath{-}1.40 pp at three seeds (per-seed \ensuremath{-}1.13 / \ensuremath{-}1.82 / \ensuremath{-}1.26) and mendeley\ensuremath{\rightarrow}mendeley 323 \ensuremath{\rightarrow} \ensuremath{-}0.7 pp, a single-seed 100-epoch reading, each against its own baseline --- the latter replacing a \ensuremath{-}2.27 pp figure we cannot trace to any released run) are its derivation set and cannot count as evidence for it; counting them in gives 4/5 across all five cells, and 3/4 restricted to same-pipeline cells (the cross-pipeline reference SHWD\ensuremath{\rightarrow}SFCHD excluded), which is the count carried in the grading table. Neither convention is a prospective pass. The read is therefore an \emph{exploratory observation, not a validated predictor}; note further that the failing domain's loss scale differs by an order of magnitude (L \ensuremath{\approx} 1,795 versus 137--283 elsewhere), so the read's cross-domain comparability is itself unestablished.

We accordingly do \textbf{not} retain the interaction term. It is removed from the theorem as redundant under H5 and mis-modulused where a curvature term is needed, and its directed replacement $\kappa^-_\Delta$ is reported as a diagnostic. What remains open is not whether the curvature object predicts the gap --- on our measurements it does not --- but whether a scale-free geometry measure or a deliberate curvature-induction experiment could give it content. The additive terms ($R_S(\theta)$, which contains $\varepsilon_{\mathrm{opt}}$, and $C_1W_1$) continue to be exercised through the budget--strategy analysis of \S{}5 and \S{}8.3; the withdrawn $C_4W_1^{3/2}$ branch is \textbf{not} exercised.

\subsection*{J.3 Proof-level detail of the derivation the main text states}
\addcontentsline{toc}{subsection}{J.3 Proof-level detail of the derivation the main text states}

\textbf{The quadratic branch needs no Taylor expansion.} Applying the \emph{defining first-order inequality} of $\mu_T$-strong convexity at $\theta^*_S$ with $y=\theta^*_T$, rearranging and completing the square,

\[(4)\;\le\;a-\tfrac{\mu_T}{2}\|\Delta\|^2\;\le\;L'W_1\|\Delta\|-\tfrac{\mu_T}{2}\|\Delta\|^2\;\le\;\frac{L'^2}{2\mu_T}W_1^2,\qquad a:=\langle\nabla R_T(\theta^*_S),\Delta\rangle,\]

where the first step is the strong-convexity inequality, the second is $a\le L'W_1\|\Delta\|$ (the inner-product estimate of Appendix C, valid in both the interior and the constrained case), and the third maximises $L'W_1d-\tfrac{\mu_T}{2}d^2$ over $d\ge0$. \textbf{No bound on $\|\nabla R_T(\theta^*_S)\|$ is used anywhere} --- that quantity is controlled only in the interior case, where it equals $\|\nabla R_T-\nabla R_S\|$, so a norm-based chain would silently exclude constrained minimisers. \textbf{No Taylor identity, no $C^2$ regularity and no $\beta$-smoothness is used} --- differentiability is part of H2's formulation, not a consequence of it --- so H3, H8 and Lemma 4 are unnecessary for this result. The linear branch is the two-KR telescoping of \S{}4.1 with the middle term dropped by source-optimality. \textbf{Sharpness} is exhibited by the point-mass family of Appendix C step 4, which attains Eq. (5)'s quadratic branch exactly; there $w$ is taken small enough that both minimisers lie in $\Theta$, that localization holds, and that the quadratic branch is the active one.

\textbf{What Corollary 1 presupposes about well-definedness.} Only that the objects are well defined: finite first moments for $W_1$, measurability and integrability of $\ell(\theta,\cdot)$, finite risks (Lemma 1), and a target minimiser --- or an infimum if attainment is not assumed, in which case the same proof gives the slightly more general form with $\inf_\vartheta R_T(\vartheta)$. The scope statement is made at the point of claim (\S{}4); the analysis and its motivation are in \S{}5--\S{}8.

\subsection*{J.4 The curvature-term retraction}
\addcontentsline{toc}{subsection}{J.4 The curvature-term retraction}

\textbf{The curvature term that is not available, and what replaces it.} Earlier versions printed a $W_1\times$ top-curvature term. It is not available, for two independent reasons (algebra in Appendix C step 5): with H5 read as genuine two-point strong convexity the quadratic term of the second-order identity is non-negative, so the displacement cost is bounded by its first-order part alone and $C_3W_1\lambda_{\max}(H_T(\theta))$ adds only slack; and if strong convexity is weakened to one-point growth the modulus is wrong --- $|x^\top H_Tx|\le\max(\lambda_{\max},-\lambda_{\min})\|x\|^2$, and the case that needs the term is exactly the one where a \emph{negative} eigenvalue exceeds $\lambda_{\max}$ in magnitude --- with $H_T=\mathrm{diag}(-10,1)$ and $\Delta\propto e_1$, $|\Delta^\top H_T\Delta|=10\|\Delta\|^2$ against $\lambda_{\max}\|\Delta\|^2=\|\Delta\|^2$. What a non-strongly-convex analysis needs is the negative part of the \textbf{path} quantity: the negative directional curvature $\kappa^-_\Delta(H_T)=\max\{0,-\widehat\Delta^\top H_T\widehat\Delta\}$ taken over the displacement segment, because over the 220 endpoint pairs the directional curvature \textbf{changes sign along the displacement in 168 of them (76\%)} and its endpoint value disagrees in sign with the ordered-pair secant modulus in 87, so a value at either end is not the object. It stays out of the bound and is reported as a same-protocol relative diagnostic (\S{}7.4), which is the use that section permits (results in Appendix J.2). For the same reason the generic-displacement estimate $2L'\sqrt{L/\mu_T}\,W_1^{3/2}$ is \textbf{withdrawn}: it is never the lower envelope of Eq. (5)'s branches, so it cannot be the active one and cannot be asymptotically sharp under strong convexity. Corollary 1 needs H1 and $\ell\ge0$ only; Proposition 1's \textbf{linear} branch needs H1 and source-optimality, while its \textbf{quadratic} branch needs H2, the two-point strong-convexity inequality for the ordered pair, localization and source stationarity and does \textbf{not} use H1.

Term (4) is the optimal-displacement cost $R_T(\theta^*_S)-R_T(\theta^*_T)$; write $\Delta := \theta^*_S-\theta^*_T$. 1. \textbf{The two conditions the argument needs, named.} (a) \emph{Localization}: $\theta^*_S$ lies in the ball $B(\theta^*_T,r_T)$ on which H5 holds --- an assumption. H10 does \textbf{not} supply it: H10 bounds $W_1$ relative to the growth scale, but a source minimiser sitting in a distant well that is only $O(LW_1)$ worse than the target well is not excluded by it, and deriving the displacement bound from local growth at $\theta^*_S$ would presuppose the very placement being established. A readable sufficient condition is the basin-separation certificate $\inf_{\theta\in\Theta\setminus B(\theta^*_T,r_T)}[R_T(\theta)-R_T(\theta^*_T)]>2LW_1$, together with $\sup_\Theta|R_S-R_T|\le LW_1$: then $R_S(\theta)>R_S(\theta^*_T)$ for every $\theta$ outside the ball, so no global source minimiser lies there. (b) \emph{Source stationarity, or the directional variational inequality}: $\nabla R_S(\theta^*_S)=0$ for an interior minimiser; the weaker condition $\langle\nabla R_S(\theta^*_S),\Delta\rangle\le0$ suffices and is what the argument uses. That inequality follows from first-order constrained optimality \textbf{when $\Theta$ is convex} and $\theta^*_T\in\Theta$; if convexity is not assumed it is an explicit assumption. Boundedness of $\Theta$ is neither sufficient nor the relevant condition, and an appeal to it ("$\Theta$ bounded and $\theta^*_T$ reachable") is withdrawn. 2. \textbf{First-order term (Lemma 2, source stationarity) --- inner-product form, which covers both cases.} Write $a=\langle\nabla R_T(\theta^*_S),\Delta\rangle$ and split it as $\langle\nabla R_T(\theta^*_S)-\nabla R_S(\theta^*_S),\Delta\rangle+\langle\nabla R_S(\theta^*_S),\Delta\rangle$. The first term is $\le L'W_1\|\Delta\|$ by Lemma 2 and Cauchy--Schwarz; the second is $0$ when $\theta^*_S$ is an interior minimiser and $\le0$ when it is a constrained minimiser satisfying the directional inequality --- the case split of Appendix C step 2. Hence $a\le L'W_1\|\Delta\|$ \textbf{in both cases}.

\emph{(Scope note: Lemma statements and the complete term-(4) algebra are reproduced here so that the bound is checkable as submitted; the per-lemma proofs, the PAC-Bayes derivation behind Remark 6 and the anisotropic details of Appendix B are carried in that theory supplement, which is \textbf{not part of this release}.)}

\subsection*{J.5 The material of \S{}4 carried in this appendix}
\addcontentsline{toc}{subsection}{J.5 The material of \S{}4 carried in this appendix}

\textbf{(a) The displacement conditions, named in full.} $(4)=R_T(\theta^*_S)-R_T(\theta^*_T)$ is non-negative by target optimality. Two conditions must be named before it can be bounded, and both are assumptions: \textbf{localization}, $\theta^*_S$ lying in the ball $B(\theta^*_T,r_T)$ on which H5's strong convexity holds --- which H10 does \emph{not} supply; and \textbf{source stationarity}, $\nabla R_S(\theta^*_S)=0$ for an interior minimiser, or the directional variational inequality $\langle\nabla R_S(\theta^*_S),\Delta\rangle\le0$, which follows from constrained optimality when $\Theta$ is convex and is an assumption otherwise. Both are exact statements about minimisers, and a reproduced endpoint satisfies neither exactly.

\textbf{(b) What released checkpoints can support (per-cell detail).} With $\varepsilon_{\mathrm{st}}:=\|\nabla R_S(\theta^*_S)\|$ measured at the source endpoint the quadratic branch becomes $\frac{(L'W_1+\varepsilon_{\mathrm{st}})^2}{2\mu_T}$, and the linear branch becomes $(4)\le 2LW_1+[R_S(\theta^*_S)-R_S(\theta^*_T)]$, in which source-optimality is \textbf{replaced by a measured term} rather than assumed: the bracket is the very quantity the telescoping produces, so evaluating it \emph{is} discharging the assumption ($\varepsilon_{\mathrm{st}}=0$ is \textbf{sufficient but not necessary}). Over the 220 endpoint pairs of the read suite (Appendix G) source optimality holds in 59 of the 164 head-compatible pairs, the directional variational inequality in 32, and \textbf{both together in only 14} --- so at computed checkpoints the exact conditions are more often false than true. The \textbf{linear} branch needs H1 and a common parameter space and nothing else; the \textbf{quadratic} branch additionally needs the two-point inequality for the ordered pair, localization and source stationarity --- and localization is not merely unmeasured there, it is unavailable: the ball form of H5 is refuted in 220 of 220 pairs. The quadratic branch is therefore \textbf{conditional} on conditions the released checkpoints do not meet in general, and the branch a released-checkpoint measurement can bear on is the linear one; where class counts differ the head is re-initialised and neither form exists (56 of the 220 pairs).

\textbf{(c) The 4.1 term-by-term reading and the elementary character of Corollary 1.} The analysis uses ten numbered assumptions (H1--H10; box in Appendix J.1, usage matrix in Appendix G). It is elementary because that is what it is ($\ell\ge0$ gives $R_T(\theta^*_T)\ge0$, then Lemma 1), and we lead with it because it is \textbf{never weaker} than any form that adds a non-negative displacement term to it. Terms: (1) is Kantorovich--Rubinstein-bounded at $LW_1$; (2) is the optimisation residual $\varepsilon_{\mathrm{opt}}$, \emph{contained in} $R_S(\theta)$ rather than added to it; (3) enters only the two-KR variant; (4) is the displacement cost of \S{}4.2. Without non-negativity the two-KR comparison gives $\mathrm{Gap}_T(\theta)\le\varepsilon_{\mathrm{opt}}+2LW_1$, also displacement-free; the two forms must not be combined, since doing so charges the source risk twice.

\textbf{(d) Algebraic restatement.} An exact decomposition (Eq. 3) whose elementary corollary is $\mathrm{Gap}_T(\theta)\le R_S(\theta)+LW_1$ --- the classical additive shape --- together with a \textbf{sharp} optimal-displacement sensitivity bound, $(4)\le\min\{2LW_1,\,L'^2W_1^2/(2\mu_T)\}$, whose quadratic branch is attained exactly.

\textbf{(e) The status of the bound.} The displacement term does \textbf{not} improve the gap bound (non-negativity makes it redundant there), and the $W_1^{3/2}$ estimate is dominated at every shift, so neither is used as a bound anywhere. The interaction term $C_3W_1\lambda_{\max}(H_T(\theta))$ is \textbf{removed} for the two independent reasons given in J.4. The bound is a \emph{decomposition and a hypothesis generator}, not a validated theory; its curvature object is one that our own measurement-hierarchy result (\S{}7) shows is \emph{not measurable} under this protocol \textbf{as an off-trajectory outcome predictor} --- it is measured, and reported, as a same-protocol relative diagnostic, which is the use \S{}7.4 permits, whereas what \S{}7 finds measurable as a proxy is a trajectory-evaluated first-order quantity --- so C1 should be read as a scaffold that makes the theory--experiment gap explicit.

\section*{Appendix K. Campaign anchors and controls}
\addcontentsline{toc}{section}{Appendix K. Campaign anchors and controls}

\subsection*{K.1 The campaign: arms, seeds and baseline anchors}
\addcontentsline{toc}{subsection}{K.1 The campaign: arms, seeds and baseline anchors}

\textbf{Two cells of Fig. S3 are seeded pairs.} On smoke\ensuremath{\rightarrow}fire-smoke at 100 epochs the lr0.005 arm reads +0.10 and +0.30 pp against its same-seed \texttt{base100} baselines (seeds 43 and 44; the earlier pair carried no seed suffix and formed no seed-wise pair); on mendeley-in the \texttt{pws} arm reads \ensuremath{-}0.20 and +0.20 pp (same seeds), mean 0.00 pp, consistent with the loss-prior null of \S{}6.3 rather than evidence of an effect. Both cells are now three-seed paired in Fig. S3.

\textbf{Table T11.} Supporting quantities for the three-seed readings, recomputed from the released \texttt{results.csv} best-epoch column (the same convention as \S{}8.2's pipeline-reproduction note) rather than quoted from an earlier record; each row names the assertion, the quantity that carries it, and the value that quantity now takes.

{\footnotesize
\setlength{\tabcolsep}{4pt}
\begin{longtable}{p{0.463\textwidth}p{0.243\textwidth}p{0.156\textwidth}p{0.050\textwidth}}
\toprule
assertion & supporting quantity & value & asserted in \\
\midrule
\endhead
strong cell (smoke\ensuremath{\rightarrow}SFCHD, 100 ep): \emph{complete separation} --- every lr0.005 run above every baseline run & the separation gap, min(lr runs) \ensuremath{-} max(baseline runs) & \textbf{+1.28 pp} (43.67 against 42.39) & \S{}8.3 \\
corroborating cell (SHWD\ensuremath{\rightarrow}SFCHD, 100 ep): three-seed effect size & strategy-side / paired-difference & \textbf{3.8\ensuremath{\sigma}} / \textbf{1.9\ensuremath{\sigma}} & \S{}8.3, \S{}9.4 \\
dota15 within-domain, 30 epochs (the non-positive counterpart cited at the same-budget control) & the two endpoints & \textbf{0.1534 \ensuremath{\rightarrow} 0.1526} (\ensuremath{-}0.08 pp) & \S{}8.3 \\
\botrule
\end{longtable}}

The strong cell's paired three-seed gain recomputes to \textbf{+1.783 pp} (paired t = \textbf{9.44}), the +1.78 pp screening reading quoted in \S{}8.3. \textbf{Three numbers for this one cell appear in this file and they are different objects, not a drifting value}: \textbf{+1.76 pp} is the ten-seed best-checkpoint level on the split that selected the checkpoint, \textbf{+1.783 pp} is the same cell's paired three-seed recomputation, and \textbf{+1.78 pp} is the screening reading of \S{}8.3; the corroborating cell recomputes to \textbf{+0.51 pp} (paired t = \textbf{3.31}, p \ensuremath{\approx} 0.081), which is what \S{}8.2 prints.

\textbf{Dose--response, label budget and controls (52 runs, tabulated in the released record; the table numbers in circulation are T7, T8 and T10--T12).} \emph{Dose--response (SHWD\ensuremath{\rightarrow}SFCHD, same pipeline and seed, gains in pp over the same-budget baseline).} At 30 epochs: lr 2.5\ensuremath{\times} +0.86, 5\ensuremath{\times} +0.87, 10\ensuremath{\times} 0.00, 20\ensuremath{\times} \ensuremath{-}0.76; at 100 epochs: 2.5\ensuremath{\times} +0.07, 5\ensuremath{\times} +0.29, 10\ensuremath{\times} +0.53, 20\ensuremath{\times} \ensuremath{-}0.38. The response has an interior optimum at \textbf{both} budgets and the optimum moves with the budget --- 2.5--5\ensuremath{\times} at 30 epochs, 10\ensuremath{\times} at 100 epochs; every tier is single-seed, so this is a dose \emph{ordering} and not a graded reading --- so the registered 5\ensuremath{\times} tier is a matched choice rather than a special value, and the largest tested tier is negative at both budgets, bounding the escape window from above. \emph{Label budget (100 epochs, fixed test set, train subsets of the same pool at 10/20/30/50\% of labels = 1,206/2,413/3,619/6,033 images).} The baseline rises monotonically with labels (0.4128 \ensuremath{\rightarrow} 0.4397 \ensuremath{\rightarrow} 0.4502 \ensuremath{\rightarrow} 0.4632, i.e. +5.04 pp in total, marginally diminishing), while the lr0.005 gain stays positive at every fraction (+0.51 / +0.29 / +0.70 / +1.25 pp). The 20\% fraction, re-run under the current pipeline, reproduces the registered baseline and strategy values to the digit (0.4397 / 0.4426), confirming that no pipeline offset separates the two. With one seed per fraction the \emph{ordering} among these four gains is not interpretable. \emph{Standard-baseline controls.} Neither a cosine schedule at the same initial learning rate (30 ep 0.4263, +0.12 pp; 100 ep 0.4413, +0.16 pp) nor a frozen backbone (first nine modules; 30 ep 0.3757, \ensuremath{-}4.94 pp; 100 ep 0.4002, \ensuremath{-}3.95 pp) reproduces the tier effect, so it is not a schedule artefact and not available from freezing. \emph{Layer-wise decay, measured rather than substituted.} The discriminative/layer-wise baseline review requested --- the empirical form of the discriminative-fine-tuning and gradual-unfreezing recipe of ULMFiT (Howard and Ruder, ACL 2018) and of LP-FT \citep{zhang2026} --- is run directly: peak lr 0.005 with a per-layer decay factor 0.1 (stem lr 5\ensuremath{\times}10\textsuperscript{-4} \ensuremath{\rightarrow} head lr 5\ensuremath{\times}10\textsuperscript{-3} over 22 layers, per-layer parameter groups; 3 seeds, same pipeline). It does \textbf{not} reproduce the gain and lands \emph{below} the same-budget baseline: at 100 epochs 43.50 (layer-wise) vs 44.42 (uniform 5\ensuremath{\times}) vs 43.91 (baseline), i.e. \textbf{\ensuremath{-}0.92 pp vs the uniform arm} (paired t = \ensuremath{-}7.24, p = 0.019) and \textbf{\ensuremath{-}0.42 pp vs the baseline} (t = \ensuremath{-}5.12, p = 0.036); at 30 epochs \textbf{\ensuremath{-}2.24 pp vs uniform} (t = \ensuremath{-}40.92, p = 0.001) and \textbf{\ensuremath{-}1.23 pp vs baseline} (t = \ensuremath{-}9.86, p = 0.010), all three seeds in the same direction. The gain therefore requires the large learning rate to be applied \emph{globally}, not merely at the head. Three decay factors (0.03, 0.1 and 0.3) were tested at two budgets at the same peak lr. \emph{Optimiser boundaries.} AdamW at lr 0.001 / 0.0025 (100 ep) gives +0.33 / +0.31 pp with no dose--response inside that range --- adaptive scaling substitutes for a higher peak learning rate at low tiers but does not reproduce the SGD dose--budget mechanism, consistent with the SGD-bounded scope of the norm-diffusion relation, which \S{}8.5 reports as descriptive rather than calibrated.

\textbf{Seed count of the whole campaign, stated because it bounds what these controls can show:} every campaign arm --- dose tiers, label budgets, cosine, frozen backbone and AdamW --- is \textbf{single-seed (shuffle-seed 42), paired against the same-seed baseline}, and its baselines are re-run under the current pipeline rather than reused from the registry; the one exception is the dota15 within-domain control carried in \S{}8.3, which has three runs on data-order seeds \textbf{s42, s43 and s44} (the third run's directory name omits the \texttt{\_s42n} suffix the other two carry, and its seed is that campaign's fixed \texttt{--shuffle-seed 42}, recoverable from its recorded configuration). Its three per-seed readings are +1.12 / +0.95 / +0.93 pp, and the suffix-less pair contributes the +1.12 pp best-epoch reading. The campaign therefore supports a statement about direction and rough magnitude at that one seed, not about seed variability. \textbf{Baseline anchors for those arms} (test-split mAP50-95, same-seed baseline in brackets): dose 2.5\ensuremath{\times}/5\ensuremath{\times}/10\ensuremath{\times}/20\ensuremath{\times} at 30 epochs +0.86/+0.87/0.00/\ensuremath{-}0.76 pp (baseline 0.4219) and at 100 epochs +0.07/+0.29/+0.53/\ensuremath{-}0.38 pp (baseline 0.4393); label budgets 10/20/30/50\% +0.51/+0.29/+0.70/+1.25 pp (baselines 0.4128/0.4397/0.4502/0.4632); cosine +0.12/+0.16 pp (0.4263/0.4413); frozen backbone \ensuremath{-}4.94/\ensuremath{-}3.95 pp; AdamW at 0.001/0.0025 +0.33/+0.31 pp.

\subsection*{K.2 Structure modules and loss priors: the block readings and the screening line}
\addcontentsline{toc}{subsection}{K.2 Structure modules and loss priors: the block readings and the screening line}

\textbf{Structure modules.} Sixteen repaired-protocol runs (eleven modules: SE/CBAM/C2fEMA/SimAM/CA/EMA/ECAplus/SA/GAM/GsConv/ADown; insertion and replacement; two protocols; checkpoint-structure-verified with parameter deltas matching the designs --- the parameter-free SimAM adds zero) contain no gain \ensuremath{\geq} +0.4 pp (range \ensuremath{-}0.41 to +0.33 pp; the +0.4 pp line is three times a typical strategy-side \ensuremath{\sigma} of \ensuremath{\approx}0.13 pp --- the same screening logic as the frozen criterion, applied here to single-run module cells, whose noise is larger than the criterion is calibrated for); SE is consistently mildly positive (+0.12/+0.28/+0.33 pp) and below the line, and the six attention modules added in the repaired extension span \ensuremath{-}0.10 to +0.22 pp on the B-pipeline protocol. Historical reports of severe insertional degradation were invalidated by checkpoint audit and excluded. \emph{Scope against supervision-first methods.} We do not test pseudo-labelling or labelled-unlabelled mixing routes (e.g. the sparse-label supervised loss of \citep{tian2025} or self-training feedback plugins); those change \emph{what supervision enters training}, whereas the cells above hold the label budget fixed and vary only the loss prior, the module or the optimisation path, so the two families are complementary and \citep{tian2025}'s ordering is consistent with ours.

\textbf{Loss priors (the per-cell readings the null of \S{}8.4 is about).} sns at 100 epochs: \ensuremath{-}0.04 pp (SHWD\ensuremath{\rightarrow}SFCHD, 3v3), +0.06 pp (smoke\ensuremath{\rightarrow}firesmoke, 3 seeds), \ensuremath{-}0.39 pp (MAFA\ensuremath{\rightarrow}mendeley, 3v3, all seeds negative, 1.5\ensuremath{\sigma} below criterion, paired t = \ensuremath{-}3.97, two-sided p \ensuremath{\approx} 0.058, disclosed per protocol; a historical +0.34 pp single reading reversed sign under the corrected pipeline and is disclosed). None reaches the frozen threshold; the three \emph{tested box-regression priors} do not exhaust the loss-design space. css fails directionally in both testable pairs (mechanism unattributed after the alignment closure); pws gains are recall-driven --- the reweighting channel --- and its single within-domain positive cell was removed as a pipeline-pairing artefact.

\subsection*{K.3 Shift-ordering caveats and per-cell detail}
\addcontentsline{toc}{subsection}{K.3 Shift-ordering caveats and per-cell detail}

Two caveats on the ordering: the three pairs are \emph{not} label-space-equivalent (smoke\ensuremath{\rightarrow}firesmoke and smoke\ensuremath{\rightarrow}SFCHD change the class vocabulary, SHWD\ensuremath{\rightarrow}SFCHD does not), so D and the task covary across the ordering; and the reported D values come from the released s-OTDD script at sample counts matched within each pair. \textbf{That composition confound has a measured shape}: a box-geometry $W_1$ decomposition over the same nine pairs --- log-width, log-height, centre-x and centre-y --- shows the shift is predominantly \emph{scale} rather than position, the scale share being 0.757-0.952 (0.860-0.886 across the three cross-dataset pairs), so these corpora differ mainly in object size. This is a decomposition of box statistics, \textbf{not} of the s-OTDD distance.

\subsection*{K.4 The norm-diffusion detail of \S{}8.5}
\addcontentsline{toc}{subsection}{K.4 The norm-diffusion detail of \S{}8.5}

\textbf{Norm diffusion.} Two normalisations must be kept apart: relative to the \emph{pretrained} origin on SHWD\ensuremath{\rightarrow}SFCHD (\ensuremath{\|}\ensuremath{\theta}\ensuremath{\|}\textsuperscript{2} \ensuremath{\approx} 21,374), the baseline inflates by \textbf{+0.62\%} at 100 epochs while the lr0.005 arm inflates by \textbf{+8.62\%}, 3.2-14\ensuremath{\times} across budgets, and the sns and pws arms sit within 0.1\% of the baseline. Relative to \emph{each cell's own same-pipeline baseline}, the lr0.005 inflation follows \ensuremath{\Delta}\ensuremath{\|}\ensuremath{\theta}\ensuremath{\|}\textsuperscript{2} \ensuremath{\approx} k\ensuremath{\cdot}N\_eff on six cells spanning N\_eff = 36 to 10,530 steps (k \ensuremath{\approx} 1.39\ensuremath{\times}10\textsuperscript{-3} \% per step). The law is descriptive, not calibrated: it is normalisation-dependent, it does not survive AdamW (the AdamW endpoint norm is released in \texttt{norm2\_adamw.json}), and it is reported as a mechanism hypothesis scoped to SGD-style implicit stabilisation, not as a law (\S{}9.4). The endpoint norms behind the +0.62\% and +8.62\% figures are released in \texttt{norm2\_extra.json}, so both recompute directly from the released values.

\subsection*{K.5 Positive controls: what bounds the nulls of \S{}8.4}
\addcontentsline{toc}{subsection}{K.5 Positive controls: what bounds the nulls of \S{}8.4}

The nulls of \S{}8.4 are statements about detection. Four controls on the same pipeline, target family and protocol fix the range in which such a null is informative rather than merely unmeasured.

\begin{enumerate}
\item \textbf{Deliberate degradation is detected from below.} The layer-wise-decay schedule at the same peak learning rate lands below both the same-budget baseline and the uniform arm at every factor tested: 0.4333/0.4083 (factor 0.03) and 0.4313/0.4003 (factor 0.3) at 100/30 epochs, against baseline 0.4396 and uniform 0.4426.
\item \textbf{Configuration-scale changes read far outside \ensuremath{\sigma}.} Frozen backbone \ensuremath{-}4.94/\ensuremath{-}3.95 pp; label budgets 10/20/30/50\% +0.51/+0.29/+0.70/+1.25 pp; cosine +0.12/+0.16 pp; AdamW +0.33/+0.31 pp (campaign arms, all single-seed).
\item \textbf{A same-pipeline second architecture resolves both sides of the screening line.} YOLO11n +1.27 pp (t = 6.97, p = 0.020) at 30 epochs and \ensuremath{-}0.01 pp (t = \ensuremath{-}0.10, p = 0.93) at 100 epochs on the same pair (\S{}4.4); \textbf{its DOTA15 runs are not read against the third cell} (different three-way split and initialisation).
\item \textbf{The headline effect is reproduced within a domain.} dota15\ensuremath{\rightarrow}dota15 (within-domain) +1.12/+0.95/+0.93 pp over matched same-seed baselines at seeds s42/s43/s44 (per-run endpoints 0.1614/0.1628/0.1616 \ensuremath{\rightarrow} 0.1726/0.1723/0.1709 pp; 3/3 positive, mean +1.00 pp).
\end{enumerate}

\textbf{What these controls do not do.} They bound what the instrument resolves at the \emph{large} end; they do not lower the detection floor of the nulls themselves, which stays at 0.56--0.62 pp for a three-seed paired test at the observed \ensuremath{\sigma}, and at +0.4 pp by construction for the single-run module screens. Every control above except the within-domain arm is single-seed, so none of them stands in for the multi-seed evidence a positive claim in \S{}8.4 would require --- which is why the module block is reported as a bounded negative rather than as evidence of absence.

\subsection*{K.6 The budget curve in full}
\addcontentsline{toc}{subsection}{K.6 The budget curve in full}

\textbf{Budget curve.} For SHWD\ensuremath{\rightarrow}SFCHD the gain decreases in its point estimates across budgets and then plateaus --- +1.00/+0.65/+0.51 pp at 30/50/100 epochs (every point a full 3v3) and +0.52 pp at 200 epochs (single run) --- so point-estimate monotonicity is claimed over 30\ensuremath{\rightarrow}100 epochs and a plateau, not a continued decline, beyond it; per-seed monotonicity is not claimed. The smoke\ensuremath{\rightarrow}SFCHD pair is direction-consistent at 30/50 epochs but ungraded (single runs).

\section*{Appendix L. Corpus roles and per-corpus split statistics}
\addcontentsline{toc}{section}{Appendix L. Corpus roles and per-corpus split statistics}

\textbf{The per-corpus attribute ranges behind \S{}5.2 (released \texttt{08\_<CJK>\_<CJK>.csv}, 26 rows of domain \ensuremath{\times} split).} Across the ten corpora: median box area \textbf{1.88\ensuremath{\times}10\textsuperscript{-4} \ldots{} 3.61\ensuremath{\times}10\textsuperscript{-1}} of image area (a \textbf{1,919\ensuremath{\times}} span), small-object share \textbf{0.00 \ldots{} 0.99}, empty-image ratio \textbf{0.00 \ldots{} 0.61}, and boxes per image \textbf{0.49 \ldots{} 139.9}; on the 20\% subsets the median box area spans 7.4\ensuremath{\times}10\textsuperscript{-5} \ldots{} 7.8\ensuremath{\times}10\textsuperscript{-2}. These are the release's own per-corpus attributes, printed here because \S{}5.2's scale-dominance reading rests on the corpora differing mainly in object size; the extended variables are not a substitute for the pairwise distances of \S{}5.2.

\textbf{Corpus-by-corpus roles and provenance (Table S1 carries the pairs, and }Table L1 below\textbf{ carries the per-corpus splits).} Each corpus enters through one of three roles --- source pretraining, same-object fine-tuning (7 pairs), or cross-domain fine-tuning --- and the shift between roles is quantified in \S{}8.2 (s-OTDD). The five cross-domain pairs are MAFA\ensuremath{\rightarrow}mendeley, MAFA\ensuremath{\rightarrow}mask and smoke\ensuremath{\rightarrow}firesmoke (B pipeline) together with \textbf{SHWD\ensuremath{\rightarrow}SFCHD and smoke\ensuremath{\rightarrow}SFCHD (SFCHD pipeline)}, the latter two carrying the paper's only criterion-passing cells. Mendeley is also the source endpoint of the mendeley target pair. The masked-face corpora (mafa, mask\_clean, mendeley) are the low-shift end and the aerial corpora (DOTA, VisDrone, AI-TOD) the high-variance end, by box statistics and images per box. Two provenance statements in place of citations we cannot establish: the five corpora marked "as recorded in the released registry" were obtained from public bundles (Kaggle, Roboflow/Mendeley packaging) whose exact provenance is recorded in the registry rather than in a citable publication; and MAFA is used here as a masked-face \emph{source} corpus, as packaged with the mask detectors we compare against, not as the recognition benchmark of \citep{kwon2021} --- we flag the difference rather than imply identity. One distinction is a confound if left implicit: the \textbf{smoke corpus is single-class} (\texttt{nc = 1}, "smoke") while the \textbf{firesmoke corpus is two-class} (fire, smoke); the smoke\ensuremath{\rightarrow}SFCHD line therefore starts from a \emph{single-class} source head adapted to the target's two classes (Appendix D), and the released weights record both stages (\texttt{smoke\_pretrain.pt}, \texttt{nc = 1}; \texttt{smoke\_pretrain\_nc2.pt}, \texttt{nc = 2}). We report both endpoint conventions for that cell where relevant and do not describe the adaptation as free of effect.

\textbf{Scope of this table.} Table L1 covers the ten corpora that carry the pretraining, same-object and cross-domain roles of \S{}8.1; the four further target corpora used by the label-axis and clean-protocol subsections (\textbf{\texttt{neu\_det}, \texttt{gdut\_hwd}, \texttt{chv}, \texttt{fsin}}) are described where they are first used and are deliberately not given roles here. \textbf{Table L1 --- the ten corpora, by role and split} (the main text points here).

{\footnotesize
\setlength{\tabcolsep}{4pt}
\begin{longtable}{p{0.109\textwidth}p{0.239\textwidth}p{0.063\textwidth}p{0.105\textwidth}p{0.059\textwidth}p{0.302\textwidth}}
\toprule
corpus & content & classes & train (20\% subset) & test/val & source \\
\midrule
\endhead
SHWD \citep{jocher2026} & safety-helmet wearing, industrial and construction scenes & 2 (hat, person) & full source split & own test split & public dataset (GitHub) \\
SFCHD \citep{ning2023} & safety clothing and helmet detection & 2 (hat, person) & 2,413 (20\% of source) & 6,033 & public dataset (arXiv:2306.02098) \\
smoke (keremberke) & smoke-only imagery & 1 (smoke) & 15,096 & 2,160 & Kaggle corpus (Roboflow COCO\ensuremath{\rightarrow}YOLO), as recorded in the released registry \\
firesmoke & fire and smoke, cleaner packaging & 2 (fire, smoke) & 56,161 (11,232) & 6,411 & as recorded in the released registry \\
mafa & masked faces & 1 (face mask) & source-pretraining corpus & own split & as recorded in the released registry \\
mask\_clean & face masks & 1 (face mask) & 6,120 (1,224) & 919 & as recorded in the released registry \\
mendeley & face masks & 1 (face mask) & 5,750 (1,150) & 800 & as recorded in the released registry \\
VisDrone \citep{andriushchenko2023} & UAV aerial scenes, pedestrian/vehicle & 10 & 6,471 (1,294) & 1,610 & public dataset (TPAMI) \\
DOTA v1.0 \citep{dziugaite2020} / DOTA-v1.5 & aerial imagery & 15 / 16 & 1,411 (282) & 458 & public dataset (CVPR 2018) \\
AI-TOD \citep{dinh2017} & aerial tiny objects & 8 & 11,214 (2,243) & 2,804 & public dataset (ICPR 2021) \\
\botrule
\end{longtable}}

\section*{Appendix M. Method, protocol and framework detail of the main text}
\addcontentsline{toc}{section}{Appendix M. Method, protocol and framework detail of the main text}

Nothing here is new: every subsection below carries the full form of the technical material the article states in compressed form, so that the article file fits the journal's page range. Each main-text pointer names its subsection, and no claim, number or disclosure differs from the article's.

\subsection*{M.1 The setup in full: task, metrics, shift quantification, flatness and notation}
\addcontentsline{toc}{subsection}{M.1 The setup in full: task, metrics, shift quantification, flatness and notation}

\textbf{The s-OTDD floor, on both released computations.} \emph{(i) The feature-level table behind \S{}5.2's D values} (ResNet-18 features, 4,000 samples, sliced Wasserstein, s-OTDD = \ensuremath{\surd}(W\textsubscript{2}\textsuperscript{2} + label term)) gives nine pairs: within-domain 2.8905 (DOTA15), 3.3524 (AI-TOD), 3.5071 (mendeley), 5.7532 (firesmoke), 6.0521 (mask), 6.4273 (VisDrone); cross-domain 7.1894 (MAFA\ensuremath{\rightarrow}mask), 9.4395 (smoke\_ker\ensuremath{\rightarrow}firesmoke), 13.0313 (MAFA\ensuremath{\rightarrow}mendeley). The lowest cross-domain reading therefore exceeds the highest within-domain reading by \textbf{12\%} (7.1894 / 6.4273 \ensuremath{-} 1 = 0.119). \emph{(ii) The earlier release computed with a different estimator} (\texttt{otdd\_results.csv}) gives within-domain 2.6593 (Paddle subsample), 3.4159 (SFCHD), 7.4686 (SHWD subsample) and cross-domain 9.865 (SHWD vs Paddle), 11.51 (fusion vs SFCHD), 12.7629 (SHWD vs SFCHD), i.e. a floor margin of \textbf{32\%} on that computation as well. Both are stated because the two estimates are not interchangeable: the floor supports that D separates shift from a measured no-shift reference at these two instrumentations, not that the two scales coincide.

\subsubsection*{3.1 Task and evaluation metrics}

Object detection requires, for every object of interest, a localisation (bounding box) and a classification decision. Accuracy is summarised by mAP. We report \textbf{mAP50} (detection; accepts IoU \ensuremath{\geq} 0.5) and \textbf{mAP50-95} (localisation tightness; averaged over IoU 0.50--0.95 with 101-point interpolation and \texttt{max\_det = 300} detections per image, the Ultralytics default). The cap is applied per image across all classes rather than per class, so a dense image can be truncated before every object is scored --- a limitation that bears on the dense-scene corpora (VisDrone, AI-TOD) in particular. The per-cell mAP50 column is carried in Table T12, because the two metrics can move in opposite directions at long budgets (\S{}5.2, \S{}8.5). All experiments use YOLOv12n \citep{kornblith2019} unless another architecture is named explicitly. Two are: the \textbf{YOLO11n} control of \S{}8.3 (the backbone swap), and a \textbf{YOLO26n} arm used for the registration-line checks, so three generations appear in the archive and the per-run table carries the architecture of each run.

\subsubsection*{3.2 Quantifying domain shift}

With $\mu_S,\mu_T$ the source/target data distributions, our primary shift measure is the 1-Wasserstein distance

\[W_1(\mu_S,\mu_T) = \inf_{\pi\in\Pi}\int \|z-z'\|\,d\pi(z,z'), \tag{1}\]

over couplings with the given marginals. Empirically we use a sliced, label-aware estimator over frozen backbone features --- sliced 1-Wasserstein geometry (Bonneel et al. \citep{pototzky2022}) combined with the label-aware dataset distance of Alvarez-Melis \& Fusi's OTDD \citep{tran2019}, henceforth s-OTDD --- decomposed into feature and label terms. The bound of \S{}4 uses this population distance; the measured D is its practical proxy on frozen features, and the constants are \emph{not} calibrated, so the experiments bear on the \emph{form} of the geometric interaction (through the flatness proxies of \S{}4.2 and \S{}7), not on numerical tightness.

\subsubsection*{3.3 Flatness (and its protocol)}

For parameters $\theta$, the flatness on domain $D$ is the worst-direction loss increment

\[S_D(\theta) = \max_{\|\epsilon\|\le \rho\|\theta\|}\;[L_D(\theta+\epsilon)-L_D(\theta)], \tag{2}\]

where the radius is \textbf{relative to the parameter norm} ($\epsilon = \rho\|\theta\|\hat v$; $\rho \in \{0.001, 0.003, 0.01\}$, primary 0.001; $\|\theta\| \approx 148$--158 in our runs, $\|\theta\|^2$ recorded per model). Near a stationary point $S_D \approx \tfrac{\rho^2\|\theta\|^2}{2}\lambda_{\max}(H_D)$; at non-stationary points the increment is first-order dominated, $S_D \approx \rho\|\theta\|\,\|\nabla R_D(\theta)\|$ --- a slope/distance proxy, verified by the \ensuremath{\varepsilon}-scaling decomposition of \S{}7.2(1) and central to the measurement-hierarchy proposition.

\subsubsection*{3.4 Notation}

$\mu_S,\mu_T$; $\theta$ (any parameter value); $\theta_{\mathrm{src}}$ (source-training endpoint, the fine-tuning initialisation); $\theta_{\pi,T}$ (endpoint after strategy $\pi$ for $T$ epochs); $\theta^*_S,\theta^*_T$ (population optima); $R_S,R_T$; $W_1$; $S_D$; $\eta,T$; $\mathrm{Gap}_T(\theta)=R_T(\theta)-R_T(\theta^*_T)$; $\Delta_\pi(T)$ (strategy gain); $\Sigma$ (saturation diagnostic: baseline gain at 100 ep minus baseline gain at 30 ep). Loss variants: \textbf{sns} (scale-normalised Shape-IoU \citep{ding2023}), \textbf{css} (centred-scale Shape-IoU), \textbf{pws} (prior-weighted Shape-IoU). \textbf{3v3} = three-seed-by-three-seed fully-permuted paired runs. Gains are in percentage points (pp) of mAP50-95 unless stated, always same-budget and same-pipeline.

\subsection*{M.2 The joint bound in full: decomposition, corollary, proposition and their conditions}
\addcontentsline{toc}{subsection}{M.2 The joint bound in full: decomposition, corollary, proposition and their conditions}

Three statements --- an exact decomposition with its \textbf{elementary} gap corollary, a \textbf{sharp} sensitivity bound on the optimal-displacement cost that does \textbf{not} improve that corollary, and a \textbf{negative result}: the curvature term the decomposition invites is redundant under the paper's own strong-convexity assumption, and where curvature is genuinely needed the standard top-eigenvalue modulus is invalid, the correct object being the path quantity $\kappa^-_\Delta$. \textbf{The bound is not a validated theory}: no graded claim in this paper rests on it. The derivations, the term-by-term reading of Eq. (3) and the hypotheses H1-H10 are in \textbf{Appendix J} (statements and scope in J.1, proof-level detail in J.3, retraction in J.4).

\subsection*{M.3 The protocol in full: settings, seed layers, registration status and the clean-protocol specification}
\addcontentsline{toc}{subsection}{M.3 The protocol in full: settings, seed layers, registration status and the clean-protocol specification}

\textbf{Figure and table numbering.} Figures are numbered F1--F7 and tables T1--T12 in the released material; in this submission the article carries one figure, \textbf{Fig. 1}, the two-stage protocol schematic. The figures \textbf{carried here} use \textbf{supplementary numbering}: the budget curve (\textbf{Fig. S1}), the shift ordering (\textbf{Fig. S2}), the dual-metric separation (\textbf{Fig. S3}), the measurement-hierarchy schematic (\textbf{Fig. S4}) and the prediction--evidence--status matrix (\textbf{Fig. S5}). \textbf{Only Fig. S3 is cited from the main text}; the other four are not, so their numbering is internal to the supplement. The released \textbf{F3} (shift ordering) and \textbf{F5} (the detector family and its insertion points) are carried here but \textbf{not cited} in either text, and are therefore not numbered. The design table of \S{}8.2 is \textbf{Table S1} (the thirteen domain pairs). Of the released tables, \textbf{T10} is cited in \textbf{this supplement only}, \textbf{T11} in this supplement and \textbf{T12} in the accounting behind C4, which is also part of this supplement; the remaining table numbers (T1--T6 and T9) belong to the released record and to earlier numbering; the tables that exist in this submission are \textbf{T7}, \textbf{T8}, \textbf{T10}, \textbf{T11} and \textbf{T12}, together with \textbf{Table S1}. \textbf{Datasets (introduced here, cited where a canonical source exists).} Thirteen domain pairs are built from ten image corpora, listed in Table S1 by pair; the per-corpus table, with the split sizes actually used (20\% subsets in parentheses), is \textbf{Table L1 (Appendix L)}, which also carries the corpus roles, provenance and the per-corpus split statistics.

One confound is stated in the main text because it conditions a headline cell: the \textbf{smoke corpus is single-class} (\texttt{nc = 1}, "smoke") while the \textbf{firesmoke corpus is two-class} (fire, smoke), so the smoke\ensuremath{\rightarrow}SFCHD line starts from a \emph{single-class} source head adapted to the target's two classes, and the released weights record both stages (\texttt{smoke\_pretrain.pt}, \texttt{nc = 1}; \texttt{smoke\_pretrain\_nc2.pt}, \texttt{nc = 2}). The three corpus roles, the cross-domain pair list, the shift ends, the provenance statements and the per-corpus split statistics are in \textbf{Appendix L}.

\textbf{Two labels the released registry keys runs by, defined here once.} \emph{A pipeline} and \emph{B pipeline} are the two training stacks: A carries the SFCHD-side cells and takes an explicit augmentation-RNG seed (\texttt{a5\_train\_obj\_aug.py}, built on the registered trainer), while B is the registered stack (\texttt{train\_obj.py}, ultralytics 8.4.120). They differ in that augmentation control and in framework bookkeeping, so cross-pipeline \emph{magnitudes} never enter a graded contrast --- every graded contrast is same-pipeline --- and a cross-pipeline number appears only where it is labelled as such: the one order-of-magnitude remark of this kind is flagged in place (\S{}8.2) and carries no graded weight. The seed mechanism is a dual layer --- a \texttt{--shuffle-seed} data-order permutation, which is what the registered \ensuremath{\sigma} varies, plus a separate \texttt{--aug-seed} augmentation-RNG layer. A run recorded \textbf{without} a seed suffix is \emph{unseeded}: no \texttt{--shuffle-seed} permutation was passed, and because the framework's own \texttt{seed} argument is not consumed by the data pipeline on these datasets (\textbf{Appendix D}), such a run has no declared data order --- a comparison against it is therefore a strategy-versus-baseline contrast at one order, not a seed-wise pair --- and the registry's naming conventions, including the \emph{BoxB} label and the two-stage protocol these stacks implement, are in \textbf{Appendix D}; the one mapping fact the argument needs is stated here: \textbf{for the two SFCHD cells both keys point at the same physical directory}, so the reported metric is measured on the split that selects the checkpoint --- as for the \texttt{aitod\_20p}, \texttt{dota15\_20p}, \texttt{dota\_20p} and \texttt{mende\_20p} cells and the source-stage yaml, while \texttt{mask\_20p}, \texttt{firesmoke\_20p} and \texttt{visdrone\_20p} carry a genuine held-out test split. Training and evaluation are disjoint at the image level in every case, so this is not train/test contamination. \textbf{For all seven of these cells the clean three-way protocol replaces that endpoint rule with the corpus's own held-out split} (Appendix M.5).

\textbf{The two rates, and why the baseline one is the one it is.} The baseline arm's peak learning rate is 0.001 and the strategy arm's is 0.005 --- the 5\ensuremath{\times} of the abstract; the dose axis of \S{}8.3 (2.5\ensuremath{\times}/5\ensuremath{\times}/10\ensuremath{\times}/20\ensuremath{\times}) is expressed as multiples of that baseline, so a multiplier always means a multiple of 0.001, and the released registry labels arms accordingly (\texttt{@lr0.001} for baseline-rate arms, \texttt{lr005} for the strategy). This paper does not derive 0.001 from a criterion; what it establishes is that 5\ensuremath{\times} is \textbf{not} a specially favourable point on the axis --- the measured dose response has an interior optimum at both budgets, and the optimum \emph{moves with the budget} (2.5--5\ensuremath{\times} at 30 epochs, 10\ensuremath{\times} at 100 epochs, \S{}8.3) --- so the registered tier is a matched choice rather than the best one.

\textbf{Pre-registration and registration status.} Blinded dual-model predictions were written before any outcome data existed, and the grading criteria were fixed in the internal grading documents after the first same-pipeline observations of the headline cells but before their final 3v3 completions, and never changed afterwards. All documents carry repository timestamps; none were back-dated. The consequence of that ordering --- the registration postdates those first observations, so we call it "prospectively frozen" rather than "pre-registered" --- and the retraction of the registration-time family-control claim are stated below.

\textbf{How the clean protocol was specified, and what it holds fixed.} The three-way protocol was built after the screening measurement that selected these two cells and before any clean-protocol run, which is the ordering the builder's recorded timestamp establishes; individual launch times are not preserved in the release. Its runs are a separate batch from the published ones, driven by yaml files written by that same carve builder. Within a cell the seeds share one source checkpoint --- the release records it per run, together with the deterministic data-order seed --- so initialisation is fixed by design rather than varied, and the initialisation component of the variance is measured separately, on independent source endpoints. What the protocol changes is the endpoint rule and not the arm pairing: the reported value is the target corpus's own held-out test split, which no checkpoint selection can touch, whereas the checkpoint itself is chosen on the validation carve-out taken from the remaining training pool.

\subsection*{M.4 The pair construction and the statistical framework in full}
\addcontentsline{toc}{subsection}{M.4 The pair construction and the statistical framework in full}

Table S1 lists the thirteen domain pairs studied: seven within-domain pairs; three B-pipeline cross-domain pairs (disjoint datasets); \textbf{two SFCHD-pipeline cross-domain pairs, SHWD\ensuremath{\rightarrow}SFCHD and smoke\ensuremath{\rightarrow}SFCHD}; and one post-freeze directional probe (visdrone\ensuremath{\rightarrow}dota15). Eight of the ten B-pipeline target domains are saturated at 30 epochs (\ensuremath{\Sigma} \ensuremath{\approx} 0); the two small-sample pairs (dota15/dota) are not. Shift D is quantified by s-OTDD over frozen features (4,000--5,000 samples/split). The thirteen pairs, their roles and their measured shift --- with the two columns that say what each row can carry, because the table is a design table and not an evidence table: "LR arm paired" is \texttt{no} wherever a cell has a single arm per side, and the tier column is this paper's own grading (\emph{screening} = three-seed reading under the frozen rule, \emph{replication} = the outcome-selected ten-seed extension and its reuse-free seven seeds, \emph{directional} = post-freeze probe, \emph{ungraded} = single-run points used in no graded claim, \textbf{*clean census}* = the clean three-way re-run of the same seven within-domain cells at a held-out endpoint, three seeds each and ten for two of them, which likewise enters no graded or family-level claim and whose per-cell values are in Appendix M.5):

{\footnotesize
\setlength{\tabcolsep}{4pt}
\begin{longtable}{p{0.025\textwidth}p{0.078\textwidth}p{0.082\textwidth}p{0.103\textwidth}p{0.066\textwidth}p{0.041\textwidth}p{0.156\textwidth}p{0.292\textwidth}}
\toprule
\# & pair & category & budgets with paired seeds & registered seeds & D (s-OTDD) & LR arm paired (seeds) & evidence tier \\
\midrule
\endhead
1 & aitod\ensuremath{\rightarrow}aitod & within-domain & 30 & 1 & 3.35 & no (1); clean 3 (s42--s44) & ungraded; \textbf{clean census (n=3)} \\
2 & dota\ensuremath{\rightarrow}dota & within-domain & 30 & 1 & n.m. & no (1); clean 3 (s42--s44) & ungraded; \textbf{clean census (n=3)} \\
3 & dota15\ensuremath{\rightarrow}dota15 & within-domain & 30, 100 & 1 & 2.89 & no (1); clean 10 (s42--s51) & ungraded; \textbf{clean census (n=10) --- a null at 30 ep, negative at 100 ep} \\
4 & firesmoke\ensuremath{\rightarrow}firesmoke & within-domain & 30, 50, 100 & 1--3 & 5.75 & no (1); clean 3 (s42--s44) & ungraded; \textbf{clean census (n=3)} \\
5 & mask\ensuremath{\rightarrow}mask & within-domain & 30, 50, 100 & 1--3 & 6.05 & no (1); clean 3 (s42--s44) & ungraded; \textbf{clean census (n=3)} \\
6 & mendeley\ensuremath{\rightarrow}mendeley & within-domain & 30, 50, 100 & 1--3 & 3.51 & no (1); clean 3 at 30 ep, 10 at 100 ep & ungraded; \textbf{clean census (n=3 / n=10)} \\
7 & visdrone\ensuremath{\rightarrow}visdrone & within-domain & 30 & 1 & 6.43 & no (1); clean 10 (s42--s51) & ungraded; \textbf{clean census (n=10) --- robustly negative} \\
8 & smoke\ensuremath{\rightarrow}firesmoke & cross-domain (B) & 30, 50, 100 & 1--3 & 9.44 & partly (1--3) & ungraded \\
9 & mafa\ensuremath{\rightarrow}mask & cross-domain (B) & 30, 100 & 1 & 7.19 & no (1) & ungraded \\
10 & mafa\ensuremath{\rightarrow}mendeley & cross-domain (B) & 30, 50, 100 & 1--3 & 13.03 & partly (1--3) & ungraded \\
11 & SHWD\ensuremath{\rightarrow}SFCHD & cross-domain (SFCHD) & 30, 50, 100 & 3, 10 & 12.76 & yes (3 \ensuremath{\rightarrow} 10) & \textbf{replication} (corroborating) \\
12 & smoke\ensuremath{\rightarrow}SFCHD & cross-domain (SFCHD) & 100 & 3, 10 & 12.99 & yes (3 \ensuremath{\rightarrow} 10) & \textbf{replication} (strong) \\
13 & visdrone\ensuremath{\rightarrow}dota15 & post-freeze probe & 100 & 10 & 8.02 & yes (10) & \emph{directional} \\
\botrule
\end{longtable}}

(n.m. = not measured for that pair; the remaining within-domain pair shares its object's row. The four pairs registered for the prospective replication of \S{}1 are \textbf{not} reported as results in this paper; those readings belong to the replication, and the roster of the four is in Appendix I, alongside their shift readings.) \textbf{On SHWD\ensuremath{\rightarrow}SFCHD's two roles}: its source endpoint was trained under the A pipeline, so its \emph{absolute} magnitudes are order-comparable with B-pipeline values only, while its \emph{cell contrast} --- lr0.005 against lr0.001 from the same source endpoint, same target subset, same budget, same seeds --- is same-pipeline in the sense that matters for the paired test, and it is on that contrast that the corroborating cell's evidence rests. We therefore treat it as a registered family cell for grading and as a cross-pipeline reference for magnitude comparisons, and the phrase 'reported separately' is withdrawn. The registered family --- what can and cannot be reconstructed from the released material, and the 26-vs-28 discrepancy --- is set out in \textbf{Appendix I}; this paper uses that 26-configuration roster for every family-level quantity and never uses the recorded count of 28 as a denominator.

\textbf{Statistical framework, in one paragraph.} The confirmatory content is the fresh-seed paired tests on the two screening-selected cells: both tiers report the \textbf{paired sign-flip permutation} alongside the paired t, and permutation is the primary test, with its n-dependent floor ($2/2^n$) printed with every value so that a floor-limited p cannot be read as strong evidence. The family-level analysis --- paired t-tests on the 26 configurations of Appendix I (Table T10) with Benjamini--Hochberg at q = 0.05 (Bonferroni, \ensuremath{\alpha}/m = 1.9\ensuremath{\times}10\textsuperscript{-3}, as sensitivity) --- is a \emph{robustness} read and not a family-wise guarantee, because the roster was enumerated after the outcomes and the extension was itself triggered by the screening result. "The frozen criterion" is the frozen effect-size rule stated below, frozen before the final outcomes: a cell passes when the mean 3-seed gain exceeds 3\ensuremath{\times} the strategy-side seed \ensuremath{\sigma}; its role is \emph{screening and tier assignment}, and its null behaviour is derived in Appendix F. Under it the registered tier is a screening result (BH q = 0.29 and 1.00), decisive only in the pre-specified ten-seed extension and on the seven fresh seeds (BH q = 1.7\ensuremath{\times}10\textsuperscript{-8}/8.1\ensuremath{\times}10\textsuperscript{-5} and 4.9\ensuremath{\times}10\textsuperscript{-6}/1.0\ensuremath{\times}10\textsuperscript{-3}); that extension is outcome-selected, which no multiplicity correction repairs. The full statement is in \textbf{Appendix F.1}.

\subsection*{M.5 The clean three-way protocol beyond SFCHD: the per-seed readings}
\addcontentsline{toc}{subsection}{M.5 The clean three-way protocol beyond SFCHD: the per-seed readings}

\textbf{Why this subsection exists.} \S{}M.3 specifies the clean three-way protocol for the two SFCHD cells. The same carve-out builder was also run over the seven within-domain corpora and over two cross-domain pairs that do \textbf{not} involve SFCHD, so that \S{}8.3's sign structure can be read on targets outside the registered family. This subsection carries those per-seed values, and the two released files named in \S{}8.3 (\texttt{clean\_threeway\_perseed.csv}, \texttt{clean\_threeway\_summary.csv}) recompute every number printed here.

\textbf{What is held fixed, and what the endpoint is.} For every cell the training split is the same 20\%-label subset as the published runs; the validation split is a carve-out taken deterministically (seed 42) from the remaining training pool; and the reported value is the target corpus's \textbf{own held-out test split}, so no checkpoint selection can touch the reported split. The endpoint is mAP50-95 (\%) of \texttt{best.pt}, the checkpoint selected on the carve-out --- selected, that is, \textbf{on the mAP@0.5:0.95 column alone} (Ultralytics 8.4.120 scores the fitness \texttt{w = [0, 0, 0, 1]}, so mAP50 does not enter the choice); \texttt{last.pt} is carried in the release as a robustness column and agrees in sign throughout, including the negative cell below. Split sizes were re-measured for this submission, and every pairwise \textbf{name-level intersection is zero} (train / validation carve-out / test): 2243/1000/2804 (aitod20), 282/300/458 (dota15\_20p), 282/300/458 (dota\_20p), 11232/2000/6411 (firesmoke20), 1224/900/919 (mask20), 1150/1000/800 (mende20), 1294/500/1610 (visdrone\_20p). Both arms are this paper's own --- baseline lr 0.001 against strategy lr 0.005 --- so each row is a same-pipeline, same-budget, same-seed pair.

\textbf{The 30-epoch within-domain cells: two arms \ensuremath{\times} three seeds, and ten seeds for two of them.} \ensuremath{\Delta} = strategy \ensuremath{-} baseline, in pp; p is the paired sign-flip permutation, the primary test, printed with its n-dependent floor. The two cells carried to ten seeds are the ones the extension was spent on: dota15 because it carried the smallest effect at three seeds, and visdrone because it is multi-class (ten classes) --- so that the extension is not an extension over single-class corpora.

{\footnotesize
\setlength{\tabcolsep}{4pt}
\begin{longtable}{p{0.226\textwidth}p{0.134\textwidth}p{0.134\textwidth}p{0.071\textwidth}p{0.064\textwidth}p{0.057\textwidth}p{0.106\textwidth}p{0.050\textwidth}}
\toprule
cell & baseline (mean \ensuremath{\pm} \ensuremath{\sigma}) & strategy (mean \ensuremath{\pm} \ensuremath{\sigma}) & \ensuremath{\Delta} & seeds won & paired t & p (permutation) & floor \\
\midrule
\endhead
aitod20 & 10.398 \ensuremath{\pm} 0.180 & 9.488 \ensuremath{\pm} 0.474 & \ensuremath{-}0.909 & 0/3 & \ensuremath{-}5.34 & 0.25 & 0.25 \\
dota15\ensuremath{\rightarrow}dota15 \textbf{(10 seeds)} & 15.214 \ensuremath{\pm} 0.118 & 15.183 \ensuremath{\pm} 0.539 & \textbf{\ensuremath{-}0.031} & \textbf{7/10} & \ensuremath{-}0.17 & 0.871 & 0.00195 \\
dota\ensuremath{\rightarrow}dota & 17.169 \ensuremath{\pm} 0.188 & 16.234 \ensuremath{\pm} 0.351 & \ensuremath{-}0.935 & 0/3 & \ensuremath{-}3.59 & 0.25 & 0.25 \\
firesmoke\ensuremath{\rightarrow}firesmoke & 39.581 \ensuremath{\pm} 0.389 & 35.197 \ensuremath{\pm} 1.139 & \ensuremath{-}4.384 & 0/3 & \ensuremath{-}6.23 & 0.25 & 0.25 \\
mask\ensuremath{\rightarrow}mask & 65.348 \ensuremath{\pm} 0.602 & 64.415 \ensuremath{\pm} 0.909 & \ensuremath{-}0.933 & 0/3 & \ensuremath{-}1.49 & 0.25 & 0.25 \\
mendeley\ensuremath{\rightarrow}mendeley & 73.075 \ensuremath{\pm} 0.278 & 70.171 \ensuremath{\pm} 2.276 & \ensuremath{-}2.903 & 0/3 & \ensuremath{-}2.52 & 0.25 & 0.25 \\
visdrone\ensuremath{\rightarrow}visdrone \textbf{(10 seeds)} & 11.681 \ensuremath{\pm} 0.108 & 10.872 \ensuremath{\pm} 0.354 & \textbf{\ensuremath{-}0.809} & \textbf{0/10} & \ensuremath{-}7.85 & \textbf{0.00195} & 0.00195 \\
\botrule
\end{longtable}}

\textbf{The 100-epoch cells: two arms \ensuremath{\times} ten seeds.} The first two rows are cross-domain; the third is within-domain. \textbf{A permutation $p$ cannot fall below its floor} ($2/2^{n}$: 0.00195 at $n=10$), so where a contrast is more extreme than the floor the table prints the floor; the paired-$t$ $p$ for that contrast would be $3.2\times10^{-11}$, and it is \textbf{not} a permutation result.

{\footnotesize
\setlength{\tabcolsep}{4pt}
\begin{longtable}{p{0.114\textwidth}p{0.087\textwidth}p{0.127\textwidth}p{0.127\textwidth}p{0.040\textwidth}p{0.060\textwidth}p{0.054\textwidth}p{0.168\textwidth}p{0.047\textwidth}}
\toprule
cell & category & baseline (mean \ensuremath{\pm} \ensuremath{\sigma}) & strategy (mean \ensuremath{\pm} \ensuremath{\sigma}) & \ensuremath{\Delta} & seeds won & paired t & p (permutation) & floor \\
\midrule
\endhead
aitod\ensuremath{\rightarrow}visdrone & cross-domain & 10.209 \ensuremath{\pm} 0.116 & 10.397 \ensuremath{\pm} 0.110 & +0.188 & 9/10 & 3.30 & 0.0059 & 0.00195 \\
visdrone\ensuremath{\rightarrow}dota15 & cross-domain & 11.626 \ensuremath{\pm} 0.167 & 15.277 \ensuremath{\pm} 0.253 & +3.650 & 10/10 & 37.75 & \textbf{0.00195} (= the floor) & 0.00195 \\
mendeley\ensuremath{\rightarrow}mendeley & within-domain & 72.549 \ensuremath{\pm} 0.488 & 70.230 \ensuremath{\pm} 1.411 & \ensuremath{-}2.319 & 0/10 & \ensuremath{-}5.69 & 0.0020 & 0.00195 \\
\botrule
\end{longtable}}

\textbf{What this carries, and what it does not.} Three limits. First, five of the seven 30-epoch rows rest on n = 3, where the permutation floor is 2/2\textsuperscript{3} = \textbf{0.25}, so those five cannot reach a conventional level however large their mean shift: they are a \textbf{sign census with seeds}, not a significance claim, and the floor is printed beside every p for exactly that reason --- the same distinction \S{}8.2 draws for the screening tier. Second, the two cells carried to ten seeds are the only 30-epoch rows with the power to separate an effect from noise, and they separate in opposite directions --- visdrone clears the conventional level at the floor, dota15 does not come near it. Third, the cross-domain evidence here is \textbf{two cells}, not a family: both are positive and both clear the conventional level, but two cells support a sign reading on targets outside SFCHD and not a broader population claim.

\textbf{The two cross-domain cells of this table are shared-training cells.} Every row is this paper's own run, with one qualification: for \texttt{visdrone\ensuremath{\rightarrow}dota15} the training recipe and the training set are \textbf{shared with the companion paper on evaluation validity} (its \texttt{T1-c}), and the two papers compute their readings separately under their own selection and endpoint rule. The fingerprint comparison shows the epoch-1 training columns agree bit-for-bit while the epoch-100 columns differ in the fourth and fifth decimal, so what the two papers hold is \textbf{one recipe, one training set and seeds, one implementation, run twice --- a run-level reproduction}; \textbf{neither paper describes it as a cross-implementation replication.} The two papers' values for this cell differ by \textbf{0.077 pp (\ensuremath{\approx}2\%)}, and that difference conflates run-level reproducibility with the endpoint-rule difference: \textbf{we do not attribute it to either}, since separating the two would require re-running both under a single endpoint rule. The same qualification applies to \texttt{aitod\ensuremath{\rightarrow}visdrone}; the within-domain rows are not shared.

\textbf{Reading.} Six of the seven within-domain cells are negative at 30 epochs and one --- dota15, at ten seeds --- is a \textbf{null} (7 of its 10 seeds favour the strategy arm); the within-domain cell read at 100 epochs is negative in all ten; and both cross-domain cells are positive, nine of ten and ten of ten. The dota15 row is reported as a null rather than folded into the census, because at three seeds its mean (\ensuremath{-}0.21 pp) was indistinguishable from the noise the extension then measured (\ensuremath{-}0.03 pp): the extension did not weaken the pattern so much as show which part of it was noise. This is the two-channel structure \S{}8.3 reports, measured at a held-out endpoint on targets that do not include SFCHD. It does not enter the registered family of Appendix I and no family-level quantity uses it.

\textbf{The same within-domain corpora under three conventions.}

{\footnotesize
\setlength{\tabcolsep}{4pt}
\begin{longtable}{p{0.133\textwidth}p{0.216\textwidth}p{0.190\textwidth}p{0.141\textwidth}p{0.214\textwidth}}
\toprule
convention (source) & same-domain cell set & sample size & column / endpoint & reading \\
\midrule
\endhead
\S{}4.6 of the main text (the archive's own endpoint rule) & same-domain controls: \texttt{dota15\ensuremath{\rightarrow}dota15}, \texttt{aitod\ensuremath{\rightarrow}aitod20}, \texttt{visdrone\ensuremath{\rightarrow}visdrone}, \texttt{mask\ensuremath{\rightarrow}mask}, \texttt{dota\ensuremath{\rightarrow}dota} & \textbf{17 cells} & the archive's own endpoint rule & mean \textbf{+0.27 pp}; median \textbf{+0.65}; \textbf{29\% negative}; \textbf{16 of 17} reaching $\|$t$\|$ \ensuremath{\geq} 2 \\
M.5, this appendix (clean three-way protocol, the table above) & the seven 30-epoch within-domain rows & \textbf{7 cells} --- five at \textbf{n = 3}, two (\texttt{dota15\ensuremath{\rightarrow}dota15}, \texttt{visdrone\ensuremath{\rightarrow}visdrone}) at \textbf{n = 10} & target corpus's own held-out test split; mAP50-95 (\%) of \texttt{best.pt} & \textbf{six negative + one null} (the null is \texttt{dota15\ensuremath{\rightarrow}dota15}, 7 of its 10 seeds favour the strategy arm) \\
\S{}8 of the main text (the $U$/$M$ account) & \texttt{aitod\ensuremath{\rightarrow}aitod20} and \texttt{mask\ensuremath{\rightarrow}mask20}, 30 epochs & \textbf{2 cells}, n = 10 each \ensuremath{\dagger} & the \textbf{held-out test} column (family \texttt{r10}, dataset \texttt{*\_3way}) & \textbf{\ensuremath{-}1.195 pp} (t = \ensuremath{-}7.85) and \textbf{\ensuremath{-}2.027 pp} (t = \ensuremath{-}5.44) \\
\botrule
\end{longtable}}

\ensuremath{\dagger} Ten paired seeds per cell, recomputable from \texttt{base/run\_table\_canonical.csv} (family \texttt{r10}, \texttt{aitod20\_3way} / \texttt{mask20\_3way}, 30 epochs): the released table gives \textbf{n = 10} for both, with means \textbf{\ensuremath{-}1.195} and \textbf{\ensuremath{-}2.027 pp}, matching the \S{}8 sentence, which prints the readings and their t values but not n.

\textbf{The three rows are three conventions, not three answers to one question.} They differ in the cell set and in the column: the first averages 17 archive cells under the archive's own endpoint rule, the second reads seven 30-epoch cells on the target corpus's own held-out test split, and the third reads two same-domain cells at 30 epochs in the two-term account. No row is a restatement of another and none contradicts another. The label-axis differences \S{}4.6 reports are rows of the differenced table: \texttt{pcb\ensuremath{\rightarrow}neu\_det} / \texttt{s2df}, 10\%\ensuremath{\rightarrow}50\% of labels, reads \textbf{\ensuremath{-}10.714 pp at MDE 3.654 pp} at 30 epochs (ratio \textbf{2.932}) and \textbf{\ensuremath{-}2.391 pp at MDE 2.395 pp} at 100 epochs (ratio \textbf{0.998}); \texttt{visdrone\ensuremath{\rightarrow}dota15} / \texttt{s2ae}, 10\%\ensuremath{\rightarrow}50\% of labels, reads \textbf{+1.978 pp at MDE 0.4189 pp} (ratio \textbf{4.722}) at 30 epochs and \textbf{+1.193 pp at MDE 0.2658 pp} (ratio \textbf{4.488}) at 100 epochs. On the \texttt{pcb} family the 100-epoch endpoint difference therefore sits \textbf{at} its own MDE while the 30-epoch one clears it by \textbf{2.9\ensuremath{\times}}.

\textbf{Provenance.} These runs are a separate batch from the published ones: the registered stack (\texttt{train\_obj.py}, ultralytics 8.4.120) with the \texttt{shapeiou} loss, each run driven by its own \texttt{*\_3way.yaml} and recorded per run in the released per-seed file together with its two arms and its data-order seed. The batch follows the carve builder and precedes the held-out readout, which is the ordering the released record establishes; individual launch times are not preserved in the release.

\section*{Appendix N. The two modules transferred from the companion paper: detail and provenance}
\addcontentsline{toc}{section}{Appendix N. The two modules transferred from the companion paper: detail and provenance}

\subsubsection*{M.5-bis The clean 200-epoch endpoint: per-run readings and the derivation of the headline}

\textbf{Why this subsection exists.} \S{}4.1 reports that the clean three-way re-run gives a 100-to-200 change of \textbf{\ensuremath{-}1.138 pp} against \textbf{\ensuremath{-}1.143 pp} on the registered column. The 20 runs behind the 200-epoch half are tabulated here \textbf{so that the number can be recomputed from released values rather than taken on assertion}. The per-run file is released as \texttt{provenance/appendix\_E3\_clean/E3\_clean200\_perrun.csv}, and the 100-epoch counterpart as \texttt{E3\_clean100\_perrun.csv}.

Protocol: \texttt{train: train/20\_percent/images}, \texttt{val: carve\_dota15\_20p\_pct/images}, \texttt{test: images/val}, 200 epochs, \texttt{visdrone\_pretrain.pt}, seeds 42--51, two arms (lr0 0.001 against 0.005).

\begin{center}
{\small
\begin{tabular}{rrrr}
\toprule
seed & base (mAP50-95) & 5\ensuremath{\times} (mAP50-95) & paired difference (pp) \\
\midrule
42 & 13.6100 & 16.2700 & +2.6600 \\
43 & 13.4900 & 16.1700 & +2.6800 \\
44 & 13.7000 & 16.4000 & +2.7000 \\
45 & 13.7400 & 16.4400 & +2.7000 \\
46 & 13.8700 & 16.1000 & +2.2300 \\
47 & 13.8100 & 16.4700 & +2.6600 \\
48 & 13.7500 & 16.2400 & +2.4900 \\
49 & 13.5900 & 16.1600 & +2.5700 \\
50 & 13.8800 & 16.1700 & +2.2900 \\
51 & 14.1200 & 16.2800 & +2.1600 \\
\textbf{mean} & 13.7560 & 16.2700 & \textbf{+2.5140} \\
\botrule
\end{tabular}}
\end{center}

The paired gain at this endpoint is \textbf{+2.5140 pp} (SD \textbf{0.2106}, n = 10, t = \textbf{37.75}), and its own MDE is \ensuremath{\gamma}(10)\ensuremath{\cdot}SD = \textbf{0.2095 pp}, i.e. \textbf{12.00\ensuremath{\times} its own MDE} --- the endpoint is resolvable in its own right, not merely as one half of a difference.

\textbf{The like-for-like change.} The clean 100-epoch counterpart reads \textbf{+3.652 pp} on the same protocol (ten seeds, \texttt{E3\_clean100\_perrun.csv}), so the clean-to-clean 100\ensuremath{\rightarrow}200 change is \textbf{-1.1380 pp}, against \textbf{\ensuremath{-}1.143 pp} on the registered column --- the two columns differ by \textbf{0.005 pp}.

\textbf{What this does and does not establish.} It establishes that the 100\ensuremath{\rightarrow}200 change survives a re-run in which checkpoint selection and reporting use different directories. It does \textbf{not} isolate the split: as \S{}4.1 states, the two clean endpoints share one training set and the two registered endpoints share another, and the two groups do not train on identical images. The clean run is a re-run of the same comparison, not an experiment that varies the split alone.

\subsection*{N.1 The registered replication: registration, criterion and archive}
\addcontentsline{toc}{subsection}{N.1 The registered replication: registration, criterion and archive}

\textbf{The registration precedes every run and this is checkable, not asserted.} The protocol was frozen at \textbf{2026-09-13 01:37:27 UTC}, at which point the run directory did not exist, so the criterion could not have been chosen after seeing a result. The registration document is \texttt{<CJK>\_<CJK>\_<CJK>\_20260913.md}; the file records a \textbf{content hash} for the text above its signature line. That value \textbf{is reproducible, and we state the convention it is computed over}: the span is the first 77 lines, LF-normalised, with no trailing LF; the search was by exhaustive enumeration over start point, end point, trailing newline and line-ending convention, and only the two equivalent solutions (first 76 lines plus a trailing LF, equivalently the first 77 lines with no trailing LF) reproduce the recorded value. \textbf{One caveat, recorded because an independent reconstruction depends on it}: "the text above the signature line" is itself a phrase whose boundary has to be resolved, and a reader who resolves it by content rather than by line count lands on a \textbf{different} span and hashes it differently --- the neighbouring readings are one line \emph{shorter} and one line longer; the specification that reproduces the recorded value is the \textbf{line-count} reading given here (77 lines, LF-normalised, no trailing LF), and the primary statement of the registration is quoted above it, so a reader who cannot reproduce the hash should compare against that statement rather than against the hash alone. The full-file hash and the byte count do match the values recorded in this submission, and the freeze timestamp and the absence of any run directory at that time are established independently.

\textbf{What the registration fixes, quoted rather than paraphrased.} Three unsaturated target domains; a 20\% label budget; YOLOv12n with the shape-IoU loss; and a criterion of \textbf{\ensuremath{\geq}2/3 pairs meeting all three of conventional significance (p \textless{} 0.01), at least eight of ten seeds in the same direction, and an effect size of \ensuremath{\geq}+0.30 pp}, under a Benjamini--Hochberg correction at \textbf{q = 0.05}. The archive holds \textbf{71 runs in total} --- three pairs \ensuremath{\times} 2 arms \ensuremath{\times} 10 seeds, a saturated control of 2 \ensuremath{\times} 5, and one near-matched control appended by amendment; an earlier recorded count of 90 does not follow from its own arithmetic, and the count used here is \textbf{71}. \textbf{All 71 are re-keyed to the checkpoint each run itself recorded} --- a property of the run's own record, not of a checkpoint chosen afterwards.

\textbf{The three pairs, as registered and as read.}

{\footnotesize
\setlength{\tabcolsep}{4pt}
\begin{longtable}{p{0.042\textwidth}p{0.118\textwidth}p{0.042\textwidth}p{0.055\textwidth}p{0.062\textwidth}p{0.097\textwidth}p{0.444\textwidth}}
\toprule
pair & registered & \ensuremath{\Delta} (pp) & paired t & p & seeds positive & verdict \\
\midrule
\endhead
T1-a & \texttt{dota15\ensuremath{\rightarrow}aitod} & +0.147 & 1.816 & 0.1027 & 6/10 & \textbf{fails} \\
T1-b & \texttt{aitod\ensuremath{\rightarrow}visdrone} & +0.168 & 4.455 & 0.0016 & 10/10 & significance and direction met; \textbf{fails the +0.30 pp magnitude} \\
T1-c & \texttt{visdrone\ensuremath{\rightarrow}dota15} & +3.727 & 43.774 & 8.5\ensuremath{\times}10\textsuperscript{-12} & 10/10 & \textbf{meets all three parts} \\
\botrule
\end{longtable}}

\textbf{These two pairs are the cells \S{}4.1 re-runs on the epoch axis, so their own-scale MDEs are recorded here.} On the released level table (\texttt{theory\_B\_MDE\_by\_cell.csv}, \texttt{test\_map50\_95}) the 100-epoch cells above carry MDE \textbf{0.2678 pp} (\texttt{visdrone\ensuremath{\rightarrow}dota15} / \texttt{t1c} / 20\% labels / 100 ep) and \textbf{0.1186 pp} (\texttt{aitod\ensuremath{\rightarrow}visdrone} / \texttt{t1b} / 20\% / 100 ep), and their 200-epoch counterparts \textbf{0.1927 pp} (\texttt{visdrone\ensuremath{\rightarrow}dota15} / \texttt{mech} / 20\% / 200 ep) and \textbf{0.1408 pp} (\texttt{aitod\ensuremath{\rightarrow}visdrone} / \texttt{mech} / 20\% / 200 ep) --- so the \texttt{aitod\ensuremath{\rightarrow}visdrone} 200-epoch level of \textbf{\ensuremath{-}0.055 pp} lies inside its own MDE while its 100-epoch level of \textbf{+0.168 pp} does not. The 100\ensuremath{\rightarrow}200 \textbf{change} itself is a \textbf{cross-family} comparison (the 100-epoch cell sits in family \texttt{t1c} / \texttt{t1b}, the 200-epoch cell in \texttt{mech}), so neither released table carries it as a row; the one number computed for it is the \texttt{aitod\ensuremath{\rightarrow}visdrone} contrast in \texttt{theory\_B\_MDE\_diff\_by\_cell\_<CJK>.md} \S{}8 (\textbf{n = 10}, \textbf{mean\_diff = \ensuremath{-}0.2230 pp}, \textbf{mde\_diff = 0.136955 pp}, \textbf{\textbar{}mean\textbar{}/MDE = 1.63\ensuremath{\times}}), quoted there and \textbf{not} in either CSV.

\textbf{Reading the outcome.} After the correction \textbf{two of the three pairs survive at q = 0.05}; on the complete three-part criterion \textbf{one of three} meets it. Across the pairs the effects move in the same direction, with a mean of \textbf{+1.347 pp over the registered seed-pairs} (the mean of the three pair means: +0.147, +0.168 and +3.727), spanning \textbf{\ensuremath{-}0.22 to +4.32 pp}. The outcome is therefore written \textbf{"criterion not met"} and not "hypothesis rejected": the registered direction is stable across all three targets while the registered magnitude is met on one, and that tension is the reported result rather than a defect in it.

\textbf{The limitation of the batch, stated with it.} It ran under the protocol in force at registration, in which \texttt{val} and \texttt{test} coincide for the cells it measures, \textbf{with one exception: T1-b}, where the two columns part by about \textbf{0.98 pp} (val-best \textbf{+1.152} vs test \textbf{+0.168}), so \textbf{it is not a clean-protocol result}: the level at which a checkpoint is read cannot be separated from the checkpoint that selected it. The three-way work of \S{}N.2 is a different batch and does not repair this one retroactively. The honest reading of the registered set is evidence about \textbf{direction on three targets}, not about magnitude on any held-out target.

\subsection*{N.2 The three-way split and the checkpoint-selection audit: the per-run detail}
\addcontentsline{toc}{subsection}{N.2 The three-way split and the checkpoint-selection audit: the per-run detail}

\textbf{The split.} Each run's target corpus is divided into three disjoint parts --- \textbf{33.31 / 29.74 / 29.86 mAP50-95} --- with train, validation carve-out and held-out test sharing \textbf{no file name}, and the reported value taken on the part no selection could see.

\textbf{Re-keying changes nothing, and that is the measurement.} Re-keying every run from its training-time selection rule to this carve-out rule moves the within-domain mean by \textbf{\ensuremath{-}0.069 pp} (t = \ensuremath{-}0.23, \textbf{p = 0.825}, 95\% CI \ensuremath{-}0.758 to +0.620 pp) --- a null, printed as a null. \textbf{806 evaluations over 41 runs} enter this audit --- the re-key test itself is taken over the \textbf{nine paired units} the carve-out rule leaves, which is the count its interval width and \emph{t} imply, of which \textbf{202 points are shared with an independently archived evaluation and agree to the digit}; that overlap is the positive control for the sweep.

\textbf{The premium, decomposed.} Hindsight capturable by selection is \textbf{0.7039 pp}. A rule taking the current maximum captures \textbf{0.2114 pp = 30.0\%} of it; averaging the top five epochs raises the captured share to \textbf{69.6\%} (0.4898 pp) at a selection-side cost of 0.7267 pp. Most of the premium is a \textbf{shape} effect rather than a mislabelled one --- on the selection side shape accounts for \textbf{54.5\%} of it --- which is why the endpoint rule, not the difficulty of the split, is what the reported level is sensitive to.

\textbf{Scope of this module.} The audit measures the \emph{endpoint rule's} contribution on this pipeline. It does not estimate the premium for any other pipeline, and the 0.077 pp difference between this paper's two readings for \texttt{visdrone\ensuremath{\rightarrow}dota15} conflates run-level reproducibility with the endpoint-rule difference and is \textbf{attributed to neither}.

\subsection*{N.3 Provenance of the transferred material}
\addcontentsline{toc}{subsection}{N.3 Provenance of the transferred material}

\textbf{The domain-shift instrument was recomputed rather than quoted.} The sliced-OTDD readings behind the shift-comparison statement were recomputed on this submission's own pipeline, with a positive control: reproducing the published within-domain cell (\texttt{MAFA\ensuremath{\rightarrow}mask}) is recovered at \textbf{7.1894}, matching the published value exactly. The cross-domain values computed for this submission are \texttt{visdrone\ensuremath{\leftrightarrow}dota15} \textbf{8.0226}, \texttt{aitod\ensuremath{\leftrightarrow}visdrone} \textbf{8.7739}, \texttt{dota15\ensuremath{\leftrightarrow}aitod} \textbf{5.1726}, \texttt{dota\ensuremath{\leftrightarrow}dota15} \textbf{0.7386} (its Wasserstein term identically zero), \texttt{dota\ensuremath{\leftrightarrow}aitod} \textbf{5.1627} and \texttt{dota\ensuremath{\leftrightarrow}visdrone} \textbf{8.0162}. \textbf{The published within-domain values are cited as published and are not recomputed or replaced here}, because this pipeline's within-domain readings do not reproduce them (for \texttt{VisDrone} it returns 3.2799 against a published 6.4273), so the instrument is used only where the positive control licenses it. The two readings are reported as measuring different constructions and are not merged.

\textbf{A recorded configuration digest that no longer matches any copy.} The third cell's configuration \textbf{as originally launched} was recorded, but \textbf{the recorded file no longer exists}: the run was relaunched at least once, and two further copies survive. The variants agree on every load-bearing setting (\texttt{model}, \texttt{data}, \texttt{epochs 100}, \texttt{batch 32}, \texttt{imgsz 640}, \texttt{lr0 0.001}, \texttt{workers 16}), so \textbf{configuration attribution is unaffected}; what is affected is that the recorded copy \textbf{cannot be checked byte-for-byte}, which we state rather than offer as verifiable.

\textbf{Environment and settings, to the limit of what was recorded.} The interpreter and package versions were not written into the frozen log, so this submission does not state them. All s-OTDD runs used ResNet-18 \texttt{IMAGENET1K\_V1}, 512-dimensional features, 224\textasciicircum{}2 inputs, \textbf{200 slices}, \texttt{seed = 42}, equal weights over the two label terms, and truncation at \texttt{n = min(n1, n2)}.

\textbf{What the transfer does and does not carry.} The two modules are this paper's from this version onward and are reported here first; the companion paper retains its own independent readings for the one cell the two share, and \textbf{neither paper describes that cell as a cross-implementation replication} --- the fingerprint comparison shows one recipe, one training set, one implementation and one seed set run twice. The companion paper's deprecation of its own \S{}8 means its \texttt{<CJK>} framing for these modules is withdrawn; the underlying released runs remain the checkable record.

\subsection*{N.4 The domain-shift measurement behind the 2\ensuremath{\times}2 design of \S{}5}
\addcontentsline{toc}{subsection}{N.4 The domain-shift measurement behind the 2\ensuremath{\times}2 design of \S{}5}

\textbf{What it is for.} \S{}5 reports a within-pair budget manipulation on two domain-difference tiers and states the measured shift of each tier (\textbf{D = 1.66} for the matched-corpus pair, \textbf{D = 5.37} for the cross-corpus pair, s-OTDD). This appendix gives the measurement those two numbers come from, so that they are not carried on assertion alone.

\textbf{Column.} The 2\ensuremath{\times}2 is reported on the \textbf{val} column, not on the \textbf{test} column the registration froze; the per-run values are \texttt{provenance/deliver/G1\_B<CJK>\_<CJK>run<CJK>.csv}, which covers \textbf{c1 and c2 only} under the column name \texttt{val\_map}, while c3 and c4 exist only as summaries.

\textbf{The instrument.} s-OTDD, reusing the released shift codebase (\texttt{shift\_families.py}) rather than a re-implementation. Protocol points, all fixed before the measurement: ResNet-18 features at 512 dimensions; images resized to 224 with ImageNet normalisation; \textbf{cap 4,000} samples per side; deterministically sorted, first $N$ taken; sliced Wasserstein-2 with \texttt{n\_slices = 200}, \texttt{seed = 42}; and s-OTDD taken as $\sqrt{W_2^2 + \text{label term}}$. $N = 282$ paired for the two tiers of the 2\ensuremath{\times}2; the same-corpus reference row uses $N = 2{,}243$.

\textbf{The measured values, decomposed.}

{\footnotesize
\setlength{\tabcolsep}{4pt}
\begin{longtable}{p{0.268\textwidth}p{0.138\textwidth}>{\raggedleft\arraybackslash}p{0.049\textwidth}>{\raggedleft\arraybackslash}p{0.105\textwidth}>{\raggedleft\arraybackslash}p{0.105\textwidth}>{\raggedleft\arraybackslash}p{0.081\textwidth}>{\raggedleft\arraybackslash}p{0.113\textwidth}}
\toprule
pair & tier & $N$ & feature $W_1$ & feature $W_2$ & label term & \textbf{D (s-OTDD)} \\
\midrule
\endhead
\texttt{dota15\_tr \ensuremath{\rightarrow} dota15\_te} & \textbf{A0, near-zero} & 282 & 0.0854 & 0.1149 & 2.7405 & \textbf{1.6594} \\
\texttt{dota15\_tr \ensuremath{\rightarrow} aitod\_tr} & \textbf{A1, high} & 282 & 0.2380 & 0.2867 & 28.7135 & \textbf{5.3662} \\
\emph{reference} \texttt{aitod\_tr \ensuremath{\rightarrow} aitod\_te} & --- & 2,243 & 0.0570 & 0.0740 & 1.6069 & \textbf{1.2698} \\
\botrule
\end{longtable}}

The two tiers are separated by \ensuremath{\approx} 3.2\ensuremath{\times}, and both same-corpus values (1.66 and 1.27) sit well below the cross-corpus 5.37, so the separation is not an artefact of the tier construction.

\textbf{The positive control, and why it matters here.} The same pipeline re-derived two \textbf{already published} shift values and reproduced each to \textbf{\ensuremath{\Delta} = 0.0000} (\texttt{shwd \ensuremath{\rightarrow} sfchd\_pool} = 11.2658; \texttt{pcb \ensuremath{\rightarrow} neu} = 23.2931). A pipeline that cannot reproduce values it did not originate is not evidence about new ones, so this control is the licence for reading the table above.

\textbf{Where the numbers live.} The full measurement record is \texttt{deliver/G1\_D<CJK>\_20261004.md} in this release; the measurement script is \texttt{code/g1\_measure\_D.py}. Two instruments for shift are used in this paper (\textbf{N.3} states the separation): the one here, and the one behind \S{}4.6's correlation table (M.1). \textbf{They are not interchangeable and are never mixed within one comparison.}

\section*{Appendix O. Related-work detail of the main text}
\addcontentsline{toc}{section}{Appendix O. Related-work detail of the main text}

\begin{quote}
\textbf{Why this appendix exists.} Two families of prior work are cited in the main text in compressed form, so that the article's \S{}2 and \S{}6 carry the positioning without carrying the enumeration. Every entry was read at \textbf{abstract level} through the channels recorded in \S{}O.3; \textbf{none was read as a full text}, and no entry here should be read as "no work exists that does X".
\end{quote}

\subsection*{O.1 The first-order-geometry family, paper by paper (cited from \S{}6)}
\addcontentsline{toc}{subsection}{O.1 The first-order-geometry family, paper by paper (cited from \S{}6)}

{\footnotesize
\setlength{\tabcolsep}{4pt}
\begin{longtable}{p{0.206\textwidth}p{0.413\textwidth}p{0.058\textwidth}p{0.235\textwidth}}
\toprule
Work & Verbatim overlap & Overlap & Difference in layer \\
\midrule
\endhead
\textbf{Andriushchenko, Croce, M\"{u}ller, Hein, Flammarion}, \emph{A Modern Look at the Relationship between Sharpness and Generalization}, ICML 2023 (arXiv:2302.07011) & \emph{"sharpness does not correlate well with generalization but rather with some training parameters like the learning rate"}; \emph{"a consistent negative correlation of sharpness with out-of-distribution error"}; \emph{"the right sharpness measure is highly data-dependent"} --- its setting includes fine-tuning CLIP and BERT & \textbf{substantial} (conclusion, premise, genre) & they \textbf{question the correlation}; we fix a \textbf{pre-registered SNR gate} per cell and report the resulting \textbf{signs}, and add one sign flip produced by the reporting column alone \\
\textbf{Dziugaite et al.}, \emph{In Search of Robust Measures of Generalization}, NeurIPS 2020 (arXiv:2010.11924) & meta-evaluation: unreliable measures \emph{"can obscure failures and successes"} & substantial (meta-evaluation stance) & same stance; different instrument \\
\textbf{Dinh, Pascanu, Bengio, Veness}, \emph{Sharp Minima Can Generalize For Deep Nets}, ICML 2017 (arXiv:1703.04933) & flatness/sharpness reads are \textbf{not invariant to reparametrization} & partial (premise) & we do not re-argue invariance; we measure a \textbf{noise floor} \\
\textbf{Kwon et al.}, \emph{ASAM}, ICML 2021 (arXiv:2102.11600) & adaptive sharpness fixes scale-invariance but \emph{does it capture generalization in modern settings?} & partial (premise) & as above \\
\textbf{Jiang et al.}, \emph{Fantastic Generalization Measures}, ICLR 2020 (arXiv:1912.02178) & the measure family itself & partial (measure layer) & as above \\
\textbf{Musgrave, Belongie, Lim}, ECCV 2020 (arXiv:2003.08505) & metric-learning evaluation caveats & partial (measure layer) & as above \\
\textbf{Dodge et al.}, EMNLP 2019 (arXiv:1909.03004); Dodge et al. 2020 (arXiv:2002.06305) & \textbf{checkpoint-selection} changes reported conclusions & \textbf{closest to (b)} & they vary the \emph{selection} rule across runs; we hold runs and seeds fixed and vary \textbf{only the reported column} on \textbf{one} cell \\
\botrule
\end{longtable}}

\textbf{What remains ours in \S{}6} (and is \emph{not} in any of the above): (1) the \textbf{gate fixed before the data existed}, applied per cell, failing on \textbf{18/18}; (2) the \textbf{\ensuremath{\varepsilon}-scaling exponent 0.56--0.81} read as a \emph{mechanism} (first-order dominated at the customary radii) rather than as another correlation; (3) the single-cell sign flip produced by \textbf{the reporting column alone} (+0.628 pp on the held-out test column vs \ensuremath{-}0.307 pp on the validation-selected best-checkpoint column).

\subsection*{O.2 The budget / scaling-law family, paper by paper (cited from \S{}2)}
\addcontentsline{toc}{subsection}{O.2 The budget / scaling-law family, paper by paper (cited from \S{}2)}

All of these are statements about \textbf{where the optimum sits} (or how it moves) under a budget change --- the layer the main text says is \emph{known}. They differ from this paper in that none reports a \textbf{paired gain} whose \textbf{sign} is read under a budget intervention on the fine-tuning axis.

{\footnotesize
\setlength{\tabcolsep}{4pt}
\begin{longtable}{p{0.368\textwidth}p{0.305\textwidth}p{0.052\textwidth}p{0.187\textwidth}}
\toprule
Work & Verbatim overlap & Axis & Why it is not this measurement \\
\midrule
\endhead
\textbf{Bjorck et al.}, \emph{Scaling Optimal LR Across Token Horizons}, ICLR 2025 (arXiv:2409.19913) & \emph{"the optimal LR changes significantly with token horizon --- longer training necessitates smaller LR"} & pretraining tokens & a law for the \textbf{optimum's location}; no paired gain, no sign \\
\textbf{Li et al.}, \emph{Rethinking the Hyperparameters for Fine-tuning}, ICLR 2020 (arXiv:2002.11770) & \emph{"the effective learning rate \ldots{} sensitive to the similarity between the source domain and target domain"} & source--target similarity & as above \\
\textbf{Mosbach, Andriushchenko, Klakow}, \emph{On the Stability of Fine-tuning BERT}, ICLR 2021 (arXiv:2006.04884) & \emph{"the role of training dataset size per se is orthogonal to fine-tuning stability. What is crucial is rather the number of training iterations"} & iterations (NLP) & mechanism-side support for the \textbf{epoch} axis; not detection, not paired, not a sign \\
\textbf{Li et al.}, 2026 (arXiv:2602.06797); \textbf{Sakai \& Imaizumi}, 2026 (arXiv:2609.35029); \textbf{Shulgin et al.}, 2026 (arXiv:2603.15958); \textbf{Bu et al.}, ICLR 2026 (arXiv:2602.07145); \textbf{Ajroldi et al.}, 2026 (arXiv:2608.28308); \textbf{Rehn et al.}, 2025 (arXiv:2510.20616) & closed-form schedules / phase transitions / transfer conditions under a fixed horizon; \emph{"both undertraining and overtraining regimes"}; \emph{"the existing heuristics for tuning B do not work"} under a fixed compute budget & various (LLM, DP) & all fit or bound an \textbf{optimum}; none reports a paired-gain sign on a detector \\
\textbf{Hernandez, Kaplan, Henighan, McCandlish}, \emph{Scaling Laws for Transfer}, 2021 (arXiv:2102.01293) & \emph{"the effective data transferred is described well in the low data regime by a power-law of parameter count and fine-tuning dataset size"} & \textbf{data} budget & a power law in the \textbf{low-data regime} with \textbf{no sign change and no bound} (its \S{}3.1 asks whether pre-training can harm performance) \\
\textbf{Springer et al.}, 2025 (arXiv:2503.19206) and \textbf{Everett \& Qiu}, 2026 (arXiv:2609.04577) & \emph{"extended pre-training can make models harder to fine-tune"}; \emph{"the preferred learning rate schedule can reverse across the overtraining axis"} & \textbf{pretraining} budget & these \textbf{do} report a reversal --- but on the pretraining axis and on optimizer preference, not on a paired fine-tuning gain \\
\textbf{Zhang, Xu \& Agrawal}, 2026 (arXiv:2606.30795) & \emph{"several variants switch from positive gains at 1\% to neutral or negative effects at larger budgets"}; \emph{"leaving less room for noisier and less stable interventions"} & \textbf{labelled target budget}, detection & the closest work; the difference is one of \textbf{role and criterion} \\
\textbf{Pardo et al.}, \emph{BAOD: Budget-Aware Object Detection}, CVPRW 2019 (arXiv:1904.05443) & budget as a \textbf{constraint the algorithm optimises under} & compute & the opposite role to ours \\
\botrule
\end{longtable}}

\subsection*{O.3 Channel statement for this appendix}
\addcontentsline{toc}{subsection}{O.3 Channel statement for this appendix}

Every entry above was reached through a jump host and read at \textbf{abstract level only}. The appendices record \textbf{two rounds} of search, and the channel state differed between them, so we state both rather than one.

\textbf{First round.} During the positioning check that produced the tables above, \textbf{\texttt{export.arxiv.org/api/query} returned HTTP 429 throughout} (empty responses, which are \textbf{fetch failures, not zero hits} --- confirmed with a known-absent control term), \textbf{OpenAlex exhausted its daily budget}, Semantic Scholar returned 429, and \textbf{no Chinese-language index was reachable}. Three of the eight claims of the novelty review were consequently left at low evidence strength (the same-domain reading, the isolated backbone swap, and the annotation-density moderator).

\textbf{Second round (re-run for those three claims).} A later session re-ran them on a restored channel: \textbf{Crossref, DataCite, the arXiv API, the arXiv landing pages and OpenAlex all answered}, and the arXiv API produced well-formed responses. Two properties of these channels were measured and are recorded because they change how a negative may be read:

\begin{itemize}
\item \textbf{A zero-hit arXiv response is 700--800 bytes} (the absent-term control returns 728 bytes) whereas a response with hits is several thousand; and an HTTP 429 with a one-byte body is a \textbf{throttled fetch}, not a zero-hit. The three are distinguishable only by reading the status code and the length together, and we report the throttled fetches as failures rather than as absences.
\item The arXiv API's \textbf{quoted-phrase} queries (\texttt{abs:"\ldots{}"}, \texttt{all:"\ldots{}"}) return HTTP 200 with 14--23 kB of \textbf{unrelated recent submissions} rather than matches; only unquoted multi-word conjunctions were used as evidence. Semantic Scholar (no key) and the dblp API (bot challenge) remained \textbf{unusable} in the second round, and those two are the highest-coverage channels for detection preprints.
\end{itemize}

\textbf{What the re-run changed, and what it did not.} It returned same-layer prior work the first round had not reached: label statistics used to predict transfer difficulty without training a model \citep{andriushchenko2023}, source-model accuracy used to predict transfer \citep{dziugaite2020}, isolated backbone comparison inside detection \citep{dinh2017,kwon2021}, and the separation of architecture from pretrained weights \citep{hu2023}. The main text cites all five and scopes its claims accordingly (\S{}2, \S{}4.4, \S{}4.6, \S{}9). It did \textbf{not} return, in either round, a paired-gain measurement of the kind this paper reports, nor a significance test on the \textbf{difference} between two backbones' budget gains, nor an annotation-density moderator of the \textbf{sign} of a gain.

Consequently the tables above are statements about \textbf{what was found in the channels reachable at the time}, and the absence of an entry is \textbf{not} evidence that no such work exists. This is a statement about coverage, not a licence: a reader who knows of work in the two unusable channels should read these sections as incomplete.

\section*{Appendix P. The resolution criterion, the size rule, and bounded nulls}
\addcontentsline{toc}{section}{Appendix P. The resolution criterion, the size rule, and bounded nulls}

\subsection*{P.1 One criterion behind the audit and the MDE}
\addcontentsline{toc}{subsection}{P.1 One criterion behind the audit and the MDE}

\textbf{Prior work.} An MDE as a \emph{gate on model-evaluation claims} is not introduced here (Arviv et al., arXiv:2607.08522, 2026; Zhuang, Li \& Fan, arXiv:2605.28873, 2026, a \textbf{paired} MDE budget). This appendix supplies the \textbf{inventory}: the criterion applied cell by cell to a released archive, with the resulting counts.

For a paired design with $n$ seeds and per-seed differences $D_i$, let $\hat\sigma$ be their sample standard deviation and let

\[\gamma(n)\;:=\;\frac{t_{1-\alpha/2,\,n-1}+t_{1-\beta,\,n-1}}{\sqrt n},\qquad \alpha=0.05,\ \beta=0.20 .\]

Then the minimum detectable difference is $\mathrm{MDE}(n,\hat\sigma)=\gamma(n)\hat\sigma$, and a cell is \textbf{resolvable} exactly when $|\bar D|/\hat\sigma\ge\gamma(n)$. The right-hand side depends on $n$ alone and the left-hand side is a sample-size-free standardised effect: the audit's threshold and the MDE of \S{}4.2 are the same object at two sample sizes.

\begin{center}
{\small
\begin{tabular}{llllllll}
\toprule
$n$ & 3 & 5 & 8 & 9 & \textbf{10} & 13 & 20 \\
\midrule
$\gamma(n)$ & 3.0965 & 1.6625 & 1.1528 & 1.0650 & \textbf{0.9947} & 0.8463 & 0.6605 \\
\botrule
\end{tabular}}
\end{center}

Calibration: $n=10$, $\hat\sigma=0.3080$ pp gives $\mathrm{MDE}=0.3064$ pp, matching the article's 0.306 pp. The per-cell counts of \S{}9 follow from this table applied row by row: \textbf{147} cells scored, \textbf{43} with $|\bar D|<\mathrm{MDE}$, of which \textbf{32 of 73} at $n=3$ and \textbf{11 of 74} at $n>3$.

\textbf{The two seed counts, read correctly.} The median $|\bar D|/\hat\sigma$ is \textbf{4.37} among the $n=3$ cells and \textbf{3.03} among the $n=10$ cells, so the two groups differ in \textbf{standardised effect} as well as in size: the 44\%-against-15\% contrast describes two groups rather than showing that more seeds resolve more. Re-basing the $n=3$ cells to $n=10$ at their own $\hat\sigma$ leaves \textbf{23\%} unresolvable against the \textbf{14\%} observed at $n=10$; the gap is the selection-effect part of the contrast, which we do not attribute to either cause.

\subsection*{P.1-bis The second released per-cell table: the differenced readings}
\addcontentsline{toc}{subsection}{P.1-bis The second released per-cell table: the differenced readings}

\textbf{Two released tables, two objects.} Both are computed over the same released runs (base \texttt{run\_table\_canonical.csv}) through the same cell parser (\texttt{scripts/\_cells.py}), and are \textbf{not the same quantity}:

\begin{itemize}
\item \texttt{theory\_B\_MDE\_by\_cell.csv} --- \textbf{147 rows}, the \textbf{level} quantities: one row per cell \texttt{(pair, family, label budget, epoch budget)} on the \texttt{test\_map50\_95} column with at least three paired seeds, and \texttt{sd} is the SD of that \textbf{one cell's own} per-seed paired differences $d_s$ (a single budget against its own baseline).
\item \texttt{theory\_B\_MDE\_diff\_by\_cell.csv} --- \textbf{402 rows}, the \textbf{differenced} quantities: one row per \textbf{comparison of two budgets} inside one \texttt{(pair, family)} group, with $D_s=d_s(\text{to})-d_s(\text{from})$, $\texttt{mean\_diff}=\overline D$, $\texttt{sd\_diff}=\mathrm{SD}(D)$ and $\texttt{mde\_diff}=\gamma(n_{\text{shared}})\cdot\texttt{sd\_diff}$, under the same $\gamma(n)$ and the same \texttt{resolvable} rule $|\overline D|\ge\texttt{mde\_diff}$ as \S{}P.1. It carries an \texttt{is\_endpoint} column (the earliest \ensuremath{\rightarrow} latest comparison inside a group) and both reporting columns.
\item \texttt{theory\_B\_MDE\_valbest\_by\_cell.csv} --- \textbf{152 rows}, the \textbf{level} quantities again but read on the \textbf{\texttt{best\_map50\_95}} column, joinable row by row with the first table on \texttt{(pair, family, label budget, epoch budget)}. Its positive control reproduces the \S{}4.3 cell (\texttt{mask} \ensuremath{\rightarrow} \texttt{mende}, \texttt{t2}, 20, 100 epochs, n = 10) at \textbf{mean \ensuremath{-}1.4275 pp, MDE 0.6708 pp, resolvable} --- \textbf{2.13\ensuremath{\times}} its own MDE on this column against \textbf{0.96\ensuremath{\times}} on \texttt{test\_map50\_95}. Its coverage is a \textbf{superset}: \textbf{147 cells shared, none unique to the test column, five unique to this one}. \textbf{Of the 147 shared cells, 23 change verdict (10 from resolvable to not, 13 the other way), so 15.6\% are column-dependent --- a different denominator from the \textasciitilde{}11\% quoted in the article, which is a fraction of the \emph{gated} cells per epoch budget rather than of the 147} --- and the median SD of the per-seed paired differences rises from 0.2744 to 0.2898. A released table is never a statement about a pair alone but about a pair \textbf{on a named column}; both are now released.
\end{itemize}

\textbf{The counts, recomputed from the CSV.} On \texttt{column=test\_map50\_95} the differenced table holds \textbf{199} rows, of which \texttt{resolvable} is \textbf{True 52 / False 147}; \textbf{61} of the 199 are \texttt{is\_endpoint=True} (\textbf{True 22 / False 39}) and the remaining \textbf{138} are interior budget pairs (\textbf{True 30 / False 108}). The other \textbf{203} rows are the \texttt{best\_map50\_95} column, which the 147-row level table never covers, so 199 + 203 = 402.

\textbf{Median effect against its own MDE.} The table's own ratio column gives a median $|\texttt{mean\_diff}|/\texttt{mde\_diff}$ of \textbf{0.4911} on the test column and \textbf{0.4733} on the val-best column (\texttt{theory\_B\_MDE\_diff\_by\_cell\_<CJK>.md} \S{}5): \textbf{more than half of the budget-to-budget differences fall inside their own MDE}. That is the contrast-level counterpart of the \textbf{43 of 147} level count of \S{}P.1, and it is \textbf{not} comparable to it by subtraction, because the two tables score different objects against different denominators.

\textbf{Recomputable path.} With \texttt{code/\_cells.py} importable, the per-cell read is \texttt{\_cells.arch\_cell(G, pair, ep, backbone, key='test\_map50\_95', lb=\textless{}label budget\textgreater{})} \ensuremath{\rightarrow} per-seed $d_s$ \ensuremath{\rightarrow} shared seeds \ensuremath{\rightarrow} $D_s=d_s(\text{to})-d_s(\text{from})$ \ensuremath{\rightarrow} \texttt{mean\_diff}/\texttt{sd\_diff}/$\gamma(n)\cdot\texttt{sd\_diff}$, $\gamma(n)$ as in \S{}P.1; the two tables share the key \texttt{(pair, family, label budget, epoch budget)} and merge row by row. Two boundary statements (\texttt{theory\_B\_MDE\_diff\_by\_cell\_<CJK>.md} \S{}8): a contrast whose two budgets sit in \textbf{different families} (the \texttt{aitod\ensuremath{\rightarrow}visdrone} 100\ensuremath{\rightarrow}200 change of \S{}4.1) is outside the within-group construction, and the \textbf{structural} contrasts of \S{}4.4 are outside it by design.

\subsection*{P.2 Why ten seeds, and why that is not an ex-ante rule}
\addcontentsline{toc}{subsection}{P.2 Why ten seeds, and why that is not an ex-ante rule}

The registered magnitude bar is $+0.30$ pp ($\S3$; carried over from the frozen criterion of \texttt{<CJK>\_P1\_v4\_20260927.md} \S{}1). At the archive's \textbf{median} $\hat\sigma=0.2744$ pp:

\begin{center}
{\small
\begin{tabular}{lllll}
\toprule
$n$ & 8 & \textbf{9} & 10 & 11 \\
\midrule
$\gamma(n)\hat\sigma$ & 0.3163 & \textbf{0.2922} & 0.2729 & 0.2571 \\
\botrule
\end{tabular}}
\end{center}

so \textbf{$n=9$ is the smallest integer at which a 0.30 pp effect is resolvable at 80\% power against the archive's median noise}, and the ten-seed bar sits \textbf{one seed above that minimum} rather than at it --- which is why the bar must not be read as derived from the noise. Two qualifications. First, $\hat\sigma$ is known only after a cell has been run, so this is a \textbf{retrospective account of the size chosen, not an ex-ante rule}. Second, the components available here are the \textbf{data-order} components, with the augmentation component (0.246 pp on the strategy arm) excluded by the registered design (\S{}F.1), so $\hat\sigma$ is a \textbf{lower bound on the noise} and the bars built on it are \textbf{looser} than nominal; the $n=3$ estimate of $\hat\sigma$ has a CI of $[0.5\times,6.3\times]$ (\S{}F).

\subsection*{P.3 Bounded nulls, and the standard they meet}
\addcontentsline{toc}{subsection}{P.3 Bounded nulls, and the standard they meet}

A \textbf{bounded null} for a modulation contrast is the statement $|\bar D|\pm \mathrm{HB}\subseteq(-b,b)$, $\mathrm{HB}=t_{1-\alpha/2,n-1}\hat\sigma/\sqrt n$. \S{}4.2's second row is the archive's clearest case, and the two numbers it supports are different:

{\footnotesize
\setlength{\tabcolsep}{4pt}
\begin{longtable}{p{0.405\textwidth}p{0.090\textwidth}p{0.435\textwidth}}
\toprule
quantity & value & what it is \\
\midrule
\endhead
observed spread of the four point estimates & \textbf{0.12 pp} & a description of this sweep (0.694 to 0.812) \\
95\% half-width at $n=10$, $\hat\sigma=0.321$ pp & \textbf{\ensuremath{\pm}0.23 pp} & the inference bound this row can carry \\
80\%-power MDE of the 30-versus-200 endpoint difference & \textbf{0.306 pp} & a different statistic: a \emph{contrast}, $\hat\sigma=0.308$ pp \\
\botrule
\end{longtable}}

\textbf{Only two of the 45 paired epoch-axis differences reach an MDE of 0.31 pp or below}: the cell above and a same-domain cell (\texttt{dota15\ensuremath{\rightarrow}dota15}, 100\ensuremath{\rightarrow}200 epochs, $-$0.066 pp at MDE 0.229 pp) --- so a bounded null at this standard is a result this design produced at one cross-domain cell, not archive-wide. All three rows are in the differenced table at their own scale: \texttt{shwd2sf\ensuremath{\rightarrow}sfchd} / \texttt{b2} / 20\% labels / 30\ensuremath{\rightarrow}200 ep carries the endpoint difference \textbf{\ensuremath{-}0.018 pp at MDE 0.3064 pp} (\texttt{ratio\_abs\_mean\_over\_mde} \textbf{0.0588}), the same-domain \texttt{dota15\ensuremath{\rightarrow}dota15} / \texttt{g3} / 20\% / 100\ensuremath{\rightarrow}200 ep \textbf{\ensuremath{-}0.066 pp at MDE 0.2293 pp} (\textbf{0.2878}), and the first row of \S{}4.2, \texttt{smoke\ensuremath{\rightarrow}sfchd} / \texttt{b2} / 20\% / 30\ensuremath{\rightarrow}200 ep, \textbf{\ensuremath{-}2.161 pp at MDE 0.3880 pp} (\textbf{5.569}) --- so the two epoch-axis endpoints at or below 0.31 pp are exactly the two rows \S{}4.2 already quotes, and the decreasing-gain row clears its own MDE by \textbf{5.6\ensuremath{\times}}.

\subsection*{P.4 What this appendix does not claim}
\addcontentsline{toc}{subsection}{P.4 What this appendix does not claim}

It does not calibrate the bars: the noise scale is a lower bound (\S{}P.2), so the thresholds are looser than nominal. It does not make the size rule prospective (\S{}P.2 is retrospective). And it does not extend the audit to the differenced readings, which have their own statistics and are carried row by row in the second released per-cell table (\S{}P.1-bis).

\subsection*{P.5 Scope of the two instruments, and the audit's own limits}
\addcontentsline{toc}{subsection}{P.5 Scope of the two instruments, and the audit's own limits}

The article uses two screening instruments. The \textbf{seed gate} (\S{}4.5) asks whether a \textbf{sign is stable across seeds} --- a difference enters the counts only if it rests on ten paired seeds, or on at least five with a unanimous sign, or reaches $|t|\ge4$. The \textbf{resolution audit} (\S{}9) asks whether a \textbf{single-budget level} is distinguishable from zero at 80\% power. Neither implies the other, and both are \textbf{filters or descriptions rather than tests}: neither carries a nominal error rate, and the family-level BH correction is robustness only. A cell may pass the gate and fail the audit; the gate's eligibility is then unaffected and only that cell's \textbf{magnitude} reading is withdrawn. Three further limits belong to the audit:

\begin{enumerate}
\item \textbf{It covers levels, not contrasts.} Every conclusion of \S{}4 is stated on a \textbf{difference} between two budgets or two backbones, and a difference has its own paired statistic and its own MDE, neither of which is in the level table of \S{}P.1, and both of which are carried row by row in the differenced table of \S{}P.1-bis. \S{}4.2 quotes one such number explicitly (SD of the per-seed endpoint differences 0.308 pp, hence MDE 0.306 pp); the audit's verdicts do not transfer to the differenced readings. For \S{}4.4's \textbf{second backbone} the differenced table carries the epoch-axis contrast, \texttt{y11\_shwd\ensuremath{\rightarrow}sfchd} / \texttt{y11} / 20\% / 30\ensuremath{\rightarrow}100 ep reading \textbf{\ensuremath{-}1.277 pp at MDE 1.255 pp} on the test column (ratio \textbf{1.017}) and \textbf{\ensuremath{-}1.280 pp at MDE 1.250 pp} on the val-best column (ratio \textbf{1.024}), and the two-end MDE of that row (\textbf{0.2787 pp} on the test column) is the yardstick inside which \S{}4.4 places the 100-epoch level. The \textbf{cross-backbone} contrast of \S{}4.4 --- one corpus, one budget, two released backbones --- is a \textbf{structural} comparison and is deliberately outside both tables, as \texttt{theory\_B\_MDE\_diff\_by\_cell\_<CJK>.md} \S{}8.2 records.
\item \textbf{Its roster is a census.} The 147 rows are every cell on the \texttt{test\_map50\_95} column with at least three paired seeds: no analysable cell is missing and no row falls outside that definition (recomputed and checked by \texttt{scripts/43\_A24\_power\_audit.py}). Because each row is an independent per-cell decision, the audit needs no multiplicity correction of its own; because the roster was assembled after outcomes were known, the audit is nonetheless a \textbf{description of this archive} and not a pre-registered family.
\item \textbf{Its threshold is uncalibrated.} The noise scale it divides by is a \textbf{lower bound} on the total noise (\S{}P.2), so the audit flags \textbf{fewer} cells than a fully calibrated threshold would, and "resolvable" here means resolvable against the noise components this design measured.
\end{enumerate}

\subsection*{P.6 The held-out sign prediction of \S{}4.5, family by family}
\addcontentsline{toc}{subsection}{P.6 The held-out sign prediction of \S{}4.5, family by family}

This table carries the per-family detail behind the 21-of-25 figure so that the failures are visible individually rather than as a residual. \textbf{Method.} For every \texttt{(pair, family)} with readings at two or more epoch budgets, the epoch-axis endpoint difference is computed within each \textbf{label budget} slice (earliest to latest epoch available, paired on the shared seeds, \texttt{test\_map50\_95}). The \textbf{lowest label budget} of each family fixes the predicted sign; every \textbf{higher} label budget is held out and scored against it. The split is fixed by budget order and not by outcomes. The companion exercise holds out the epoch budget instead: the label-axis endpoint difference at 30 epochs predicts the sign at 100 epochs.

\textbf{Result.} 21 of 25 held-out slices correct; 4 of the 6 families carrying all five label budgets unanimous; 7 of 7 families correct on the cross-epoch label-axis exercise. The independent units are the \textbf{families} (seven), not the slices.

{\footnotesize
\setlength{\tabcolsep}{4pt}
\begin{longtable}{p{0.178\textwidth}p{0.100\textwidth}p{0.144\textwidth}p{0.105\textwidth}>{\raggedleft\arraybackslash}p{0.117\textwidth}>{\raggedleft\arraybackslash}p{0.233\textwidth}}
\toprule
pair / family & label budgets used & sign fixed at lowest slice & \textbf{held-out slices} & \textbf{predicted correct} & verdict \\
\midrule
\endhead
\texttt{pcb\ensuremath{\rightarrow}neu\_det} / \texttt{s2df} & 10/20/30/40/50 & negative & 4 & \textbf{4} & all five same sign \\
\texttt{chv\ensuremath{\rightarrow}gdut\_hwd} / \texttt{s2hv} & 10/20/30/40/50 & negative & 4 & \textbf{4} & all five same sign \\
\texttt{fsin\ensuremath{\rightarrow}smoke\_keremberke} / \texttt{s2sm} & 10/20/30/40/50 & positive & 4 & \textbf{4} & all five same sign \\
\texttt{visdrone\ensuremath{\rightarrow}dota15} / \texttt{s2ae} & 10/20/30/40/50 & negative & 4 & \textbf{4} & all five same sign \\
\texttt{mafa\ensuremath{\rightarrow}mask\_clean} / \texttt{s2mk} & 10/20/30/40/50 & \texttt{+0.057} at 10\% & 4 & 3 & fails at 30\% (\texttt{\ensuremath{-}0.547}) \\
\texttt{shwd\ensuremath{\rightarrow}sfchd} / \texttt{a0} & 10/20/30/40/50 & \texttt{\ensuremath{-}0.200} at 10\% & 4 & 2 & fails at 40\% (\texttt{+0.233}) and 50\% (\texttt{+0.327}) \\
\texttt{shwd2sf\ensuremath{\rightarrow}sfchd} / \texttt{a0} & 20/40 & \texttt{\ensuremath{-}0.320} at 20\% & 1 & 0 & fails at 40\% (\texttt{+0.097}) \\
\textbf{total} & --- & --- & \textbf{25} & \textbf{21} & --- \\
\botrule
\end{longtable}}

\textbf{Falsification condition, and its current status.} The exercise is refuted in a family whose two label-budget slices carry endpoint differences of \textbf{opposite sign that are each separately resolvable} ($|d|$ at least its own MDE). \textbf{None} of the failures above meets that: each sign-reversing pair has at least one slice below its own 80\%-power MDE under the $\gamma(n)$ convention of Appendix P, and the shortfalls are large (\texttt{mafa\ensuremath{\rightarrow}mask\_clean} 10\% by about 75, 30\% by about 8). The two shortfalls are rows of the differenced table: \texttt{mafa\ensuremath{\rightarrow}mask\_clean} / \texttt{s2mk} at 10\% labels, 30\ensuremath{\rightarrow}100 ep, reads \textbf{+0.057 pp at MDE 4.271 pp} (\texttt{ratio\_abs\_mean\_over\_mde} \textbf{0.01327}) and its 30\% slice reads \textbf{\ensuremath{-}0.547 pp at MDE 4.459 pp} (\textbf{0.12261}). The falsification condition is therefore \textbf{not yet instantiated on this archive}; what the failures establish is the weaker statement that the sign is not unanimous within a family, which is why the result is stated as "four of the six families that carry all five label budgets unanimous, with named sign-reversing exceptions" --- not "the sign is a property of the pair".

\textbf{One limitation of the test itself.} The test was decided after the archive existed. Its split is prior (budget order, outcome-blind), but its existence is not, and we record that distinction rather than claiming a pre-registered prediction.

\subsection*{P.7 The saturated control, per batch and per column}
\addcontentsline{toc}{subsection}{P.7 The saturated control, per batch and per column}

\S{}4.3 reports a damage on one column and its independent batch on the same column; this table carries the four numbers so that the column dependence can be audited.

{\footnotesize
\setlength{\tabcolsep}{4pt}
\begin{longtable}{p{0.053\textwidth}p{0.053\textwidth}p{0.105\textwidth}p{0.053\textwidth}p{0.289\textwidth}p{0.324\textwidth}}
\toprule
batch & family & label budget & epochs & \texttt{test\_map50\_95} (held-out test) & \texttt{best\_map50\_95} (validation-selected) \\
\midrule
\endhead
first & \texttt{t2} & 20\% & 100 & \ensuremath{-}1.156 pp (t = \ensuremath{-}3.03, n = 10) & \textbf{\ensuremath{-}1.428 pp (t = \ensuremath{-}6.69, n = 10)} \\
second & \texttt{r10} & 20\% & 100 & \textbf{\ensuremath{-}1.170 pp (t = \ensuremath{-}3.39, n = 13)} & +0.004 pp (t = +0.02, n = 13) \\
\botrule
\end{longtable}}

\textbf{How to read it.} \S{}4.3 declares the validation-selected column; on it the first batch supplies the damage while the second supplies a \textbf{null}, and the two batches agree only on the held-out test column. Neither column is wrong --- the article's \S{}3 convention assigns val-best to \S{}4.3 and \S{}4.4 while \S{}4.1--\S{}4.2 use the test column. The same dependence is what \S{}6(b) reports for the primary pair, where one cell reads \textbf{+0.628 pp} on test and \textbf{\ensuremath{-}0.307 pp} on val-best with runs, seeds and fine-tuning untouched. The saturated-target damage is therefore \textbf{carried by one batch on one column}, the independent batch supporting it only on the other. The two test-column readings are level quantities, scored in \S{}P.1's table: \texttt{mask\ensuremath{\rightarrow}mende} / \texttt{t2} / 20\% / 100 ep reads \textbf{\ensuremath{-}1.156 pp at MDE 1.199 pp} (\texttt{resolvable} False) and the independent second batch \texttt{mask\ensuremath{\rightarrow}mende20} / \texttt{r10} / no budget suffix / 100 ep reads \textbf{\ensuremath{-}1.170 pp at MDE 1.054 pp} (\texttt{resolvable} True). The \textbf{declared val-best level has no row in either released table} (a level, not a contrast); its own MDE is in \texttt{theory\_B\_MDE\_diff\_by\_cell\_<CJK>.md} \S{}8.4 (\textbf{MDE 0.6708 pp}, \textbar{}mean\textbar{}/MDE \textbf{2.13}, against the same cell's test column at \textbf{0.96}).

\subsection*{P.8 Resolvability as a function of the reporting column}
\addcontentsline{toc}{subsection}{P.8 Resolvability as a function of the reporting column}

\textbf{Prior work.} No discovery of the bias is claimed. The increment is narrower and countable: the cell set is fixed by \texttt{(pair, family, label budget, epoch budget)} over cells whose dataset name carries an \textbf{explicit label budget} (the derived field must be non-null, which excludes the \texttt{r10}-family cells whose names have no budget suffix); the disagreement between two columns is resolved \textbf{against each cell's own MDE}, and the twelve disagreeing cells (five at 30 epochs and seven at 100) are then decomposed \textbf{exhaustively} into the point-estimate channel and the variance channel.

\textbf{The object.} For a cell $c=(\text{pair},\text{family},\text{budget},\text{epochs})$ we take the per-seed paired gain $d_s = 100\,[\text{lr005}_s-\text{base}_s]$ in percentage points, its mean $\bar\Delta$, its sample SD $\hat\sigma_d$ over the common seeds, and the same 80\%-power minimum detectable difference used by the audit of \S{}9 and Appendix P.1, $\mathrm{MDE} = (t_{1-\alpha/2,\,n-1}+t_{1-\beta,\,n-1})\,\hat\sigma_d/\sqrt n$ with $\alpha=0.05$, $\beta=0.20$. A cell is \textbf{resolvable} when $|\bar\Delta|\ge\mathrm{MDE}$. The exercise runs once per column, on cells where \textbf{both} columns are complete and $n\ge 3$; the two columns are never mixed.

\textbf{The count.}

{\footnotesize
\setlength{\tabcolsep}{4pt}
\begin{longtable}{p{0.124\textwidth}>{\raggedleft\arraybackslash}p{0.332\textwidth}>{\raggedleft\arraybackslash}p{0.363\textwidth}>{\raggedleft\arraybackslash}p{0.093\textwidth}}
\toprule
epoch budget & cells with both columns, $n\ge3$ & \textbf{verdict changes with the column} & share \\
\midrule
\endhead
30 & 45 & \textbf{5} & \textbf{11.1\%} \\
100 & 63 & \textbf{7} & \textbf{11.1\%} \\
\botrule
\end{longtable}}

Two features matter. First, the \textbf{share is the same at both budgets (about 11\%) even though the columns' mean gap is not}: the difference between the columns' cell gains, $\bar\Delta_{\text{test}}-\bar\Delta_{\text{val-best}}$, is $-0.043$ pp at 30 epochs ($t=-0.46$, $p=0.65$) against $-0.482$ pp at 100 epochs ($t=-3.12$, $p=0.003$) --- the verdict flip and the mean gap are \textbf{two different statements} and are not combined here. Second, \textbf{all seven} of the 100-epoch verdict changes run in the same direction: \textbf{unresolvable on the test column and resolvable on the val-best column}; \textbf{none moves the other way}.

\textbf{Reporting clusters, not independent experiments.} The seven 100-epoch cells fall in \textbf{four $(\text{pair},\text{family})$ clusters} (three share one pair and family; slices inside a family share data source, pretraining and seed set). They are \textbf{reporting strata, not four independently validated replications}: the released run table carries no image-membership manifest or sample-content hash, so neither cluster independence nor the absence of image-level overlap is established. The count of evidence is \textbf{four reporting strata}, not seven.

\textbf{Robustness.} Raising the minimum seed count from three to six moves the 100-epoch count only from $7/63$ to $4/28$, i.e. the share rises from 11.1\% to 14.3\%; the pattern is therefore not an artefact of the lowest-inclusion cells. At 30 epochs the 45 cells give 5 flips, so the phenomenon is present at a short budget as well, and no reading across budgets is offered.

\textbf{Limits.} (i) The cells are \textbf{not matched across epoch budgets} --- the 45 at 30 epochs and the 63 at 100 are different (pair, family, budget) sets --- so nothing here is monotone in the budget. (ii) The cluster count is \textbf{four}, which bounds the weight the count carries: reporting strata, not independent experiments. (iii) The base table records a \emph{seed kind} per run, and in \textbf{9 of the 165 cells that carry both arms} the two arms' seed kinds are \textbf{not the same set (165 is the two-arm count over cells with a parseable label budget, three fewer than the 168 two-arm cells P.9 reports before that filter)}; here "the same seed" is therefore a pairing asserted from the run records rather than one verified independently, which is why the verdict change is reported as an observed property of the two columns and is not attributed to a mechanism.

\textbf{Why it is checkable.} The released per-run tables carry both columns for every run, so a reader can recompute both verdicts per cell without re-running anything. Where the two columns' paired-difference spreads differ, the \textbf{same} $\bar\Delta$ crosses one column's threshold and not the other's; that is the sense in which the verdict belongs to the convention.

\subsection*{P.9 Boundaries of the released cell definition, and what they do not affect}
\addcontentsline{toc}{subsection}{P.9 Boundaries of the released cell definition, and what they do not affect}

The per-run tables are keyed on a cell defined as (pair, family, label budget, epoch budget). Four properties bound how far a cell can be read; each constrains \emph{interpretation} rather than the arithmetic, and none changes any reported number.

\textbf{1. The cell key does not carry the box-regression loss prior.} \texttt{loss} is one of the axes this paper varies, but it is not part of the key. Measured on the released table: in \textbf{40} (cell, seed, arm) slots the same key holds runs with \textbf{more than one} \texttt{loss} value, and in \textbf{40 of 40} of those slots the readings differ between the two losses. (The same phenomenon counted per (cell, seed), pooling both arms, occurs in \textbf{20} slots; we use the per-arm count because that is the unit the parser selects in.) Selection inside a slot is deterministic but was, until this round, incidental: the released parser breaks ties by run-name length, and the \texttt{shapeiou} records are the shortest, so it took \texttt{shapeiou} in every observed case. That rule is now explicit in the parser (\texttt{shapeiou} preferred, then the outcome rank, then name length), and the explicit rule reproduces the released cell values exactly: \textbf{0 of 338 cells (both columns) change}. The boundary to carry forward is that a cell in this release is a \texttt{shapeiou} cell, and a comparison across loss priors is not available at this key.

\textbf{2. The three-way-protocol flag is collinear with one family.} Of the \textbf{389} runs carrying \texttt{threeway = 1}, \textbf{363 (93.3\%)} sit in the family \texttt{r10} and the remaining 26 in \texttt{g4}. The flag is therefore close to a family label in this archive, so any contrast drawn between the clean three-way protocol and the published protocol is at the same time a contrast between \texttt{r10} and the other families. This does not weaken the protocol argument, which rests on the construction of the split rather than on a between-family contrast, but the two readings are not separable here and should not be presented as if they were.

\textbf{3. A residual set of runs records a seed that its run name contradicts.} Those runs are named with one seed but carry another in their configuration; they are excluded from every seed-gated count in this paper, and no reported number rests on them --- \textbf{except that one of them is itself a released cell of the archive, which is why the exclusion is stated here rather than left implicit}. The residual is small and we do not resolve it here.

\textbf{4. What fraction of the label axis is actually labelled.} The label budget is the measured independent variable of \S{}4 and \S{}4.6, so how much of the archive carries a \emph{checkable} budget label is a bound on that axis. Of the \textbf{2,362} rows in the released table (2,336 distinct runs; the same base as \S{}3), \textbf{453 (19.2\%) carry no parseable label budget}. \textbf{442 of those are recoverable to 20\%} from configuration evidence --- a saved dataset YAML whose \texttt{train} path contains \texttt{20\_percent}, or the generated split specification together with its execution log. (The recovery itself is carried out on the analysis table of 2,570 run-rows that the accompanying analysis package ships, a superset of the 2,362 runs above; \textbf{266} of the recovered rows are usable for pairing there. The two row counts are different objects and we do not mix them.) A further \textbf{13 rows} carried a parser artefact, in which a run name read as a budget what was not one (\texttt{g1p} as 1\%, \texttt{dose2p5x} as 2\%); these are now left \textbf{unlabelled rather than guessed}. The recovery does not change any cell count: the complete 5x2 label-by-epoch rectangles remain \textbf{6}, the cells with both arms remain \textbf{168}, and the resolvability audit of \S{}9 is unaffected because it does not require a budget label. In \textbf{25} rows the stored \texttt{seed\_used} differs from the seed embedded in the run name; every one of those rows stores \texttt{42} while the name carries a different seed, and every one is an incomplete or restarted run. Consequently in \textbf{5} (cell, seed) slots the base and strategy arms do not share the same embedded seed, so for those slots "the same seed" is a pairing asserted from the stored field rather than one verified from the run names. Four of the five slots are residual runs and one is a released cell, so the effect on reported values is small; the boundary is recorded because the paired-by-seed convention is stated as a property of the whole archive.

\textbf{What follows from the first three.} The released per-run table supports recomputation of every reported number, and it does so under a key that is narrower than the design: loss prior is fixed by selection rather than by the key, the protocol flag is confounded with a family, and a small residual set carries a seed field that its name contradicts. A reader who needs a cell that varies loss prior, or that separates protocol from family, will not find it in this release.

\subsection*{P.10 Arm-level accounting for a falling budget gain}
\addcontentsline{toc}{subsection}{P.10 Arm-level accounting for a falling budget gain}

\textbf{Prior work and scope.} Writing a change in a difference as the difference of two arm-level changes is an identity, not a method, and prior work on budgeted training already reports instances of both catch-up and degradation; the claim here is a \textbf{census} over a released archive, not a new estimator, and $C$ and $A$ are measurements rather than mechanism. A gain that falls as the epoch budget grows can be accounted for in two ways, separable without any new run. Fix a (pair, family, label budget) triple and two epoch budgets $l<h$; for every seed with all four runs present write $B_l,I_l,B_h,I_h$ for the two arms, and define

$C=B_h-B_l$, $A=I_h-I_l$, $D=C-A$.

$C$ is the absolute movement of the baseline arm, $A$ that of the intervention arm, and $D$ equals the fall in the gain, $G_l-G_h$. All quantities are in mAP50-95 points, intervals are conditional $t$ intervals on same-seed paired differences, the two columns are computed separately and never pooled, and \textbf{every decomposition is taken inside a single family} (pairing across families would fold a protocol difference into what is read as an epoch effect).

\textbf{Coverage.} 86 decompositions are computable on the released table, and the identity $D=C-A=G_l-G_h$ holds to \textbf{3.6e-15} in every one of them. Of these, 31 have $D\le0$, i.e. no fall in the gain, and are not classified below. The remaining \textbf{55} split as follows.

{\footnotesize
\setlength{\tabcolsep}{4pt}
\begin{longtable}{p{0.393\textwidth}p{0.479\textwidth}>{\raggedleft\arraybackslash}p{0.058\textwidth}}
\toprule
composition of a falling gain & meaning & count \\
\midrule
\endhead
\textbf{baseline catch-up, $C>A\ge0$} & both arms improve, the baseline more & \textbf{47} \\
\textbf{intervention-arm degradation, $C>0>A$} & baseline improves while the intervention arm falls & \textbf{5} \\
\textbf{both arms fall, $A$ further} & $C\le0$ and $A<0$ & \textbf{3} \\
\botrule
\end{longtable}}

Where the gain does fall, the dominant account is \textbf{that both arms improve and the baseline improves more} (47 of 55), while an absolute degradation of the intervention arm occurs in five cases. This is a statement about the archive obtained by an identity, and it is what a bare "gain falls with budget" curve cannot make: the curve is consistent with either account, and here one is rare.

\textbf{One instance where the columns disagree about the arms.} For \texttt{pcb-\textgreater{}neu\_det / s2df / 50\% / 30-\textgreater{}100} the test column gives $A=-0.302$ with $D=+1.435$, while val-best gives $A=+9.152$ with $D=-0.971$: the columns disagree not only about the sign of the gain but about \textbf{which arm moved}, so a reporting convention here decides a question about the training dynamics.

\textbf{Limits.} (i) The paired counts are mostly 3 to 10, and the intervals cover data-order randomness only, not population training variance. (ii) The epoch endpoints are \textbf{not uniform} across entries (30-\textgreater{}100, 30-\textgreater{}200, 50-\textgreater{}400 and so on), so magnitudes must not be compared across rows. (iii) Eight of the twenty largest entries come from one pair and family, \texttt{pcb-\textgreater{}neu\_det / s2df}, so the count of independent units is far below the count of entries. (iv) A single cell can yield opposite arm-level readings under the two columns, so "which arm moved" is not a property of the cell alone. (v) $C$ and $A$ are measured quantities and do not identify the headroom or mismatch terms that \S{}8 discusses; we do not read them as mechanism.

\section*{Appendix Q. Limitations in full}
\addcontentsline{toc}{section}{Appendix Q. Limitations in full}

\textbf{Detailed limitations; the article retains one sentence per item and points here.}

\textbf{Q.1 The power audit, in full.}

\begin{itemize}
\item \textbf{Resolvability (our own power audit).} Every cell in the archive was scored against its own 80\%-power minimum detectable difference (paired \emph{t}, \ensuremath{\alpha} = 0.05, the same convention as \S{}4.2): \textbf{43 of 147 cells have \textbar{}\ensuremath{\Delta}\textbar{} below their MDE and therefore cannot support a sign claim} --- \textbf{32 of 73 cells at three seeds (44\%)} against \textbf{11 of 74 at more than three seeds (15\%)}. \textbf{Scope of this count: it is about single-budget \emph{levels}, not about the \emph{differenced} quantities that every \S{}4 conclusion lives on} --- the audit does not transfer to differenced readings, and Appendix P.5 states that limit. The headline cells are unaffected: the largest \S{}4.1 readings sit at \textbf{13.9\ensuremath{\times}} (the \texttt{t1c} cell at 100 epochs) and \textbf{13.4\ensuremath{\times}} (the \texttt{mech} cell at 200 epochs) their own level MDE, and the two primary \S{}4.1 readings sit at \textbf{4.6\ensuremath{\times} and 1.4\ensuremath{\times}}; those are \textbf{level} ratios, whereas the \emph{differenced} statement the \S{}4.1 table actually makes sits at \textbf{4.6\ensuremath{\times} and 1.6\ensuremath{\times}}, which is the number the conclusion should be read against, but the label-axis counts of \S{}4.6 and the 100-epoch backbone point of \S{}4.4 are in the affected band and are qualified where they are reported. The per-cell table is released with this paper.
\end{itemize}

\textbf{Q.2 Resolvability as a function of the reporting column.}

\begin{itemize}
\item \textbf{Resolvability is itself a function of the reporting column.} Switching between the test and val-best columns changes the resolvable / unresolvable verdict of a cell --- about \textbf{11\%} of gated cells at both 30 and 100 epochs, against \textbf{four (pair, family) reporting strata} (independence and image-level overlap between them are not established by the released run table) --- so whether a cell \emph{can be adjudicated} is a property of the convention, not of the cell. Detail, including the per-cell list and the robustness checks, is in Appendix P.8.
\end{itemize}

\textbf{Q.3 The metric as a second reporting convention.}

\begin{itemize}
\item \textbf{The metric is a second reporting convention, and the two metrics can disagree in sign.} mAP50 and mAP50-95 give \textbf{different signs on the same cell}: at 100 epochs the strategy gain turns negative in mAP50 while staying non-negative in mAP50-95 in \textbf{4 of the 15} long-budget cells (\textbf{Fig. S3}, Table T12). This paper therefore \textbf{takes mAP50-95 as its primary metric} and bounds the separability it claims between the two --- the divergence is reported as a property of the metric pair rather than read as evidence about the intervention, and it is, after the column of the previous bullet, a second reason that a sign reading here belongs to a stated convention.
\end{itemize}

\textbf{Q.4 The audit as one criterion.}

\begin{itemize}
\item \textbf{The same audit as one criterion.} The audit's threshold and \S{}4.2's MDE are the same object seen twice: with $\hat\sigma$ the between-seed SD of the paired differences and $\gamma(n) := (t_{1-\alpha/2,\,n-1}+t_{1-\beta,\,n-1})/\sqrt n$ a function of $n$ alone ($\gamma(3) = 3.10$, $\gamma(10) = 0.99$), a cell is resolvable \textbf{iff} $|\bar\Delta|/\hat\sigma \ge \gamma(n)$, so the counts above and \S{}4.2's MDE are one judgement taken at two sample sizes. Two consequences are derived in Appendix P: the \textbf{44\%-against-15\% contrast is not a mechanical consequence of $n$}, because the median $|\bar\Delta|/\hat\sigma$ is \textbf{4.37} among the three-seed cells against \textbf{3.03} among the ten-seed cells; and the ten-seed bar used for the registered 0.30 pp comparison sits \textbf{one seed above the smallest integer} that can resolve it at the archive's median noise ($n=9$, a fact about the archive's noise and not a rule we followed; the per-cell MDEs this audit uses are computed per cell and are not all at ten seeds --- they run as high as 1.255 and 4.27 pp), so the size we chose is a \textbf{retrospective account, not an ex-ante rule}.
\end{itemize}

\textbf{Q.5 What the audit covers, and what it does not.}

\begin{itemize}
\item \textbf{What the audit covers, and what it does not.} It is a statement about \textbf{single-budget levels}, not about the \textbf{differenced} quantities every \S{}4 conclusion lives on, and it is not the \textbf{seed gate} of \S{}4.5 --- the gate asks whether a \textbf{sign is stable across seeds}, the audit whether a \textbf{level is distinguishable from zero}. A cell can pass the gate and fail the audit; where that happens the gate's eligibility stands and only the right to read that cell's \textbf{magnitude} is withdrawn (Appendix P.5).
\end{itemize}

\textbf{Q.6 Headroom and shift not fully separated.}

\begin{itemize}
\item \textbf{Not fully separated.} Headroom and shift covary on the graded cells. One moderator does track the sign across corpora: the \textbf{annotation density of the target domain}. Among the nine corpora whose label statistics we can measure, the three sparsest (\texttt{mask\_clean} 0.10, \texttt{mendeley} 0.59, \texttt{firesmoke} 0.68 instances per image) carry \textbf{negative} mean gains --- eleven of eleven cells negative on the first, eight of eight on the second --- while the two densest (\texttt{dota15} 80.9, \texttt{visdrone} 70.7) carry \textbf{+2.16 and +0.95 pp} with one and two negative cells out of 33 and 8. The corpus-level correlation is \textbf{r \ensuremath{\approx} 0.34}, \texttt{neu\_det} (2.29 instances per image, +2.70 pp) is a clear exception, and density is \textbf{confounded with the pair} because each corpus supplies a fixed set of pairs. A related but distinct notion of density is already used in detection transfer: for LiDAR, the \textbf{beam-density difference between collectors} is treated as a \emph{source of the domain gap} and removed by a density-insensitive framework \citep{hu2023}. What is reported here is a different quantity in a different role --- \textbf{annotation density per image}, read as a candidate moderator of the sign of the gain, and confounded with the pair. We therefore report density as a \textbf{candidate} moderator, not an identified one. It is also not the first corpus statistic used to predict transfer behaviour: \textbf{label statistics predict transfer difficulty without training a model}, via the conditional entropy between source and target label assignments \citep{tran2019}, and \textbf{source-model accuracy predicts how well a model transfers} \citep{kornblith2019}. What those works do not do is treat a \textbf{countable per-corpus annotation statistic} as a candidate moderator of the \textbf{sign} of a paired gain, which is the role read here; with nine corpora, a correlation of 0.34 and a confound with the pair, the reading is offered as a signpost for a design that varies density deliberately, not as an identified effect.
\end{itemize}

\section*{Appendix R. Fig. 1 in full: column conventions, audit placement and exceptions}
\addcontentsline{toc}{section}{Appendix R. Fig. 1 in full: column conventions, audit placement and exceptions}

\textbf{The caption is a schematic; this is what it abbreviates.} Column convention: the intervention and measurement results (\S{}4.1, \S{}4.2, \S{}4.5, \S{}4.6) are reported on the \textbf{held-out test split} (the corpus's own \texttt{test:} entry) --- \textbf{except the two \S{}4.2 rows and \S{}4.1's T1-c, whose released \texttt{test:} key also selected the checkpoint, and \S{}4.1's \texttt{mech} cell, whose released \texttt{test:} and \texttt{val:} entries are the same directory, so for those the reading is not an independent held-out split} (Appendix M.3) --- whereas the two control subsections (\S{}4.3, \S{}4.4) use the \textbf{validation-selected best checkpoint's validation reading}; \S{}4.4's isolated backbone swap is where the two conventions put the \textbf{cross-backbone difference} on opposite sides of zero on the primary target's corpus, although each backbone's own gain is positive under both conventions. A third reading appears once, in \S{}4.6: the same-domain contrast there is stated under the archive's own endpoint rule rather than under either column. A fourth exception is the 2\ensuremath{\times}2 sensitivity test of \S{}5, reported on the \textbf{val} column --- the same column as \S{}4.3 and \S{}4.4 (Supplementary N.4). \textbf{Of the two, the test column is this paper's primary column and the val-best column is its reported sensitivity}; the subsections that read the val-best column are the paper's \textbf{sensitivity group}, and their numbers are labelled as sensitivity wherever they are quoted. The same cell can differ between the two, and by more than a rounding step --- on \texttt{shwd2sf\ensuremath{\rightarrow}sfchd20} (100 epochs, thirteen paired seeds) the test column reads \textbf{+0.628 pp} while the best-checkpoint validation column reads \textbf{\ensuremath{-}0.307 pp}, a sign flip produced by the choice of column alone (\S{}6(b)); every number in this paper states which of the two it uses, and the released per-run tables carry both columns. Sign convention: gains are always differenced as a \textbf{later or higher budget minus an earlier or lower budget}, so a \textbf{negative} difference is the direction the budget account predicts; the two \S{}4.2 rows and the \S{}4.5 label axis are on the same convention. Audit placement: the power audit of \S{}9 and of Appendix P.1 is a \textbf{level} audit, and the \S{}4 conclusions are \textbf{differenced} readings that it does not cover (Appendix P.5); where a headline reading is quoted with its own MDE, the table is Appendix S. Exceptions are collected here: (i) the four endpoint-rule exceptions above; (ii) the \texttt{val}-column 2\ensuremath{\times}2; (iii) \S{}4.6's archive-endpoint-rule row; (iv) the \texttt{dota15\ensuremath{\rightarrow}dota15} within-domain control, whose third run's directory name omits the seed suffix that the other two carry (\S{}K.1).

\section*{Appendix S. Headline numbers with their column and their MDE ratio}
\addcontentsline{toc}{section}{Appendix S. Headline numbers with their column and their MDE ratio}

\textbf{Why this table exists.} A reader otherwise assembles it from P.1's level audit, P.1-bis's differenced audit, P.7's per-batch columns and Q.1's scope note. \textbf{What it is not}: a complete inventory. Two quantities the article cites most often --- \S{}4.1's \ensuremath{-}1.143 pp and \ensuremath{-}0.223 pp \textbf{differences} --- are single-budget differenced readings that no released table prints as a row, so they are absent here and are given in N.1 instead. \texttt{MDE = gamma(n) * sd} at alpha = 0.05, beta = 0.20 (paired \emph{t}), computed per cell; \texttt{ratio = \textbar{}reading\textbar{} / MDE}. \textbf{Only rows whose reading and MDE are both recomputed from \texttt{base/run\_table\_canonical.csv} are listed}; the remaining headline readings stay in P.1/P.1-bis until their cell's MDE is computed the same way.

{\footnotesize
\setlength{\tabcolsep}{4pt}
\begin{longtable}{p{0.367\textwidth}p{0.261\textwidth}p{0.077\textwidth}p{0.071\textwidth}>{\raggedleft\arraybackslash}p{0.047\textwidth}>{\raggedleft\arraybackslash}p{0.053\textwidth}}
\toprule
headline reading & cell (pair / family / label budget / epochs) & column & reading (pp) & MDE (pp) & ratio \\
\midrule
\endhead
\S{}4.1 visdrone\ensuremath{\rightarrow}dota15 100 ep & visdrone\ensuremath{\rightarrow}dota15 / \texttt{t1c} / 20\% / 100 & held-out test & +3.727 & 0.268 & \textbf{13.9x} \\
\S{}4.1 visdrone\ensuremath{\rightarrow}dota15 200 ep & visdrone\ensuremath{\rightarrow}dota15 / \texttt{mech} / 20\% / 200 & held-out test & +2.584 & 0.193 & \textbf{13.4x} \\
\S{}4.1 same cell, val-best & visdrone\ensuremath{\rightarrow}dota15 / \texttt{t1c} / 20\% / 100 & val-best & +3.733 & 0.275 & 13.6x \\
\S{}4.1 same cell, val-best & visdrone\ensuremath{\rightarrow}dota15 / \texttt{mech} / 20\% / 200 & val-best & +2.588 & 0.183 & 14.2x \\
\S{}4.1 aitod\ensuremath{\rightarrow}visdrone 100 ep & aitod\ensuremath{\rightarrow}visdrone / \texttt{t1b} / 20\% / 100 & held-out test & +0.168 & 0.119 & \textbf{1.4x} \\
\S{}4.1 same cell, val-best & aitod\ensuremath{\rightarrow}visdrone / \texttt{t1b} / 20\% / 100 & val-best & +1.152 & 0.184 & 6.2x \\
\S{}4.2 smoke\ensuremath{\rightarrow}SFCHD ladder, 30 ep & smoke\ensuremath{\rightarrow}sfchd / \texttt{b2} / 20\% / 30 & held-out test & +3.656 & 0.435 & 8.4x \\
\S{}4.2 same ladder, 50 ep & smoke\ensuremath{\rightarrow}sfchd / \texttt{b2} / 20\% / 50 & held-out test & +2.674 & 0.267 & 10.0x \\
\S{}4.2 same ladder, 100 ep & smoke\ensuremath{\rightarrow}sfchd / \texttt{b2} / 20\% / 100 & held-out test & +1.666 & 0.200 & 8.3x \\
\S{}4.2 same ladder, 200 ep & smoke\ensuremath{\rightarrow}sfchd / \texttt{b2} / 20\% / 200 & held-out test & +1.495 & 0.252 & 5.9x \\
\S{}4.5 pcb\ensuremath{\rightarrow}neu\_det, 10\% level, 100 ep & pcb\ensuremath{\rightarrow}neu\_det / \texttt{s2df} / 10\% / 100 & held-out test & +2.381 & 1.738 & 1.4x \\
\S{}4.5 visdrone\ensuremath{\rightarrow}dota15, 10\% level, 30 ep & visdrone\ensuremath{\rightarrow}dota15 / \texttt{s2ae} / 10\% / 30 & held-out test & +1.995 & 0.266 & 7.5x \\
\S{}4.5 same pair, \textbf{10\%\ensuremath{\rightarrow}50\% labels} (the differenced reading) & visdrone\ensuremath{\rightarrow}dota15 / \texttt{s2ae} & held-out test & +1.978 & 0.419 & 4.7x \\
\S{}4.5 pcb\ensuremath{\rightarrow}neu\_det, \textbf{10\%\ensuremath{\rightarrow}50\% labels} (the differenced reading) & pcb\ensuremath{\rightarrow}neu\_det / \texttt{s2df} & held-out test & \ensuremath{-}2.391 & 2.396 & \textbf{1.0x} \\
\botrule
\end{longtable}}

\textbf{What the table makes plain, on the rows it carries.} Every \textbf{level} reading here sits at \textbf{1.4x or more} of its cell's own MDE, and the two \S{}4.1 levels sit at 13-14x. The \textbf{differenced} label-axis readings are the exception: \texttt{pcb\ensuremath{\rightarrow}neu\_det}'s 10\%\ensuremath{\rightarrow}50\% change sits at \textbf{1.0x}, i.e. exactly at its own MDE, which is why \S{}9 reports the label-axis cells as descriptions rather than as resolvable per-cell claims, while \texttt{visdrone\ensuremath{\rightarrow}dota15}'s sits at 4.7x. The level-versus-difference distinction is visible in the columns, since \S{}4.1's endpoint difference is not a level and its own MDE is computed on the per-seed differences.

\section*{Appendix T. The reporting cap and the two stratifications, measured on released checkpoints}
\addcontentsline{toc}{section}{Appendix T. The reporting cap and the two stratifications, measured on released checkpoints}

\textbf{What this appendix adds.} Three readings the article refers to but that were not in the released set before: the \textbf{reporting cap} re-measured on four cell families and both columns, and the same cells \textbf{stratified by object area} and \textbf{by class}. All of it is evaluation only - no checkpoint was retrained - and all of it is on the held-out \texttt{test} column at \texttt{max\_det = 300} unless stated otherwise. \textbf{The per-run results are released}, so that no number here rests on our own table: \texttt{provenance/appendix\_T/T1\_maxdet\_levels.csv} (the cap ladder, one row per run, split and cap), \texttt{T2\_area\_AP.csv} (the area-stratified AP, one row per run and area bin) and \texttt{T3\_perclass\_AP.csv} (the per-class AP, one row per run and class), with the conventions and the three limitations that must be read with them stated in that directory's \texttt{README.md}.

\subsection*{T.1 The reporting cap (\texttt{max\_det}) does not move the paired gains}
\addcontentsline{toc}{subsection}{T.1 The reporting cap (\texttt{max\_det}) does not move the paired gains}

The article's claims are paired strategy-minus-baseline gains, so that is the quantity to check. Each cell family was re-read at \texttt{max\_det} = 100, 300 and 1000 from the same checkpoints (10 paired seeds per family and arm, 146 of 160 scheduled evaluations completed, no failures).

\textbf{Two conventions must be read with this table, because they are what make two of its numbers differ from the same cells printed elsewhere in this paper.}

\begin{enumerate}
\item \textbf{The "val" rows are not this paper's val-best readings.} They are the \texttt{val:} directory re-read with this appendix's own evaluation setting (\texttt{conf = 0.001}, COCO-style 101-point AP, \texttt{max\_det} as stated). The \texttt{t1c} cell's val-best reading as reported in \S{}4.1 and Appendix S is \textbf{+3.733 pp}; the +7.063 pp below is the same directory read in the other convention. The two are different objects and neither supersedes the other.
\item \textbf{The paired gains are the quantity to compare, and they are cap-stable.} The \emph{levels} are not comparable across columns or conventions, which is exactly the point of the section.
\end{enumerate}

{\footnotesize
\setlength{\tabcolsep}{4pt}
\begin{longtable}{p{0.278\textwidth}p{0.194\textwidth}>{\raggedleft\arraybackslash}p{0.119\textwidth}>{\raggedleft\arraybackslash}p{0.053\textwidth}>{\raggedleft\arraybackslash}p{0.106\textwidth}p{0.128\textwidth}}
\toprule
family / column (cell key: \texttt{\textless{}pair\textgreater{}/\textless{}family\textgreater{}/\textless{}labels\textgreater{}/\textless{}epochs\textgreater{}}) & evidence & gain at 300 & gain at 1000 & gain at 100 & sign or significance changed? \\
\midrule
\endhead
\texttt{t1c} / test --- \texttt{visdrone-\textgreater{}dota15} / \texttt{t1c} / 20\% / 100 & independent held-out & \textbf{+3.726 pp (p = 1.7e-11)} & +3.659 pp & +3.614 pp & no \\
\texttt{t1c} / val --- same cell, val-directory re-read & not this paper's val-best (see convention 1) & +7.063 pp & +6.948 pp & +7.028 pp & no \\
\texttt{mech} / val(=test) --- \texttt{visdrone-\textgreater{}dota15} / \texttt{mech} / 20\% / 200 & same directory as its own test & +2.701 pp & +2.682 pp & +2.688 pp & no \\
\texttt{d15d15} / val --- \texttt{dota15-\textgreater{}dota15} / \texttt{r10} / 20\% / 30 & same-domain control & -1.441 pp & -1.516 pp & -1.430 pp & no \\
\texttt{d15d15} / test & same-domain control & -0.006 pp (p = 0.973) & -0.008 pp & -0.006 pp & no \\
\texttt{d15ai} / val --- \texttt{dota15-\textgreater{}aitod} / \texttt{d15ai} / 20\% / 100 & cross-domain control, n = 10 & -0.713 pp (p = 0.0090) & -0.693 pp & -1.014 pp & no \\
\texttt{d15ai} / test & cross-domain control, n = 10 & +0.264 pp (p = 0.034) & +0.298 pp & \textbf{+0.154 pp (p = 0.19)} & \textbf{yes at 100 only} \\
\texttt{aitodvis} / test --- \texttt{aitod-\textgreater{}visdrone} / \texttt{aitodvis} / 20\% / 100 & n = 3 only & +0.147 pp (p = 0.19) & +0.150 pp & +0.177 pp (p = 0.19) & no \\
\botrule
\end{longtable}}

\textbf{Why 14 of 160 evaluations are absent}: the \texttt{d15ai} family's released checkpoints begin at seed 45 and one run (\texttt{d15\_ai\_base100\_s46n}) has no \texttt{args.yaml}, so 12 runs could not be read and two more could not resolve their data path; the table reports the cells that could be read. \textbf{The paired-gain drift over the whole table is 0.002-0.115 pp} (the figure 0.067-0.115 pp holds for the \texttt{t1c} rows alone). \textbf{The cap changes the levels and not the pairing.} The level drift from 300 to 1000 is up to \textbf{0.51 pp} (levels are in the released \texttt{T1\_maxdet\_levels.csv}, not in this table), which exceeds \texttt{t1c}'s 0.268 pp MDE, while the paired-gain drift is \textbf{0.067-0.115 pp}. So a reader comparing \emph{levels} across cells is reading a quantity the cap moves, and a reader comparing \emph{paired gains} is not. \textbf{The fragile end is 100, not 300}: at \texttt{max\_det} = 100 the \texttt{d15ai} test-column gain loses significance (+0.154 pp, p = 0.19), which is the only place in the table where the conclusion depends on the cap.

\subsection*{T.2 Stratified by object area: the gain grows with object size}
\addcontentsline{toc}{subsection}{T.2 Stratified by object area: the gain grows with object size}

COCO area bins (small \textless{} 32\textasciicircum{}2 px, medium 32\textasciicircum{}2-96\textasciicircum{}2, large \textgreater{} 96\textasciicircum{}2), paired over 10 seeds, \texttt{max\_det = 300}, \texttt{conf = 0.001}, test column. AP is the COCO definition (101-point interpolation, IoU 0.50:0.05:0.95), computed directly because \texttt{pycocotools} is not present in the evaluation environment - a limitation, since this is therefore not an independent implementation of the metric.

{\footnotesize
\setlength{\tabcolsep}{4pt}
\begin{longtable}{p{0.453\textwidth}>{\raggedleft\arraybackslash}p{0.098\textwidth}>{\raggedleft\arraybackslash}p{0.147\textwidth}>{\raggedleft\arraybackslash}p{0.098\textwidth}>{\raggedleft\arraybackslash}p{0.098\textwidth}}
\toprule
family & all & \textbf{small} & medium & large \\
\midrule
\endhead
\texttt{t1c} (visdrone-\textgreater{}dota15, 100 ep) & +2.39 pp & \textbf{+0.89 pp} & +3.66 pp & +5.89 pp \\
\texttt{mech} (visdrone-\textgreater{}dota15, 200 ep) & +2.11 pp & \textbf{+1.05 pp} & +3.45 pp & +3.83 pp \\
\texttt{d15d15} (same-domain control, 30 ep) & -0.54 pp & -0.66 pp & -0.56 pp & +0.07 pp \\
\botrule
\end{longtable}}

\textbf{The \texttt{all} column here is not the article's headline reading for the same cells} (T2 uses \texttt{conf = 0.001} and this appendix's own AP implementation, so \texttt{t1c}'s \texttt{all} is +2.39 pp against the headline +3.727 pp). The two are different objects: the headline number is the evaluator's mAP50-95 at its own confidence threshold, and this column is the same split under the stratified read. \textbf{Two cells independently give the same shape}: the intervention's gain rises monotonically with object size (small -\textgreater{} medium -\textgreater{} large), and the small-object level itself is low (0.044 AP on \texttt{t1c}). The same-domain control sits near zero in every bin, so the shape is a property of the intervention rather than of the metric. The consequence for the article is that the single-cell mAP50-95 gain of about +2.4 pp is \textbf{carried mainly by medium and large objects}, and any statement of it as a percentage should say so.

\subsection*{T.3 Stratified by class: the gain is concentrated where the instances are few}
\addcontentsline{toc}{subsection}{T.3 Stratified by class: the gain is concentrated where the instances are few}

Per-class AP is taken from the evaluator's own per-class output (not recomputed), \texttt{max\_det = 300}, 10 seeds.

{\footnotesize
\setlength{\tabcolsep}{4pt}
\begin{longtable}{p{0.440\textwidth}p{0.261\textwidth}p{0.228\textwidth}}
\toprule
 & \texttt{t1c} & \texttt{mech} \\
\midrule
\endhead
classes with a positive gain & \textbf{15 of 16} & \textbf{15 of 16} \\
largest & \texttt{soccer-ball-field} \textbf{+9.79 pp} & \texttt{swimming-pool} \textbf{+6.85 pp} \\
\texttt{small-vehicle} (the dominant class by instance count) & \textbf{+0.99 pp} & \textbf{+1.10 pp} \\
\texttt{container-crane} & 0.00 pp (base also 0.000) & 0.00 pp \\
\botrule
\end{longtable}}

\textbf{The direction is stable across classes but the magnitude is not.} The classes that carry most of the gain are those with few instances (playing fields, swimming pools, harbours), while \texttt{small-vehicle} - the class that dominates these dense corpora - gains about 1 pp. This agrees with T.2 from an independent stratification: both say the aggregate gain does not come from the small objects or from the dominant class. \textbf{The family correction leaves every class significant, so what varies is the size and not the direction}: Benjamini--Hochberg over the \textbf{15 testable classes} per cell (the sixteenth, \texttt{container-crane}, has zero AP in both arms in both cells and is uninformative) keeps \textbf{all fifteen below q = 0.05 in both cells} --- even the smallest, \texttt{small-vehicle} at \textbf{+0.99 pp}, reaches t = 9.6 --- because the per-seed spread of a per-class gain is small. The class-level result is therefore that the \textbf{direction is reliable everywhere and the magnitude is not}.

\section*{Appendix U. The pointwise $r_D$ measured at the released endpoints}
\addcontentsline{toc}{section}{Appendix U. The pointwise $r_D$ measured at the released endpoints}

\textbf{What was measured.} The article's \S{}6 states a regime judgement from the ratio $r_D=\|\nabla R_D(\theta)\|/(\|H_D(\theta)\|\,\rho\|\theta\|)$ and previously said it was not measured, because the numerator is a gradient norm at fine-tuned endpoints. It is now measured: the released endpoint script records that norm together with the Lanczos spectral estimate, so the ratio follows per checkpoint. \textbf{Evaluation only --- no checkpoint was retrained.} Settings: \texttt{--batch 4 --nbatches 4 --lanczos 20}, \texttt{split val}, \texttt{rho = 0.001}.

\textbf{Coverage.} 99 checkpoints, i.e. the same cells the \textbf{path} form uses, so the two forms are compared on one sample. Per-run (the independent unit) at \texttt{rho = 0.001}:

{\footnotesize
\setlength{\tabcolsep}{4pt}
\begin{longtable}{p{0.268\textwidth}>{\raggedleft\arraybackslash}p{0.122\textwidth}>{\raggedleft\arraybackslash}p{0.122\textwidth}>{\raggedleft\arraybackslash}p{0.122\textwidth}>{\raggedleft\arraybackslash}p{0.122\textwidth}>{\raggedleft\arraybackslash}p{0.122\textwidth}}
\toprule
quantity & min & p25 & median & p75 & max \\
\midrule
\endhead
\texttt{r\_D} & \textbf{0.0081} & \textbf{0.0224} & \textbf{0.0283} & \textbf{0.0452} & \textbf{0.1029} \\
\texttt{r\_D} at \texttt{rho = 0.003} & 0.0027 & 0.0075 & 0.0094 & 0.0151 & 0.0343 \\
\texttt{r\_D} at \texttt{rho = 0.01} & 0.00081 & 0.0022 & 0.0028 & 0.0045 & 0.0103 \\
\botrule
\end{longtable}}

\textbf{The result, and the honest bound on it.} Under the article's own conventions ($r_D\le 1/16$ curvature-dominated, $r_D\ge1$ slope-dominated) \textbf{86.9\% of runs are below $1/16$ and none reaches $1$}, while the \textbf{path} form on the same cells has median \textbf{10.8}. So the two forms disagree in direction by about three orders of magnitude, and the disagreement is not a defect of either: the pointwise ratio asks about the endpoint's own Hessian, the path ratio about a secant modulus over the displacement, and \textbf{a regime judgement is therefore relative to the form asked} ($\S$6). \textbf{The curvature reading is an \textasciitilde{}87\% statement, not a universal one}: 13.1\% of runs exceed $1/16$ (largest 0.1029, still an order of magnitude below $1$). By target corpus the medians run from 0.016 (visdrone) to 0.039 (mende).

\textbf{Two things a reader must know about the sample.}

\begin{enumerate}
\item \textbf{The released driver drops 92 of the 180 endpoint rows.} \texttt{p1\_rD\_run.py} resolves each row's data yaml by searching only \texttt{/root/datasets\_mask} and \texttt{/root/datasets}, but the sfchd rows' yamls live in \texttt{/root/} and \texttt{/root/xcal/}, so an entire target domain (55 of the 99 runs) is skipped; the ceiling under the published command is 88 rows / 44 runs. Those rows were recovered with a driver identical in every measurement argument that only extends the search roots and checks that each yaml's image directories exist. The driver's resolver is fixed in the released copy.
\item \textbf{180 rows are not 180 independent readings.} The endpoint table's rows differ in \texttt{(draw, stride)}, but the driver passes neither, so rows of one run are recomputations of one observation --- 26 of the 27 multi-row runs are bitwise identical. \textbf{The independent unit is the run, n = 99}, which is the basis of the table above.
\end{enumerate}

\textbf{Released}: \texttt{provenance/appendix\_U/rD\_pointwise.csv} (per row) and \texttt{rD\_pointwise\_perrun.csv} (per run), with the columns the computation needs.

\subsection*{U.1 The independent evaluation stack (official \texttt{pycocotools})}
\addcontentsline{toc}{subsection}{U.1 The independent evaluation stack (official \texttt{pycocotools})}

The article's readings all come from one implementation (Ultralytics 8.4.120 \texttt{val}), and T.2's stratified AP is our own recomputation, which T.2 itself flags as not an independent implementation. Both cells' headline runs were therefore re-scored from their \textbf{predictions} by the \textbf{official \texttt{pycocotools.COCOeval}} (predictions emitted as COCO JSON at \texttt{max\_det = 300}, \texttt{conf = 0.001}, \texttt{iou = 0.7}; scoring done by the official library, not by us):

{\footnotesize
\setlength{\tabcolsep}{4pt}
\begin{longtable}{p{0.195\textwidth}p{0.104\textwidth}>{\raggedleft\arraybackslash}p{0.326\textwidth}>{\raggedleft\arraybackslash}p{0.287\textwidth}}
\toprule
cell & arm & official \texttt{pycocotools} AP & Ultralytics, same runs \\
\midrule
\endhead
\texttt{t1c} (100 ep) & base & \textbf{0.1105} & 0.1156 \\
 & strategy & \textbf{0.1453} & --- \\
\texttt{mech} (200 ep) & base & \textbf{0.1329} & --- \\
 & strategy & \textbf{0.1548} & --- \\
\botrule
\end{longtable}}

\textbf{Paired gains}: \texttt{t1c} \textbf{+3.48 pp} (t = +44.1) under the official stack against \textbf{+3.73 pp} under Ultralytics on the same runs, and \texttt{mech} \textbf{+2.19 pp}. So the \textbf{level} moves by 0.51 pp (1.9x the \texttt{t1c} MDE of 0.268 pp) while the \textbf{paired gain} moves by 0.25 pp, inside its own MDE: the \S{}4 conclusions are a property of the contrast, not of the scoring implementation. Two bounds: \texttt{mech}'s \texttt{test} column is its own selection column (not an independent split, as R states), so that cell is not independent evidence in either stack; and the GT was converted from the YOLO labels rather than cross-checked against third-party COCO annotations.

\subsection*{U.2 The stratified readings under the official implementation}
\addcontentsline{toc}{subsection}{U.2 The stratified readings under the official implementation}

T.2's area and class strata are this package's own recomputation, and T.2 says so. The same predictions were therefore scored by the \textbf{official \texttt{pycocotools.COCOeval}} (predictions already in COCO form from the runs themselves), giving the strata in the official convention (\texttt{max\_det = 100}, which \texttt{COCOeval}'s area split uses; T.2 is at 300).

\begin{center}
{\small
\begin{tabular}{llrrrr}
\toprule
cell & arm & all & small & medium & large \\
\midrule
\texttt{t1c} & base & 0.1105 & 0.0344 & 0.1139 & 0.1687 \\
 & strategy & 0.1453 & 0.0405 & 0.1382 & 0.2487 \\
 & \textbf{gain} & \textbf{+3.48} & \textbf{+0.61} & \textbf{+2.43} & \textbf{+8.00} \\
\texttt{mech} & base & 0.1329 & 0.0458 & 0.1279 & 0.2195 \\
 & strategy & 0.1548 & 0.0479 & 0.1479 & 0.2611 \\
 & \textbf{gain} & \textbf{+2.20} & \textbf{+0.21} & \textbf{+2.00} & \textbf{+4.17} \\
\botrule
\end{tabular}}
\end{center}

\textbf{Direction is implementation-independent; magnitude is not.} Both implementations put the gain on the larger objects and leave the dominant small class near 1 pp, so the article's warning that a bare percentage overstates what a small-object detector gains holds in both. But \texttt{large} for \texttt{t1c} is \textbf{+8.00 pp} here against \textbf{+5.89 pp} in T.2 --- a \textbf{2.11 pp} difference, about eight times that cell's MDE --- so T.2's figures are statements \textbf{in this package's recomputation}, not constants across implementations. The per-class block is released in \texttt{provenance/appendix\_U/coco\_stratified\_official.csv}, and the two conventions (\texttt{max\_det} 100 against 300) are \textbf{not separated} from the implementation difference in this comparison.

\bibliography{sn-bibliography}
\end{document}
